% MATH 521 | Final Practice Workbook
% University of Wisconsin--Madison, Spring 2026
% Workbook build: Homework 1 through Homework 12

\documentclass[12pt]{article}

\usepackage[
  paperwidth=8.5in,
  paperheight=11in,
  top=1.0in,
  bottom=1.0in,
  left=1.0in,
  right=1.0in
]{geometry}

\usepackage[T1]{fontenc}
\usepackage{lmodern}
\usepackage{amsmath, amssymb, amsthm, mathtools}
\usepackage{parskip}
\usepackage{titlesec}
\usepackage{enumitem}
\usepackage{xcolor}
\usepackage{tcolorbox}
\tcbuselibrary{skins, breakable}
\usepackage{hyperref}

\hypersetup{
  colorlinks=true,
  linkcolor=blue!60!black,
  urlcolor=blue!60!black
}

\definecolor{probblue}{RGB}{215,232,255}
\definecolor{pass1green}{RGB}{218,245,218}
\definecolor{pass2yellow}{RGB}{255,252,215}
\definecolor{pass3orange}{RGB}{255,236,210}
\definecolor{pass4purple}{RGB}{240,218,255}
\definecolor{pass5gray}{RGB}{238,238,238}
\definecolor{solgreen}{RGB}{208,240,210}
\definecolor{reflectcyan}{RGB}{208,245,252}

\setcounter{tocdepth}{2}
\setlist[enumerate]{leftmargin=2.2em, itemsep=0.25em, topsep=0.25em}
\setlist[itemize]{leftmargin=2.0em, itemsep=0.25em, topsep=0.25em}

\titleformat{\section}
  [block]{\Large\bfseries\centering}{}{0em}{}[\vspace{2pt}\titlerule]
\titleformat{\subsection}
  [block]{\large\bfseries}{}{0em}{}[\vspace{1pt}\titlerule]
\titlespacing{\section}{0pt}{1.2em}{1em}
\titlespacing{\subsection}{0pt}{0.8em}{0.8em}

\newcommand{\R}{\mathbb{R}}
\newcommand{\Q}{\mathbb{Q}}
\newcommand{\Z}{\mathbb{Z}}
\newcommand{\N}{\mathbb{N}}
\newcommand{\closure}[1]{\overline{#1}}
\newcommand{\ball}[2]{B_{#1}(#2)}
\newcommand{\writespace}[1]{\vspace*{#1}}

\newcommand{\fillfield}[1]{%
  \par\noindent\textbf{#1:}\enspace\hrulefill\par\vspace{0.25in}%
}

\newcommand{\fillfieldbig}[2]{%
  \par\noindent\textbf{#1:}\par
  \writespace{#2}
  \noindent\rule{\linewidth}{0.35pt}\par\vspace{0.15in}%
}

\newcommand{\workquestion}[2]{%
  \par\noindent\textbf{#1}\par
  \writespace{#2}
  \noindent\rule{\linewidth}{0.35pt}\par\vspace{0.18in}%
}

\newenvironment{problemstatementbox}[1]{%
  \begin{tcolorbox}[
    enhanced,
    breakable,
    colback=probblue,
    colframe=blue!55!black,
    fonttitle=\bfseries,
    title={Problem Statement\hfill #1}
  ]
}{%
  \end{tcolorbox}
}

\newcommand{\firstdiagnosis}[1]{%
  \begin{tcolorbox}[
    enhanced,
    colback=pass1green,
    colframe=green!55!black,
    fonttitle=\bfseries,
    title={Pass 1 -- First Diagnosis\hfill #1},
    height=8.45in
  ]
  \fillfield{Topic}
  \fillfieldbig{What is being asked?}{1.05in}
  \fillfieldbig{Definitions likely relevant}{1.15in}
  \fillfieldbig{Theorem/proof method likely relevant}{1.15in}
  \fillfieldbig{What should the first line of the proof be?}{1.1in}
  \end{tcolorbox}
}

\newcommand{\recipecard}[1]{%
  \begin{tcolorbox}[
    enhanced,
    colback=pass2yellow,
    colframe=yellow!65!black,
    fonttitle=\bfseries,
    title={Pass 2 -- Proof Recipe Card\hfill #1},
    height=8.45in
  ]
  \fillfieldbig{Hypotheses}{0.75in}
  \fillfieldbig{Conclusion}{0.75in}
  \fillfieldbig{Definitions to unpack}{0.95in}
  \fillfieldbig{Useful theorem/lemma}{0.85in}
  \fillfieldbig{Key inequality/identity}{0.8in}
  \fillfieldbig{One-sentence proof strategy}{0.9in}
  \end{tcolorbox}
}

\newenvironment{selfexplanationbox}[1]{%
  \begin{tcolorbox}[
    enhanced,
    colback=pass3orange,
    colframe=orange!65!black,
    fonttitle=\bfseries,
    title={Pass 3 -- Self-Explanation\hfill #1},
    height=8.45in
  ]
}{%
  \end{tcolorbox}
}

\newenvironment{skeletonbox}[1]{%
  \begin{tcolorbox}[
    enhanced,
    breakable,
    colback=pass4purple,
    colframe=purple!55!black,
    fonttitle=\bfseries,
    title={Pass 4 -- Skeleton Proof\hfill #1}
  ]
}{%
  \end{tcolorbox}
}

\newcommand{\fullproofpage}[2]{%
  \begin{tcolorbox}[
    enhanced,
    colback=white,
    colframe=gray!60!black,
    fonttitle=\bfseries,
    title={Pass 5 -- Full Proof Attempt\hfill #1},
    height=8.45in
  ]
  \textbf{#2}

  \writespace{7.65in}
  \end{tcolorbox}
}

\newenvironment{solutionbox}[1]{%
  \begin{tcolorbox}[
    enhanced,
    breakable,
    colback=solgreen,
    colframe=green!45!black,
    fonttitle=\bfseries,
    title={#1}
  ]
}{%
  \end{tcolorbox}
}

\newcommand{\reflectionpage}[1]{%
  \begin{tcolorbox}[
    enhanced,
    colback=reflectcyan,
    colframe=cyan!55!black,
    fonttitle=\bfseries,
    title={Reflection \& Variation\hfill #1},
    height=8.45in
  ]
  \fillfieldbig{What was the key definition?}{1.15in}
  \fillfieldbig{What was the key proof move?}{1.15in}
  \fillfieldbig{What step would I forget under time pressure?}{1.15in}
  \fillfieldbig{What nearby exam variation could appear?}{1.35in}
  \end{tcolorbox}
}

\begin{document}

\hypersetup{pageanchor=false}
\begin{titlepage}
\centering
\vspace*{2.35cm}
{\Huge\bfseries MATH 521\par}
\vspace{0.35cm}
{\Large Real Analysis\par}
\vspace{1.1cm}
\rule{0.55\linewidth}{1pt}
\vspace{1.1cm}

{\huge\bfseries Final Practice Book\par}
\vspace{0.7cm}
{\Large Write-On Workbook\par}
\vspace{1.5cm}

{\large University of Wisconsin--Madison\par}
{\large Spring 2026\par}
\vspace{1.7cm}

{\normalsize\itshape Current build: Homework 1 through Homework 12\par}
\vspace{2.8cm}

{\large\bfseries Name:\enspace\hrulefill\par}
\vspace{0.9cm}
{\large\bfseries Date:\enspace\hrulefill\par}
\end{titlepage}
\hypersetup{pageanchor=true}

\pagenumbering{roman}
\tableofcontents
\newpage

\section*{How to Use This Workbook}
\addcontentsline{toc}{section}{How to Use This Workbook}

This workbook is for active proof practice. For each problem, work through the diagnosis,
recipe card, self-explanation, skeleton, and full proof attempt before reading the reference
solution.

\begin{description}[leftmargin=2.2cm, labelwidth=2.0cm, labelsep=0.2cm, itemsep=4pt]
  \item[Pass 1] Diagnose the topic, task, relevant definitions, and likely proof method.
  \item[Pass 2] Convert the problem into hypotheses, conclusion, tools, and a proof plan.
  \item[Pass 3] Answer targeted questions that expose the proof's logic.
  \item[Pass 4] Reconstruct the proof from a skeleton without reading the full solution.
  \item[Pass 5] Write a complete proof from scratch in the blank writing space.
\end{description}

Only after Pass 5 should you compare your work to the official or reference solution.

\newpage
\pagenumbering{arabic}

\section{Homework 1}

% ---------------------------------------------------------------------------
% HW1 Problem 1
% ---------------------------------------------------------------------------
\newpage
\subsection{Problem 1}

\begin{problemstatementbox}{Homework 1, Problem 1}
Using induction, prove that for all \(n\in\N\) we have
\[
  3+11+\cdots+(8n-5)=4n^2-n.
\]
\end{problemstatementbox}

\newpage
\firstdiagnosis{Homework 1, Problem 1}

\newpage
\recipecard{Homework 1, Problem 1}

\newpage
\begin{selfexplanationbox}{Homework 1, Problem 1}
\workquestion{1. What is the statement \(P(n)\) that should be proved by induction?}{0.95in}
\workquestion{2. What is the base case, and what are both sides when \(n=1\)?}{0.9in}
\workquestion{3. What is the next term added when passing from \(n\) to \(n+1\)?}{0.95in}
\workquestion{4. How does \(4n^2-n+(8n+3)\) simplify?}{0.9in}
\workquestion{5. Why is that expression equal to \(4(n+1)^2-(n+1)\)?}{0.9in}
\end{selfexplanationbox}

\newpage
\begin{skeletonbox}{Homework 1, Problem 1}
\workquestion{Define \(P(n)\) precisely:}{0.8in}
\workquestion{Base case \(n=1\): verify left-hand side equals right-hand side:}{0.85in}
\workquestion{Inductive hypothesis: assume \(P(n)\) for some \(n\in\N\):}{0.85in}
\workquestion{Add the next term \(8(n+1)-5=\) to both sides:}{0.85in}
\workquestion{Simplify the right-hand side and match the formula for \(P(n+1)\):}{0.95in}
\workquestion{Invoke the principle of mathematical induction:}{0.75in}
\end{skeletonbox}

\newpage
\fullproofpage{Homework 1, Problem 1}{Use this page for the complete induction proof.}

\newpage
\begin{solutionbox}{Student-Derived Solution (No official solution available)}
For \(n\in\N\), let \(P(n)\) denote the statement
\[
  3+11+\cdots+(8n-5)=4n^2-n.
\]
We argue by induction.

Base case. For \(n=1\), the left-hand side is \(3\), while the right-hand side is
\[
  4\cdot 1^2-1=3.
\]
Hence \(P(1)\) holds.

Inductive step. Assume \(P(n)\) holds for some \(n\in\N\), i.e.
\[
  3+11+\cdots+(8n-5)=4n^2-n.
\]
We must show that \(P(n+1)\) holds. The next term in the sum is
\[
  8(n+1)-5=8n+3.
\]
Adding this term to both sides of the inductive hypothesis gives
\[
  3+11+\cdots+(8n-5)+(8n+3)
  =
  (4n^2-n)+(8n+3).
\]
The right-hand side simplifies to
\[
  4n^2+7n+3.
\]
On the other hand,
\[
  4(n+1)^2-(n+1)
  =
  4(n^2+2n+1)-n-1
  =
  4n^2+7n+3.
\]
Therefore
\[
  3+11+\cdots+(8(n+1)-5)=4(n+1)^2-(n+1),
\]
so \(P(n+1)\) holds. By the principle of mathematical induction, \(P(n)\) is true for all \(n\in\N\).
\end{solutionbox}

\newpage
\reflectionpage{Homework 1, Problem 1}

% ---------------------------------------------------------------------------
% HW1 Problem 2
% ---------------------------------------------------------------------------
\newpage
\subsection{Problem 2}

\begin{problemstatementbox}{Homework 1, Problem 2}
Compute the following sum:
\[
  \frac1{1\cdot4}+\frac1{4\cdot7}+\cdots+\frac1{(3n-2)(3n+1)}.
\]
Prove that your answer is true for all \(n\in\N\) using induction.
\end{problemstatementbox}

\newpage
\firstdiagnosis{Homework 1, Problem 2}

\newpage
\recipecard{Homework 1, Problem 2}

\newpage
\begin{selfexplanationbox}{Homework 1, Problem 2}
\workquestion{1. What closed form does the telescoping computation suggest?}{0.95in}
\workquestion{2. What is the partial fraction decomposition of \(\frac1{(3k-2)(3k+1)}\)?}{1.0in}
\workquestion{3. What is the induction statement \(P(n)\)?}{0.9in}
\workquestion{4. In the inductive step, what is the next term?}{0.85in}
\workquestion{5. What algebra turns \(\frac n{3n+1}+\frac1{(3n+1)(3n+4)}\) into \(\frac{n+1}{3n+4}\)?}{1.0in}
\end{selfexplanationbox}

\newpage
\begin{skeletonbox}{Homework 1, Problem 2}
\workquestion{First compute the suspected formula using partial fractions/telescoping:}{1.0in}
\workquestion{State \(P(n)\):}{0.8in}
\workquestion{Base case \(n=1\):}{0.75in}
\workquestion{Inductive hypothesis: assume the sum through \(n\) equals \(\frac n{3n+1}\):}{0.85in}
\workquestion{Add the term \(\frac1{(3n+1)(3n+4)}\) and combine fractions:}{1.0in}
\workquestion{Conclude \(P(n+1)\) and invoke induction:}{0.75in}
\end{skeletonbox}

\newpage
\fullproofpage{Homework 1, Problem 2}{Use this page for the telescoping computation and induction proof.}

\newpage
\begin{solutionbox}{Student-Derived Solution (No official solution available)}
For \(n\in\N\), let \(P(n)\) denote the statement
\[
  \frac1{1\cdot4}+\frac1{4\cdot7}+\cdots+\frac1{(3n-2)(3n+1)}
  =
  \frac n{3n+1}.
\]
We prove \(P(n)\) for all \(n\in\N\) by induction.

Base case. \(P(1)\) says
\[
  \frac1{1\cdot4}=\frac1{3\cdot1+1},
\]
which is true.

Inductive step. Assume \(P(n)\) is true for some \(n\in\N\), that is,
\[
  \frac1{1\cdot4}+\frac1{4\cdot7}+\cdots+\frac1{(3n-2)(3n+1)}
  =
  \frac n{3n+1}.
\]
Adding the next term \(\frac1{(3n+1)(3n+4)}\) to both sides gives
\[
  \frac1{1\cdot4}+\cdots+\frac1{(3n-2)(3n+1)}+\frac1{(3n+1)(3n+4)}
  =
  \frac n{3n+1}+\frac1{(3n+1)(3n+4)}.
\]
The right-hand side is
\[
  \frac1{3n+1}\left(n+\frac1{3n+4}\right)
  =
  \frac1{3n+1}\cdot\frac{n(3n+4)+1}{3n+4}.
\]
Since
\[
  n(3n+4)+1=3n^2+4n+1=(3n+1)(n+1),
\]
we obtain
\[
  \frac n{3n+1}+\frac1{(3n+1)(3n+4)}
  =
  \frac{n+1}{3n+4}
  =
  \frac{n+1}{3(n+1)+1}.
\]
Thus \(P(n+1)\) holds. By induction,
\[
  \frac1{1\cdot4}+\frac1{4\cdot7}+\cdots+\frac1{(3n-2)(3n+1)}
  =
  \frac n{3n+1}
\]
for all \(n\in\N\).
\end{solutionbox}

\newpage
\reflectionpage{Homework 1, Problem 2}

% ---------------------------------------------------------------------------
% HW1 Problem 3
% ---------------------------------------------------------------------------
\newpage
\subsection{Problem 3}

\begin{problemstatementbox}{Homework 1, Problem 3}
Let \(f:A\to B\) be a function.
\begin{enumerate}[label=(\alph*)]
  \item Prove that \(f\) is injective if and only if \(f^{-1}(f(X))=X\) for all \(X\subseteq A\).
  \item Prove that \(f\) is surjective if and only if \(f(f^{-1}(Y))=Y\) for all \(Y\subseteq B\).
\end{enumerate}
\end{problemstatementbox}

\newpage
\firstdiagnosis{Homework 1, Problem 3}

\newpage
\recipecard{Homework 1, Problem 3}

\newpage
\begin{selfexplanationbox}{Homework 1, Problem 3}
\workquestion{1. Why is \(X\subseteq f^{-1}(f(X))\) always true?}{0.95in}
\workquestion{2. Where is injectivity used to prove \(f^{-1}(f(X))\subseteq X\)?}{0.95in}
\workquestion{3. To prove injectivity from the set identity, why choose \(X=\{a_1\}\)?}{0.95in}
\workquestion{4. Why is \(f(f^{-1}(Y))\subseteq Y\) always true?}{0.95in}
\workquestion{5. Where is surjectivity used to prove \(Y\subseteq f(f^{-1}(Y))\)?}{0.95in}
\end{selfexplanationbox}

\newpage
\begin{skeletonbox}{Homework 1, Problem 3}
\textbf{Part (a), injective \(\Rightarrow\) identity.}
\workquestion{Fix \(X\subseteq A\). Prove \(X\subseteq f^{-1}(f(X))\):}{0.85in}
\workquestion{Let \(a\in f^{-1}(f(X))\). Choose \(x\in X\) with \(f(a)=f(x)\), then use injectivity:}{0.95in}

\textbf{Part (a), identity \(\Rightarrow\) injective.}
\workquestion{Assume \(f(a_1)=f(a_2)\). Let \(X=\{a_1\}\) and show \(a_2\in f^{-1}(f(X))=X\):}{0.95in}

\textbf{Part (b).}
\workquestion{For surjective \(\Rightarrow\) identity, prove both inclusions for \(Y\subseteq B\):}{0.95in}
\workquestion{For identity \(\Rightarrow\) surjective, apply the identity to \(Y=\{b\}\) for arbitrary \(b\in B\):}{0.95in}
\end{skeletonbox}

\newpage
\fullproofpage{Homework 1, Problem 3}{Use this page for part (a).}

\newpage
\fullproofpage{Homework 1, Problem 3}{Use this page for part (b).}

\newpage
\begin{solutionbox}{Student-Derived Solution (No official solution available)}
(a) Suppose first that \(f\) is injective. Let \(X\subseteq A\). If \(x\in X\), then \(f(x)\in f(X)\), so \(x\in f^{-1}(f(X))\). Hence \(X\subseteq f^{-1}(f(X))\).

Conversely, let \(a\in f^{-1}(f(X))\). Then \(f(a)\in f(X)\), so there exists \(x\in X\) such that \(f(a)=f(x)\). Since \(f\) is injective, \(a=x\in X\). Thus \(f^{-1}(f(X))\subseteq X\), so \(f^{-1}(f(X))=X\).

Now suppose \(f^{-1}(f(X))=X\) for all \(X\subseteq A\). Let \(a_1,a_2\in A\) and assume \(f(a_1)=f(a_2)\). Set \(X=\{a_1\}\). Then \(f(a_2)=f(a_1)\in f(X)\), so
\[
  a_2\in f^{-1}(f(X))=X=\{a_1\}.
\]
Hence \(a_2=a_1\), and \(f\) is injective.

(b) Suppose first that \(f\) is surjective. Let \(Y\subseteq B\). For any function, \(f(f^{-1}(Y))\subseteq Y\): if \(b\in f(f^{-1}(Y))\), then \(b=f(a)\) for some \(a\in f^{-1}(Y)\), hence \(f(a)\in Y\), so \(b\in Y\).

For the reverse inclusion, let \(y\in Y\). Since \(f\) is surjective, there exists \(a\in A\) such that \(f(a)=y\). Then \(a\in f^{-1}(Y)\), so \(y=f(a)\in f(f^{-1}(Y))\). Thus \(Y\subseteq f(f^{-1}(Y))\), and \(f(f^{-1}(Y))=Y\).

Conversely, suppose \(f(f^{-1}(Y))=Y\) for all \(Y\subseteq B\). Fix \(b\in B\). Let \(Y=\{b\}\). Then
\[
  f(f^{-1}(\{b\}))=\{b\}.
\]
Thus \(f^{-1}(\{b\})\neq\varnothing\), so there exists \(a\in A\) with \(f(a)=b\). Since \(b\in B\) was arbitrary, \(f\) is surjective.
\end{solutionbox}

\newpage
\reflectionpage{Homework 1, Problem 3}

% ---------------------------------------------------------------------------
% HW1 Problem 4
% ---------------------------------------------------------------------------
\newpage
\subsection{Problem 4}

\begin{problemstatementbox}{Homework 1, Problem 4}
Let \(f:A\to B\) be a function. Prove that the following two statements are equivalent:
\begin{enumerate}[label=(\alph*)]
  \item The function \(f\) is surjective.
  \item For any set \(C\) and for any functions \(g:B\to C\) and \(h:B\to C\) such that \(g\circ f=h\circ f\), we have \(g=h\).
\end{enumerate}
\end{problemstatementbox}

\newpage
\firstdiagnosis{Homework 1, Problem 4}

\newpage
\recipecard{Homework 1, Problem 4}

\newpage
\begin{selfexplanationbox}{Homework 1, Problem 4}
\workquestion{1. Why does surjectivity let you test equality of \(g,h:B\to C\) after composing with \(f\)?}{1.0in}
\workquestion{2. In the forward direction, after fixing \(b\in B\), what element of \(A\) do you need?}{0.9in}
\workquestion{3. To prove the reverse direction, what does it mean for \(f\) not to be surjective?}{0.9in}
\workquestion{4. Why choose \(C=\{0,1\}\) in the counterexample construction?}{0.95in}
\workquestion{5. How can \(g\circ f=h\circ f\) hold while \(g\neq h\)?}{0.95in}
\end{selfexplanationbox}

\newpage
\begin{skeletonbox}{Homework 1, Problem 4}
\textbf{Surjective \(\Rightarrow\) cancellation.}
\workquestion{Assume \(g\circ f=h\circ f\). Fix \(b\in B\), choose \(a\in A\) with \(f(a)=b\):}{0.95in}
\workquestion{Compute \(g(b)=g(f(a))=(g\circ f)(a)=\cdots=h(b)\):}{0.85in}

\textbf{Cancellation \(\Rightarrow\) surjective.}
\workquestion{Assume \(f\) is not surjective; choose \(b_0\in B\setminus f(A)\):}{0.85in}
\workquestion{Set \(C=\{0,1\}\), define \(g\) constantly \(0\), and define \(h\) to differ only at \(b_0\):}{1.0in}
\workquestion{Show \(g\circ f=h\circ f\) because \(f(a)\neq b_0\) for all \(a\):}{0.95in}
\workquestion{Contradict the cancellation property and conclude \(f\) is surjective:}{0.85in}
\end{skeletonbox}

\newpage
\fullproofpage{Homework 1, Problem 4}{Use this page for the forward direction.}

\newpage
\fullproofpage{Homework 1, Problem 4}{Use this page for the reverse direction.}

\newpage
\begin{solutionbox}{Student-Derived Solution (No official solution available)}
\(\Rightarrow\) Assume that \(f\) is surjective. Let \(C\) be any set, and let \(g,h:B\to C\) satisfy \(g\circ f=h\circ f\). We show \(g=h\). Fix \(b\in B\). Since \(f\) is surjective, there exists \(a\in A\) such that \(f(a)=b\). Then
\[
  g(b)=g(f(a))=(g\circ f)(a)=(h\circ f)(a)=h(f(a))=h(b).
\]
Since \(b\in B\) was arbitrary, \(g=h\).

\(\Leftarrow\) Assume the cancellation property in (b). Suppose, for contradiction, that \(f\) is not surjective. Then there exists \(b_0\in B\setminus f(A)\). Let \(C=\{0,1\}\), and define \(g,h:B\to C\) by
\[
  g(b)=0 \quad \text{for all } b\in B,
\]
and
\[
  h(b)=
  \begin{cases}
    1, & b=b_0,\\
    0, & b\neq b_0.
  \end{cases}
\]
Then \(g\neq h\), since \(g(b_0)=0\) but \(h(b_0)=1\). However, for every \(a\in A\), we have \(f(a)\neq b_0\), so
\[
  (g\circ f)(a)=g(f(a))=0=h(f(a))=(h\circ f)(a).
\]
Thus \(g\circ f=h\circ f\), and by the assumed property this forces \(g=h\), a contradiction. Therefore \(f\) is surjective.
\end{solutionbox}

\newpage
\reflectionpage{Homework 1, Problem 4}

% ---------------------------------------------------------------------------
% HW1 Problem 5
% ---------------------------------------------------------------------------
\newpage
\subsection{Problem 5}

\begin{problemstatementbox}{Homework 1, Problem 5}
Let \(f:B\to C\) be a function. Prove that the following two statements are equivalent:
\begin{enumerate}[label=(\alph*)]
  \item The function \(f\) is injective.
  \item For any set \(A\) and for any functions \(g:A\to B\) and \(h:A\to B\) such that \(f\circ g=f\circ h\), we have \(g=h\).
\end{enumerate}
\end{problemstatementbox}

\newpage
\firstdiagnosis{Homework 1, Problem 5}

\newpage
\recipecard{Homework 1, Problem 5}

\newpage
\begin{selfexplanationbox}{Homework 1, Problem 5}
\workquestion{1. Why is this a left-cancellation property for injective maps?}{0.95in}
\workquestion{2. In the forward direction, what equality do we get after evaluating at \(a\in A\)?}{0.9in}
\workquestion{3. Where is injectivity used to conclude \(g(a)=h(a)\)?}{0.85in}
\workquestion{4. In the reverse direction, why does choosing \(A=\{1\}\) suffice?}{0.95in}
\workquestion{5. How do the test maps \(g(1)=b_1\), \(h(1)=b_2\) force \(b_1=b_2\)?}{0.95in}
\end{selfexplanationbox}

\newpage
\begin{skeletonbox}{Homework 1, Problem 5}
\textbf{Injective \(\Rightarrow\) cancellation.}
\workquestion{Assume \(f\circ g=f\circ h\). Fix \(a\in A\) and write the equality after evaluating at \(a\):}{0.9in}
\workquestion{Use injectivity of \(f\) to conclude \(g(a)=h(a)\):}{0.85in}
\workquestion{Since \(a\) was arbitrary, conclude \(g=h\):}{0.75in}

\textbf{Cancellation \(\Rightarrow\) injective.}
\workquestion{Fix \(b_1,b_2\in B\) with \(f(b_1)=f(b_2)\):}{0.8in}
\workquestion{Define \(g,h:\{1\}\to B\) by \(g(1)=b_1\), \(h(1)=b_2\):}{0.9in}
\workquestion{Show \(f\circ g=f\circ h\), apply the property, and conclude \(b_1=b_2\):}{0.95in}
\end{skeletonbox}

\newpage
\fullproofpage{Homework 1, Problem 5}{Use this page for both directions.}

\newpage
\begin{solutionbox}{Official Solution (Professor)}
\((a)\Rightarrow(b)\): Suppose that \(f\) is injective. We want to show that property (b) is true.

Let \(A\) be a set and \(g,h:A\to B\) be functions such that \(f\circ g=f\circ h\). This means that for any \(a\in A\), we have
\[
  f(g(a))=f(h(a)).
\]
In other words, \(g(a)\) and \(h(a)\) are two points in \(B\) which \(f\) maps to the same point in \(C\). As \(f\) is injective, we must have
\[
  g(a)=h(a).
\]
This holds for any \(a\in A\), and so the functions \(g\) and \(h\) are equal.

\((a)\Leftarrow(b)\): Suppose that \(f\) satisfies property (b). We want to show that \(f\) is injective.

Fix \(b_1,b_2\in B\) such that \(f(b_1)=f(b_2)\). Consider the functions \(g,h:\{1\}\to B\) given by
\[
  g(1)=b_1 \qquad \text{and} \qquad h(1)=b_2.
\]
Then we have \(f\circ g=f\circ h\), since
\[
  f(g(1))=f(b_1)=f(b_2)=f(h(1)).
\]
Therefore, by property (b) we must have \(g=h\), and so
\[
  b_1=g(1)=h(1)=b_2.
\]
Thus \(b_1=b_2\), as desired.
\end{solutionbox}

\newpage
\reflectionpage{Homework 1, Problem 5}

% ---------------------------------------------------------------------------
% HW1 Problem 6
% ---------------------------------------------------------------------------
\newpage
\subsection{Problem 6}

\begin{problemstatementbox}{Homework 1, Problem 6}
Let \(f:A\to B\) be a function. Suppose that there exists a function \(g:B\to A\) such that \((g\circ f)(a)=a\) for all \(a\in A\) and \((f\circ g)(b)=b\) for all \(b\in B\). Prove that \(f\) is bijective.
\end{problemstatementbox}

\newpage
\firstdiagnosis{Homework 1, Problem 6}

\newpage
\recipecard{Homework 1, Problem 6}

\newpage
\begin{selfexplanationbox}{Homework 1, Problem 6}
\workquestion{1. What do the equations \(g\circ f=\mathrm{id}_A\) and \(f\circ g=\mathrm{id}_B\) mean pointwise?}{1.0in}
\workquestion{2. To prove injectivity, what should we assume about \(f(a_1)\) and \(f(a_2)\)?}{0.9in}
\workquestion{3. Why does applying \(g\) to both sides prove \(a_1=a_2\)?}{0.9in}
\workquestion{4. To prove surjectivity, how should an arbitrary \(b\in B\) be written as \(f(a)\)?}{0.95in}
\workquestion{5. Why do injective plus surjective imply bijective?}{0.85in}
\end{selfexplanationbox}

\newpage
\begin{skeletonbox}{Homework 1, Problem 6}
\textbf{Injectivity.}
\workquestion{Assume \(f(a_1)=f(a_2)\). Apply \(g\) and use \(g(f(a))=a\):}{0.95in}
\workquestion{Conclude \(a_1=a_2\), so \(f\) is injective:}{0.75in}

\textbf{Surjectivity.}
\workquestion{Fix \(b\in B\). Set \(a=\)}{0.75in}
\workquestion{Use \(f(g(b))=b\) to show \(b\in f(A)\):}{0.85in}
\workquestion{Conclude every \(b\in B\) is hit, and finish:}{0.75in}
\end{skeletonbox}

\newpage
\fullproofpage{Homework 1, Problem 6}{Use this page for the complete proof.}

\newpage
\begin{solutionbox}{Official Solution (Professor)}
First, we will show that \(f\) is injective. Suppose that \(f(a_1)=f(a_2)\) for some \(a_1,a_2\in A\). Then
\[
  a_1=g(f(a_1))=g(f(a_2))=a_2.
\]

Next, we will show that \(f\) is surjective. Fix \(b\in B\). Set \(a=g(b)\). Then \(a\in A\) and
\[
  f(a)=f(g(b))=b.
\]

Together, we conclude that \(f\) is bijective.
\end{solutionbox}

\newpage
\reflectionpage{Homework 1, Problem 6}

\section{Homework 2}

% ---------------------------------------------------------------------------
% HW2 Problem 1
% ---------------------------------------------------------------------------
\newpage
\subsection{Problem 1}

\begin{problemstatementbox}{Homework 2, Problem 1}
Prove that if \(A\) and \(B\) are both countable sets and \(A\cap B=\varnothing\), then \(A\cup B\) is countable. (Hint: We proved that \(\Z\) is countable in class. Try using a similar idea here.)
\end{problemstatementbox}

\newpage
\firstdiagnosis{Homework 2, Problem 1}

\newpage
\recipecard{Homework 2, Problem 1}

\newpage
\begin{selfexplanationbox}{Homework 2, Problem 1}
\workquestion{1. What does it mean for \(A\) and \(B\) to be countable in terms of bijections with \(\N\)?}{0.95in}
\workquestion{2. Why does disjointness make the odd/even coding well-defined?}{0.95in}
\workquestion{3. Which natural numbers should encode elements of \(A\), and which should encode elements of \(B\)?}{0.95in}
\workquestion{4. How is the inverse map \(\phi:\N\to A\cup B\) defined on odd and even integers?}{1.0in}
\workquestion{5. Why do the two composition identities prove bijectivity?}{0.9in}
\end{selfexplanationbox}

\newpage
\begin{skeletonbox}{Homework 2, Problem 1}
\workquestion{Choose bijections \(g:A\to\N\) and \(h:B\to\N\):}{0.8in}
\workquestion{Define \(f:A\cup B\to\N\) by odd values on \(A\) and even values on \(B\):}{0.95in}
\workquestion{Define \(\phi:\N\to A\cup B\) separately for odd and even \(n\):}{1.0in}
\workquestion{Verify \(f\circ\phi=\mathrm{id}_{\N}\):}{0.95in}
\workquestion{Verify \(\phi\circ f=\mathrm{id}_{A\cup B}\):}{0.95in}
\workquestion{Apply HW1\#6 and conclude \(A\cup B\) is countable:}{0.75in}
\end{skeletonbox}

\newpage
\fullproofpage{Homework 2, Problem 1}{Use this page for the complete proof.}

\newpage
\begin{solutionbox}{Official Solution (Professor)}
As \(A\) and \(B\) are countable, there exist functions \(g:A\to\N\) and \(h:B\to\N\) that are bijective.

Define the function \(f:A\cup B\to\N\) by
\[
  f(x)=
  \begin{cases}
    2g(x)-1, & x\in A,\\
    2h(x), & x\in B.
  \end{cases}
\]
As \(A\cap B=\varnothing\), these two cases are mutually exclusive and so \(f\) is indeed a function.

We want to show that \(f\) is bijective. To this end, consider the function \(\phi:\N\to A\cup B\) given by
\[
  \phi(n)=
  \begin{cases}
    g^{-1}\left(\frac{n+1}{2}\right), & n\text{ odd},\\
    h^{-1}\left(\frac n2\right), & n\text{ even}.
  \end{cases}
\]
For \(n\in\N\) odd, note that \(\phi(n)=g^{-1}\left(\frac{n+1}{2}\right)\in A\), and so
\[
  (f\circ\phi)(n)
  =
  f\left(g^{-1}\left(\frac{n+1}{2}\right)\right)
  =
  2g\left(g^{-1}\left(\frac{n+1}{2}\right)\right)-1
  =
  n.
\]
Similarly, for \(n\in\N\) even, \(\phi(n)=h^{-1}(\frac n2)\in B\), and
\[
  (f\circ\phi)(n)
  =
  f\left(h^{-1}\left(\frac n2\right)\right)
  =
  2h\left(h^{-1}\left(\frac n2\right)\right)
  =
  n.
\]
Together, \((f\circ\phi)(n)=n\) for all \(n\in\N\).

On the other hand, for \(x\in A\), \(f(x)=2g(x)-1\) is odd, and so
\[
  (\phi\circ f)(x)
  =
  \phi(2g(x)-1)
  =
  g^{-1}\left(\frac{2g(x)-1+1}{2}\right)
  =
  g^{-1}(g(x))
  =
  x.
\]
Likewise, for \(x\in B\), \(f(x)=2h(x)\) is even, and so
\[
  (\phi\circ f)(x)
  =
  \phi(2h(x))
  =
  h^{-1}\left(\frac{2h(x)}{2}\right)
  =
  h^{-1}(h(x))
  =
  x.
\]
Together, \((\phi\circ f)(x)=x\) for all \(x\in A\cup B\). Combining the previous two paragraphs, \(f\) is bijective by HW1\#6.
\end{solutionbox}

\newpage
\reflectionpage{Homework 2, Problem 1}

% ---------------------------------------------------------------------------
% HW2 Problem 2
% ---------------------------------------------------------------------------
\newpage
\subsection{Problem 2}

\begin{problemstatementbox}{Homework 2, Problem 2}
Let \(A_n\) be a countable set for each \(n\in\N\), and suppose that \(A_m\cap A_n=\varnothing\) whenever \(m\neq n\). Prove that the union \(A=A_1\cup A_2\cup A_3\cup\cdots\) is countable. (Hint: We proved that \(\N\times\N\) is countable in class. Try using a similar idea here.)
\end{problemstatementbox}

\newpage
\firstdiagnosis{Homework 2, Problem 2}

\newpage
\recipecard{Homework 2, Problem 2}

\newpage
\begin{selfexplanationbox}{Homework 2, Problem 2}
\workquestion{1. For each \(A_n\), what enumeration or bijection should be chosen?}{0.95in}
\workquestion{2. Why does an element of the union have a unique first coordinate \(n\)?}{0.95in}
\workquestion{3. What map \(h:\N\times\N\to A\) naturally lists all elements?}{0.95in}
\workquestion{4. How does disjointness prove injectivity of \(h\)?}{0.95in}
\workquestion{5. How does countability of \(\N\times\N\) finish the proof?}{0.9in}
\end{selfexplanationbox}

\newpage
\begin{skeletonbox}{Homework 2, Problem 2}
\workquestion{For each \(n\), choose a bijection \(f_n:\N\to A_n\):}{0.8in}
\workquestion{Define \(h:\N\times\N\to A\) by \(h(n,k)=\)}{0.8in}
\workquestion{Prove \(h\) is surjective by taking \(x\in A\) and locating \(x\in A_n\):}{0.95in}
\workquestion{Prove \(h\) is injective using pairwise disjointness and injectivity of \(f_n\):}{1.0in}
\workquestion{Compose with a bijection \(\N\to\N\times\N\) and conclude \(A\) is countable:}{0.9in}
\end{skeletonbox}

\newpage
\fullproofpage{Homework 2, Problem 2}{Use this page for the complete proof.}

\newpage
\begin{solutionbox}{Student-Derived Solution (No official solution available)}
For each \(n\in\N\), since \(A_n\) is countable, choose a bijection
\[
  f_n:\N\to A_n.
\]
Define \(h:\N\times\N\to A\) by
\[
  h(n,k)=f_n(k).
\]
We show that \(h\) is bijective.

Surjectivity. Let \(x\in A\). Since \(A=\bigcup_{n\in\N}A_n\), there exists \(n\in\N\) such that \(x\in A_n\). Since \(f_n\) is surjective, there exists \(k\in\N\) such that \(f_n(k)=x\). Hence \(h(n,k)=x\).

Injectivity. Suppose \(h(n,k)=h(m,\ell)\). Then
\[
  f_n(k)=f_m(\ell).
\]
The left side lies in \(A_n\) and the right side lies in \(A_m\). If \(n\neq m\), this contradicts the assumption that \(A_n\cap A_m=\varnothing\). Hence \(n=m\). With \(n=m\), we have \(f_n(k)=f_n(\ell)\), and since \(f_n\) is injective, \(k=\ell\). Therefore \((n,k)=(m,\ell)\), so \(h\) is injective.

Thus \(h\) is a bijection \(\N\times\N\to A\). Since \(\N\times\N\) is countable, there exists a bijection \(\phi:\N\to\N\times\N\). Then \(h\circ\phi:\N\to A\) is a bijection, so \(A\) is countable.
\end{solutionbox}

\newpage
\reflectionpage{Homework 2, Problem 2}

% ---------------------------------------------------------------------------
% HW2 Problem 3
% ---------------------------------------------------------------------------
\newpage
\subsection{Problem 3}

\begin{problemstatementbox}{Homework 2, Problem 3}
Recall that given two sets \(A\) and \(B\), we write \(A\sim B\) if and only if there exists a bijective function \(f:A\to B\). Prove that \(\sim\) is an equivalence relation.
\end{problemstatementbox}

\newpage
\firstdiagnosis{Homework 2, Problem 3}

\newpage
\recipecard{Homework 2, Problem 3}

\newpage
\begin{selfexplanationbox}{Homework 2, Problem 3}
\workquestion{1. What are the three properties of an equivalence relation?}{0.9in}
\workquestion{2. Which bijection proves \(A\sim A\)?}{0.85in}
\workquestion{3. Why does a bijection \(A\to B\) give a bijection \(B\to A\)?}{0.95in}
\workquestion{4. Why is the composition of bijections bijective?}{0.95in}
\workquestion{5. How do these three facts match reflexivity, symmetry, and transitivity?}{0.9in}
\end{selfexplanationbox}

\newpage
\begin{skeletonbox}{Homework 2, Problem 3}
\workquestion{Reflexivity: define the identity map \(f:A\to A\) and show it is bijective:}{0.9in}
\workquestion{Symmetry: start with a bijection \(g:A\to B\) and use \(g^{-1}:B\to A\):}{0.95in}
\workquestion{Transitivity: compose bijections \(h_1:A\to B\) and \(h_2:B\to C\):}{0.95in}
\workquestion{State that \(\sim\) is reflexive, symmetric, and transitive:}{0.75in}
\end{skeletonbox}

\newpage
\fullproofpage{Homework 2, Problem 3}{Use this page for the complete proof.}

\newpage
\begin{solutionbox}{Official Solution (Professor)}
Reflexivity: Let \(A\) be a set. We want to show that \(A\sim A\). Consider the function \(f:A\to A\) given by \(f(a)=a\). Notice that \((f\circ f)(a)=a\) for all \(a\in A\), and so \(f\) is bijective by HW1\#6. Therefore \(A\sim A\).

Symmetry: Suppose that \(A,B\) are sets such that \(A\sim B\). Then there exists a bijective function \(g:A\to B\). Consider the function \(g^{-1}:B\to A\). By Lecture 3 we know that \((g^{-1}\circ g)(a)=a\) for all \(a\in A\) and \((g\circ g^{-1})(b)=b\) for all \(b\in B\), and so \(g^{-1}\) is bijective by HW1\#6. Therefore \(B\sim A\).

Transitivity: Suppose that \(A,B,C\) are sets such that \(A\sim B\) and \(B\sim C\). Then there exist bijective functions \(h_1:A\to B\) and \(h_2:B\to C\). By Lecture 3, the composition \(h_2\circ h_1:A\to C\) is also bijective, and so \(A\sim C\).
\end{solutionbox}

\newpage
\reflectionpage{Homework 2, Problem 3}

% ---------------------------------------------------------------------------
% HW2 Problem 4
% ---------------------------------------------------------------------------
\newpage
\subsection{Problem 4}

\begin{problemstatementbox}{Homework 2, Problem 4}
For \(a,b\in\Z\), define \(a\sim b\) if and only if \(b-a\) is divisible by \(2\).
\begin{enumerate}[label=(\alph*)]
  \item Prove that \(\sim\) is an equivalence relation on \(\Z\).
  \item Find the equivalence classes.
\end{enumerate}
\end{problemstatementbox}

\newpage
\firstdiagnosis{Homework 2, Problem 4}

\newpage
\recipecard{Homework 2, Problem 4}

\newpage
\begin{selfexplanationbox}{Homework 2, Problem 4}
\workquestion{1. What does it mean for an integer to be divisible by \(2\)?}{0.9in}
\workquestion{2. Why is \(a-a=0\) enough for reflexivity?}{0.85in}
\workquestion{3. Why does \(b-a=2k\) imply \(a-b\) is divisible by \(2\)?}{0.9in}
\workquestion{4. How does adding \(b-a\) and \(c-b\) prove transitivity?}{0.95in}
\workquestion{5. Why are the equivalence classes exactly the even and odd integers?}{0.95in}
\end{selfexplanationbox}

\newpage
\begin{skeletonbox}{Homework 2, Problem 4}
\textbf{Part (a).}
\workquestion{Reflexive: show \(a-a\) is divisible by \(2\):}{0.75in}
\workquestion{Symmetric: assume \(b-a=2k\), then write \(a-b=\)}{0.8in}
\workquestion{Transitive: assume \(b-a=2k\) and \(c-b=2\ell\), then compute \(c-a\):}{0.9in}

\textbf{Part (b).}
\workquestion{Show \([0]=\{2k:k\in\Z\}\):}{0.9in}
\workquestion{Show \([1]=\{2k+1:k\in\Z\}\):}{0.9in}
\workquestion{Explain why these are the only two classes:}{0.75in}
\end{skeletonbox}

\newpage
\fullproofpage{Homework 2, Problem 4}{Use this page for part (a).}

\newpage
\fullproofpage{Homework 2, Problem 4}{Use this page for part (b).}

\newpage
\begin{solutionbox}{Student-Derived Solution (No official solution available)}
(a) We show that \(\sim\) is reflexive, symmetric, and transitive on \(\Z\).

Reflexive. Let \(a\in\Z\). Then \(a-a=0=2\cdot0\), so \(a-a\) is divisible by \(2\). Hence \(a\sim a\).

Symmetric. Suppose \(a\sim b\). Then \(b-a=2k\) for some \(k\in\Z\). Hence
\[
  a-b=-(b-a)=-2k=2(-k),
\]
so \(a-b\) is divisible by \(2\). Thus \(b\sim a\).

Transitive. Suppose \(a\sim b\) and \(b\sim c\). Then \(b-a=2k\) and \(c-b=2\ell\) for some \(k,\ell\in\Z\). Adding,
\[
  c-a=(c-b)+(b-a)=2\ell+2k=2(k+\ell),
\]
so \(c-a\) is divisible by \(2\). Hence \(a\sim c\).

Therefore \(\sim\) is an equivalence relation.

(b) The equivalence classes are the parity classes. First,
\[
  [0]=\{n\in\Z:n-0\text{ is divisible by }2\}=\{2k:k\in\Z\}=2\Z.
\]
Similarly,
\[
  [1]=\{n\in\Z:n-1\text{ is divisible by }2\}=\{2k+1:k\in\Z\}=2\Z+1.
\]
Every integer is either even or odd, so these are the only equivalence classes.
\end{solutionbox}

\newpage
\reflectionpage{Homework 2, Problem 4}

% ---------------------------------------------------------------------------
% HW2 Problem 5
% ---------------------------------------------------------------------------
\newpage
\subsection{Problem 5}

\begin{problemstatementbox}{Homework 2, Problem 5}
Let \((F,+,\cdot)\) be a field.
\begin{enumerate}[label=(\alph*)]
  \item Prove that for any \(x\in F\), the function \(f:F\to F\), \(f(y)=x+y\) is bijective.
  \item Prove that for any \(x\in F\setminus\{0\}\), the function \(g:F\to F\), \(g(y)=x\cdot y\) is bijective.
\end{enumerate}
\end{problemstatementbox}

\newpage
\firstdiagnosis{Homework 2, Problem 5}

\newpage
\recipecard{Homework 2, Problem 5}

\newpage
\begin{selfexplanationbox}{Homework 2, Problem 5}
\workquestion{1. What operation undoes adding a fixed \(x\)?}{0.85in}
\workquestion{2. Which field axiom gives the additive inverse \(-x\)?}{0.85in}
\workquestion{3. What operation undoes multiplying by a fixed nonzero \(x\)?}{0.9in}
\workquestion{4. Which field axiom gives \(x^{-1}\)?}{0.85in}
\workquestion{5. Why does exhibiting a two-sided inverse function prove bijectivity?}{0.95in}
\end{selfexplanationbox}

\newpage
\begin{skeletonbox}{Homework 2, Problem 5}
\textbf{Part (a).}
\workquestion{Define \(h:F\to F\) by \(h(z)=\)}{0.75in}
\workquestion{Compute \((f\circ h)(z)\) using associativity, inverse, and identity axioms:}{0.95in}
\workquestion{Compute \((h\circ f)(y)\):}{0.9in}
\workquestion{Conclude \(f\) is bijective by HW1\#6:}{0.75in}

\textbf{Part (b).}
\workquestion{Define \(h:F\to F\) by \(h(z)=\)}{0.75in}
\workquestion{Compute \((g\circ h)(z)\) and \((h\circ g)(y)\):}{1.0in}
\workquestion{Conclude \(g\) is bijective:}{0.75in}
\end{skeletonbox}

\newpage
\fullproofpage{Homework 2, Problem 5}{Use this page for part (a).}

\newpage
\fullproofpage{Homework 2, Problem 5}{Use this page for part (b).}

\newpage
\begin{solutionbox}{Official Solution (Professor)}
(a) Fix \(x\in F\). By (A5), there exists an additive inverse \(-x\in F\) for \(x\), so that
\[
  x+(-x)=0=-x+x.
\]
Consider the function \(h:F\to F\) given by \(h(z)=-x+z\). Then, for any \(z\in F\),
\[
  (f\circ h)(z)=f(-x+z)=x+(-x+z)
  =
  (x+(-x))+z
  =
  0+z
  =
  z.
\]
Similarly, for \(y\in F\),
\[
  (h\circ f)(y)=h(x+y)=-x+(x+y)
  =
  (-x+x)+y
  =
  0+y
  =
  y.
\]
Therefore \(f\) is bijective by HW1\#6.

(b) Fix \(x\in F\setminus\{0\}\). As \(x\neq0\), there exists a multiplicative inverse \(x^{-1}\in F\) for \(x\), so that
\[
  x\cdot x^{-1}=1=x^{-1}\cdot x.
\]
Consider the function \(h:F\to F\) given by \(h(z)=x^{-1}\cdot z\). Then, for any \(z\in F\),
\[
  (g\circ h)(z)=g(x^{-1}\cdot z)=x\cdot(x^{-1}\cdot z)
  =
  (x\cdot x^{-1})\cdot z
  =
  1\cdot z
  =
  z.
\]
Similarly, for \(y\in F\),
\[
  (h\circ g)(y)=h(x\cdot y)=x^{-1}\cdot(x\cdot y)
  =
  (x^{-1}\cdot x)\cdot y
  =
  1\cdot y
  =
  y.
\]
Therefore \(g\) is bijective by HW1\#6.
\end{solutionbox}

\newpage
\reflectionpage{Homework 2, Problem 5}

% ---------------------------------------------------------------------------
% HW2 Problem 6
% ---------------------------------------------------------------------------
\newpage
\subsection{Problem 6}

\begin{problemstatementbox}{Homework 2, Problem 6}
Let \((F,+,\cdot)\) be a field. Suppose that \(F=\{0,1,a,b\}\) is a set with exactly four elements, where \(0\) and \(1\) denote the identities for \(+\) and \(\cdot\) respectively, and \(a,b\) denote the remaining two elements of \(F\). Fill in the addition and multiplication tables below, and justify your answer.
\[
\begin{array}{c|cccc}
+&0&1&a&b\\ \hline
0&&&&\\
1&&&&\\
a&&&&\\
b&&&&
\end{array}
\qquad
\begin{array}{c|cccc}
\cdot&0&1&a&b\\ \hline
0&&&&\\
1&&&&\\
a&&&&\\
b&&&&
\end{array}
\]
(Hint: Using the previous problem, show that each row and column in the addition table contains every element of \(F\) exactly once, like a Sudoku puzzle. Show that the same is true for the nonzero rows and columns of the multiplication table. Note that for each entry, there is only one correct answer.)
\end{problemstatementbox}

\newpage
\firstdiagnosis{Homework 2, Problem 6}

\newpage
\recipecard{Homework 2, Problem 6}

\newpage
\begin{selfexplanationbox}{Homework 2, Problem 6}
\workquestion{1. How does Problem 5 imply each addition row/column is a permutation of \(F\)?}{0.95in}
\workquestion{2. How does Problem 5 imply each nonzero multiplication row/column is a permutation of \(F\setminus\{0\}\)?}{1.0in}
\workquestion{3. Which entries are forced immediately by additive and multiplicative identities?}{0.95in}
\workquestion{4. Why must \(a\cdot b=1\) in the multiplication table?}{0.95in}
\workquestion{5. How does distributivity rule out the wrong addition-table entries?}{1.0in}
\end{selfexplanationbox}

\newpage
\begin{skeletonbox}{Homework 2, Problem 6}
\workquestion{Fill the \(0\) and \(1\) rows/columns in the multiplication table:}{0.95in}
\workquestion{Use the nonzero row/column permutation rule to force \(a\cdot b=1\), \(a^2=b\), and \(b^2=a\):}{1.0in}
\workquestion{Fill the \(0\) row/column in the addition table:}{0.85in}
\workquestion{Use distributivity to rule out \(1+a=0\), hence force \(1+a=b\):}{1.0in}
\workquestion{Similarly force \(1+b=a\) and \(a+b=1\):}{0.95in}
\workquestion{Write the completed tables:}{0.9in}
\end{skeletonbox}

\newpage
\fullproofpage{Homework 2, Problem 6}{Use this page for the multiplication table.}

\newpage
\fullproofpage{Homework 2, Problem 6}{Use this page for the addition table and justification.}

\newpage
\begin{solutionbox}{Official Solution (Professor)}
There are many ways to do this problem. Here's one option that works.

Step 1: Start with the multiplication table. We can fill in the \(0\) and \(1\) rows and columns using what we know from lecture:
\[
\begin{array}{c|cccc}
\cdot&0&1&a&b\\ \hline
0&0&0&0&0\\
1&0&1&a&b\\
a&0&a&&\\
b&0&b&&
\end{array}
\]

Step 2: The only possibility for the \(a\cdot b\) entry is \(1\), since the other three elements already appear in this row and column. This in turn leaves only one option for the remaining entries:
\[
\begin{array}{c|cccc}
\cdot&0&1&a&b\\ \hline
0&0&0&0&0\\
1&0&1&a&b\\
a&0&a&b&1\\
b&0&b&1&a
\end{array}
\]

Step 3: Now turn to the addition table. We can fill in the \(0\) row and column using (A4):
\[
\begin{array}{c|cccc}
+&0&1&a&b\\ \hline
0&0&1&a&b\\
1&1&&&\\
a&a&&&\\
b&b&&&
\end{array}
\]

Step 4: Consider the entry for \(1+a\). There are two elements that we can choose from, \(0\) and \(b\), since the other two elements already appear in this row and column. Suppose \(1+a=0\). This forces
\[
  0=0\cdot b=(1+a)\cdot b=1\cdot b+a\cdot b=b+1,
\]
by distributivity and the multiplication table. But this would mean that \(0\) appears twice in a row/column, which is a contradiction. Therefore \(1+a=b\). Similarly, if \(1+b=0\), then
\[
  0=0\cdot b=(1+b)\cdot a=1\cdot a+b\cdot a=a+1,
\]
which contradicts the current entry for \(a+1\). So \(1+b=a\).

Step 5: Consider the entry for \(a+b\). There are two elements that we can choose from, \(0\) and \(1\), since the other two elements already appear in this row and column. Suppose \(a+b=0\). This forces
\[
  0=0\cdot b=(a+b)\cdot b=a\cdot b+b\cdot b=1+a,
\]
by distributivity and the multiplication table. This contradicts the current entry for \(1+a\), so \(a+b=1\). This leaves the completed addition table:
\[
\begin{array}{c|cccc}
+&0&1&a&b\\ \hline
0&0&1&a&b\\
1&1&0&b&a\\
a&a&b&0&1\\
b&b&a&1&0
\end{array}
\]
\end{solutionbox}

\newpage
\reflectionpage{Homework 2, Problem 6}

\section{Homework 3}

% ---------------------------------------------------------------------------
% HW3 Problem 1
% ---------------------------------------------------------------------------
\newpage
\subsection{Problem 1}

\begin{problemstatementbox}{Homework 3, Problem 1}
Let \((F,+,\cdot)\) be a field.
\begin{enumerate}[label=(\alph*)]
  \item Prove that the multiplicative identity is unique: If \(b\in F\) satisfies \(b\cdot a=a\cdot b=a\) for all \(a\in F\), then \(b=1\).
  \item Prove that multiplicative inverses are unique: If \(a\in F\setminus\{0\}\) and \(b\in F\) satisfy \(b\cdot a=a\cdot b=1\), then \(b=a^{-1}\).
  \item Prove that if \(a\cdot b=a\cdot c\) and \(a\neq0\), then \(b=c\).
  \item Prove that if \(a,b\in F\) then \(-a+(-b)=-(a+b)\).
  \item Prove that if \(a,b\in F\) and \(a,b\neq0\), then \(a^{-1}\cdot b^{-1}=(ab)^{-1}\).
\end{enumerate}
\end{problemstatementbox}

\newpage
\firstdiagnosis{Homework 3, Problem 1}

\newpage
\recipecard{Homework 3, Problem 1}

\newpage
\begin{selfexplanationbox}{Homework 3, Problem 1}
\workquestion{1. In part (a), what two ways can you compute \(b\cdot1\)?}{0.9in}
\workquestion{2. In part (b), why is multiplying by \(a^{-1}\) legitimate?}{0.9in}
\workquestion{3. In part (c), what field axioms justify cancellation?}{0.9in}
\workquestion{4. In part (d), how do you prove an element is the additive inverse of \(a+b\)?}{0.95in}
\workquestion{5. In part (e), what product must be checked to show \(a^{-1}b^{-1}\) is \((ab)^{-1}\)?}{0.95in}
\end{selfexplanationbox}

\newpage
\begin{skeletonbox}{Homework 3, Problem 1}
\textbf{Parts (a)--(c).}
\workquestion{For (a), use the usual identity with \(a=b\) and the alternative identity with \(a=1\):}{0.95in}
\workquestion{For (b), start from \(b\cdot a=1\) and multiply on the right by \(a^{-1}\):}{0.95in}
\workquestion{For (c), multiply \(a\cdot b=a\cdot c\) on the left by \(a^{-1}\):}{0.9in}

\textbf{Parts (d)--(e).}
\workquestion{For (d), compute \(((-a)+(-b))+(a+b)\) and rearrange using commutativity/associativity:}{1.0in}
\workquestion{For (e), compute \((a^{-1}b^{-1})(ab)\) and use commutativity/associativity:}{1.0in}
\workquestion{Invoke uniqueness from earlier parts where needed:}{0.75in}
\end{skeletonbox}

\newpage
\fullproofpage{Homework 3, Problem 1}{Use this page for parts (a)--(c).}

\newpage
\fullproofpage{Homework 3, Problem 1}{Use this page for parts (d)--(e).}

\newpage
\begin{solutionbox}{Official Solution (Professor)}
(a) We know that \(1,b\in F\) satisfy:
\begin{itemize}
  \item \(a\cdot1=1\cdot a=a\) for all \(a\in F\);
  \item \(a\cdot b=b\cdot a=a\) for all \(a\in F\).
\end{itemize}
Taking \(a=b\) in the first bullet point, we see that
\[
  b\cdot1=b.
\]
Taking \(a=1\) in the second bullet point, we obtain
\[
  b\cdot1=1.
\]
Comparing the previous two equalities, we conclude \(b=1\).

(b) As \(a\neq0\), by (M5) there exists \(a^{-1}\in F\) so that
\[
  a\cdot a^{-1}=a^{-1}\cdot a=1.
\]
Multiplying \(b\cdot a=1\) by \(a^{-1}\) on the right, we obtain
\[
  (b\cdot a)\cdot a^{-1}=1\cdot a^{-1}
  \Longrightarrow
  b\cdot(a\cdot a^{-1})=a^{-1}
  \Longrightarrow
  b\cdot1=a^{-1}
  \Longrightarrow
  b=a^{-1}.
\]

(c) As \(a\neq0\), by (M5) there exists \(a^{-1}\in F\) so that
\[
  a\cdot a^{-1}=a^{-1}\cdot a=1.
\]
Multiplying \(a\cdot b=a\cdot c\) by \(a^{-1}\) on the left, we obtain
\[
  a^{-1}\cdot(a\cdot b)=a^{-1}\cdot(a\cdot c)
  \Longrightarrow
  (a^{-1}\cdot a)\cdot b=(a^{-1}\cdot a)\cdot c
  \Longrightarrow
  1\cdot b=1\cdot c
  \Longrightarrow
  b=c.
\]

(d) Fix \(a,b\in F\). Then
\[
  ((-a)+(-b))+(a+b)
  =
  ((-a)+(-b))+(b+a)
\]
\[
  =
  -a+((-b)+b)+a
  =
  -a+0+a
  =
  -a+a
  =
  0.
\]
In other words, \((-a)+(-b)\) is the additive inverse of \(a+b\). As additive inverses are unique, we conclude
\[
  -a+(-b)=-(a+b).
\]

(e) Fix \(a,b\neq0\). Then
\[
  (a^{-1}\cdot b^{-1})\cdot(ab)
  =
  (a^{-1}\cdot b^{-1})\cdot(ba)
  =
  a^{-1}\cdot(b^{-1}b)\cdot a
  =
  a^{-1}\cdot1\cdot a
  =
  a^{-1}\cdot a
  =
  1.
\]
In other words, \(a^{-1}\cdot b^{-1}\) is the multiplicative inverse of \(ab\). As multiplicative inverses are unique by part (b), we conclude
\[
  a^{-1}\cdot b^{-1}=(ab)^{-1}.
\]
\end{solutionbox}

\newpage
\reflectionpage{Homework 3, Problem 1}

% ---------------------------------------------------------------------------
% HW3 Problem 2
% ---------------------------------------------------------------------------
\newpage
\subsection{Problem 2}

\begin{problemstatementbox}{Homework 3, Problem 2}
Given two points \((x_1,y_1)\) and \((x_2,y_2)\) in \(\R^2=\{(x,y):x,y\in\R\}\), we write
\[
  (x_1,y_1)<(x_2,y_2)
  \quad \Longleftrightarrow \quad
  x_1<x_2\text{ or }(x_1=x_2\text{ and }y_1<y_2).
\]
\begin{enumerate}[label=(\alph*)]
  \item Prove that \(<\) is an order relation on \(\R^2\).
  \item Prove that \((\R^2,<)\) does not satisfy the least upper bound property. (Hint: Consider the set of all points on the \(y\)-axis.)
\end{enumerate}
\end{problemstatementbox}

\newpage
\firstdiagnosis{Homework 3, Problem 2}

\newpage
\recipecard{Homework 3, Problem 2}

\newpage
\begin{selfexplanationbox}{Homework 3, Problem 2}
\workquestion{1. What are trichotomy and transitivity for an order relation?}{0.9in}
\workquestion{2. Why does the first coordinate decide most comparisons in lexicographic order?}{0.95in}
\workquestion{3. What set on the \(y\)-axis should be used in part (b)?}{0.85in}
\workquestion{4. Why is that set bounded above by \((1,0)\)?}{0.9in}
\workquestion{5. Why must any upper bound have positive first coordinate, and why can it be lowered?}{1.0in}
\end{selfexplanationbox}

\newpage
\begin{skeletonbox}{Homework 3, Problem 2}
\textbf{Part (a).}
\workquestion{For trichotomy, split into \(x<x'\), \(x>x'\), and \(x=x'\):}{0.95in}
\workquestion{For transitivity, split into the cases from the definition of \(<\):}{1.0in}
\workquestion{Conclude \(<\) is an order relation:}{0.75in}

\textbf{Part (b).}
\workquestion{Let \(Y=\{(0,y):y\in\R\}\). Show \(Y\) is nonempty and bounded above:}{0.9in}
\workquestion{Assume \((a,b)\) is a least upper bound. Prove \(a>0\):}{0.95in}
\workquestion{Show \((a/2,b)\) is also an upper bound and is strictly smaller:}{0.95in}
\workquestion{State the contradiction to leastness:}{0.75in}
\end{skeletonbox}

\newpage
\fullproofpage{Homework 3, Problem 2}{Use this page for part (a).}

\newpage
\fullproofpage{Homework 3, Problem 2}{Use this page for part (b).}

\newpage
\begin{solutionbox}{Official Solution (Professor)}
(a) By definition of an order relation, there are two properties to check.

Trichotomy: Fix \((x,y),(x',y')\in\R^2\). We prove that in every possible case, exactly one of the statements \(<\), \(=\), or \(>\) is true.
\begin{itemize}
  \item Case: \(x<x'\). Then \((x,y)<(x',y')\).
  \item Case: \(x>x'\). Then \((x,y)>(x',y')\).
  \item Case: \(x=x'\). If \(y<y'\), then \((x,y)<(x',y')\). If \(y>y'\), then \((x,y)>(x',y')\). If \(y=y'\), then \((x,y)=(x',y')\).
\end{itemize}

Transitivity: Suppose \((x_1,y_1)<(x_2,y_2)\) and \((x_2,y_2)<(x_3,y_3)\). By definition of \(<\), there are four possible cases.
\begin{itemize}
  \item If \(x_1<x_2\) and \(x_2<x_3\), then \(x_1<x_3\).
  \item If \(x_1<x_2\), \(x_2=x_3\), and \(y_2<y_3\), then \(x_1<x_3\).
  \item If \(x_1=x_2\), \(y_1<y_2\), and \(x_2<x_3\), then \(x_1<x_3\).
  \item If \(x_1=x_2\), \(y_1<y_2\), \(x_2=x_3\), and \(y_2<y_3\), then \(x_1=x_3\) and \(y_1<y_3\).
\end{itemize}
So, in all cases, \((x_1,y_1)<(x_3,y_3)\).

(b) Suppose toward a contradiction that \((\R^2,<)\) satisfies the least upper bound property. Set
\[
  Y=\{(0,y):y\in\R\}.
\]
Then \(Y\) is nonempty, since \((0,0)\in Y\), and \(Y\) is bounded above by \((1,0)\). Therefore there exists a least upper bound \((a,b)\in\R^2\) for \(Y\).

Claim: \(a>0\). If \(a<0\), then \((0,0)\in Y\) would satisfy \((0,0)>(a,b)\), contradicting that \((a,b)\) is an upper bound for \(Y\). If \(a=0\), then \((0,b+1)\in Y\) and \((0,b+1)>(a,b)\), again contradicting that \((a,b)\) is an upper bound. Thus \(a>0\).

As \(a>0\), \((a/2,b)\) is also an upper bound for \(Y\), since
\[
  0<a/2 \quad \Longrightarrow \quad (0,y)<(a/2,b)\quad \text{for all }y\in\R.
\]
But \(a/2<a\), so \((a/2,b)<(a,b)\). This contradicts that \((a,b)\) was the least upper bound.
\end{solutionbox}

\newpage
\reflectionpage{Homework 3, Problem 2}

% ---------------------------------------------------------------------------
% HW3 Problem 3
% ---------------------------------------------------------------------------
\newpage
\subsection{Problem 3}

\begin{problemstatementbox}{Homework 3, Problem 3}
Recall that the set of complex numbers \(\mathbb C\) consists of all numbers of the form \(a+ib\) for \(a,b\in\R\). We can add and multiply complex numbers as follows:
\[
  (a+ib)+(c+id)=(a+c)+i(b+d)
\]
and
\[
  (a+ib)\cdot(c+id)=(ac-bd)+i(ad+bc).
\]
With these operations, it turns out \((\mathbb C,+,\cdot)\) is a field (but you do not need to prove this).
\begin{enumerate}[label=(\alph*)]
  \item The definition of a field says that for any \(a+ib\neq0+i0\), there exists \(c+id\in\mathbb C\) so that \((a+ib)\cdot(c+id)=1+i0\). What are \(c,d\in\R\) in terms of \(a,b\in\R\)? Prove your answer.
  \item Prove that there is no order relation \(<\) on \(\mathbb C\) that makes \((\mathbb C,+,\cdot,<)\) an ordered field. (Hint: Can \(i=0+i1\) be positive? Negative? Zero?)
\end{enumerate}
\end{problemstatementbox}

\newpage
\firstdiagnosis{Homework 3, Problem 3}

\newpage
\recipecard{Homework 3, Problem 3}

\newpage
\begin{selfexplanationbox}{Homework 3, Problem 3}
\workquestion{1. What equations result from \((a+ib)(c+id)=1+i0\)?}{0.95in}
\workquestion{2. Why is \(a^2+b^2\neq0\) when \(a+ib\neq0+i0\)?}{0.9in}
\workquestion{3. What are the formulas for \(c\) and \(d\)?}{0.85in}
\workquestion{4. In an ordered field, why is \(x^2>0\) for every nonzero \(x\)?}{0.95in}
\workquestion{5. Why does \(i^2=-1\) contradict ordered-field behavior?}{0.95in}
\end{selfexplanationbox}

\newpage
\begin{skeletonbox}{Homework 3, Problem 3}
\textbf{Part (a).}
\workquestion{Set \(c=\frac{a}{a^2+b^2}\), \(d=-\frac{b}{a^2+b^2}\):}{0.85in}
\workquestion{Verify \(ac-bd=1\):}{0.85in}
\workquestion{Verify \(ad+bc=0\):}{0.85in}

\textbf{Part (b).}
\workquestion{Assume \(\mathbb C\) is an ordered field. Since \(i\neq0\), use the ordered-field fact that \(i^2>0\):}{0.95in}
\workquestion{Use \(i^2=-1\) and \(1>0\) to get a contradiction:}{0.95in}
\end{skeletonbox}

\newpage
\fullproofpage{Homework 3, Problem 3}{Use this page for part (a).}

\newpage
\fullproofpage{Homework 3, Problem 3}{Use this page for part (b).}

\newpage
\begin{solutionbox}{Official Solution (Professor)}
(a) Set
\[
  c=\frac{a}{a^2+b^2},
  \qquad
  d=-\frac{b}{a^2+b^2}.
\]
Then
\[
  ac-bd=\frac{a^2}{a^2+b^2}+\frac{b^2}{a^2+b^2}=1,
\]
and
\[
  ad+bc=-\frac{ab}{a^2+b^2}+\frac{ab}{a^2+b^2}=0.
\]
Therefore \((a+ib)\cdot(c+id)=1\).

(b) Suppose toward a contradiction that \((\mathbb C,+,\cdot,<)\) is an ordered field. As \(i\neq0\), a proposition from lecture gives
\[
  0<i^2=-1.
\]
On the other hand, we also proved in lecture that
\[
  1>0,
  \qquad \text{and so} \qquad
  -1<0.
\]
Altogether we have both \(-1>0\) and \(-1<0\), which is a contradiction.
\end{solutionbox}

\newpage
\reflectionpage{Homework 3, Problem 3}

% ---------------------------------------------------------------------------
% HW3 Problem 4
% ---------------------------------------------------------------------------
\newpage
\subsection{Problem 4}

\begin{problemstatementbox}{Homework 3, Problem 4}
Let \(A\subseteq\R\) be a nonempty and bounded set, let \(c\in\R\) be a number, and define
\[
  B=\{c\cdot a:a\in A\}.
\]
\begin{enumerate}[label=(\alph*)]
  \item Prove that if \(c>0\), then \(\sup B=c\cdot\sup A\).
  \item Prove that if \(c<0\), then \(\sup B=c\cdot\inf A\).
\end{enumerate}
\end{problemstatementbox}

\newpage
\firstdiagnosis{Homework 3, Problem 4}

\newpage
\recipecard{Homework 3, Problem 4}

\newpage
\begin{selfexplanationbox}{Homework 3, Problem 4}
\workquestion{1. What two facts prove that a candidate is a supremum?}{0.9in}
\workquestion{2. Why does multiplying by \(c>0\) preserve inequalities?}{0.9in}
\workquestion{3. Why does multiplying by \(c<0\) reverse inequalities?}{0.9in}
\workquestion{4. In part (a), how does an upper bound \(L\) for \(B\) produce an upper bound \(L/c\) for \(A\)?}{1.0in}
\workquestion{5. In part (b), how does an upper bound \(L\) for \(B\) produce a lower bound \(L/c\) for \(A\)?}{1.0in}
\end{selfexplanationbox}

\newpage
\begin{skeletonbox}{Homework 3, Problem 4}
\textbf{Part (a), \(c>0\).}
\workquestion{Show \(c\sup A\) is an upper bound for \(B\):}{0.9in}
\workquestion{Let \(L\) be any upper bound for \(B\). From \(ca\le L\), prove \(a\le L/c\):}{0.95in}
\workquestion{Use \(\sup A\le L/c\) to conclude \(c\sup A\le L\):}{0.85in}

\textbf{Part (b), \(c<0\).}
\workquestion{Show \(c\inf A\) is an upper bound for \(B\):}{0.9in}
\workquestion{Let \(L\) be any upper bound for \(B\). From \(ca\le L\), prove \(a\ge L/c\):}{0.95in}
\workquestion{Use \(L/c\le\inf A\), then multiply by \(c<0\):}{0.85in}
\end{skeletonbox}

\newpage
\fullproofpage{Homework 3, Problem 4}{Use this page for part (a).}

\newpage
\fullproofpage{Homework 3, Problem 4}{Use this page for part (b).}

\newpage
\begin{solutionbox}{Official Solution (Professor)}
(a) As \(A\) is nonempty and bounded, \(\sup A\in\R\) exists.

Claim: \(c\sup A\) is an upper bound for \(B\). As \(\sup A\) is an upper bound for \(A\), we know that
\[
  a\le\sup A \qquad \text{for all }a\in A,
\]
and so
\[
  ca\le c\sup A \qquad \text{for all }a\in A.
\]
Therefore \(c\sup A\) is an upper bound for \(B\).

Claim: If \(L\) is an upper bound for \(B\), then \(L\ge c\sup A\). As \(L\) is an upper bound for \(B\), we have
\[
  ca\le L \qquad \text{for all }a\in A,
\]
and so
\[
  a\le \frac Lc \qquad \text{for all }a\in A.
\]
In other words, \(\frac Lc\) is an upper bound for \(A\). As \(\sup A\) is the least upper bound, \(\sup A\le\frac Lc\), and therefore \(c\sup A\le L\). Collecting the previous two claims, \(\sup B=c\sup A\).

(b) As \(A\) is nonempty and bounded, \(\inf A\in\R\) exists.

Claim: \(c\inf A\) is an upper bound for \(B\). As \(\inf A\) is a lower bound for \(A\), we know that
\[
  a\ge\inf A \qquad \text{for all }a\in A,
\]
and so
\[
  ca\le c\inf A \qquad \text{for all }a\in A.
\]
Therefore \(c\inf A\) is an upper bound for \(B\).

Claim: If \(L\) is an upper bound for \(B\), then \(L\ge c\inf A\). As \(L\) is an upper bound for \(B\), we have
\[
  ca\le L \qquad \text{for all }a\in A,
\]
and so
\[
  a\ge \frac Lc \qquad \text{for all }a\in A.
\]
Thus \(\frac Lc\) is a lower bound for \(A\). As \(\inf A\) is the greatest lower bound, \(\inf A\ge\frac Lc\), and therefore \(c\inf A\le L\). Collecting the previous two claims, \(\sup B=c\inf A\).
\end{solutionbox}

\newpage
\reflectionpage{Homework 3, Problem 4}

% ---------------------------------------------------------------------------
% HW3 Problem 5
% ---------------------------------------------------------------------------
\newpage
\subsection{Problem 5}

\begin{problemstatementbox}{Homework 3, Problem 5}
Let \(A,B\subseteq\R\) be nonempty and bounded sets, and define
\[
  S=\{a+b:a\in A\text{ and }b\in B\}.
\]
Prove that \(\sup S=\sup A+\sup B\).
\end{problemstatementbox}

\newpage
\firstdiagnosis{Homework 3, Problem 5}

\newpage
\recipecard{Homework 3, Problem 5}

\newpage
\begin{selfexplanationbox}{Homework 3, Problem 5}
\workquestion{1. Why do \(\sup A\) and \(\sup B\) exist?}{0.85in}
\workquestion{2. Why is \(\sup A+\sup B\) an upper bound for \(S\)?}{0.95in}
\workquestion{3. How does the proof by contradiction use \(L<\sup A+\sup B\)?}{0.95in}
\workquestion{4. Why can you find \(a_0\in A\) with \(a_0>\sup A-\varepsilon\)?}{0.95in}
\workquestion{5. How do the choices of \(a_0,b_0\) contradict \(L\) being an upper bound?}{0.95in}
\end{selfexplanationbox}

\newpage
\begin{skeletonbox}{Homework 3, Problem 5}
\workquestion{Let \(\alpha=\sup A\) and \(\beta=\sup B\):}{0.75in}
\workquestion{Show \(\alpha+\beta\) is an upper bound for \(S\):}{0.9in}
\workquestion{Let \(L\) be an upper bound and suppose \(L<\alpha+\beta\). Set \(\varepsilon=\)}{0.85in}
\workquestion{Use non-leastness of \(\alpha-\varepsilon\) and \(\beta-\varepsilon\) to choose \(a_0,b_0\):}{1.0in}
\workquestion{Show \(a_0+b_0>L\), contradiction:}{0.9in}
\workquestion{Conclude \(\sup S=\alpha+\beta\):}{0.75in}
\end{skeletonbox}

\newpage
\fullproofpage{Homework 3, Problem 5}{Use this page for the complete proof.}

\newpage
\begin{solutionbox}{Official Solution (Professor)}
As \(A,B\subseteq\R\) are nonempty and bounded, \(\sup A\) and \(\sup B\) exist.

Claim: \(\sup A+\sup B\) is an upper bound for \(S\). As \(\sup A\) is an upper bound for \(A\), we know
\[
  a\le\sup A \qquad \text{for all }a\in A.
\]
Similarly, \(\sup B\) is an upper bound for \(B\), and so
\[
  b\le\sup B \qquad \text{for all }b\in B.
\]
Adding these two inequalities, we conclude that
\[
  a+b\le\sup A+\sup B
  \qquad \text{for all }a\in A,\ b\in B.
\]
Therefore \(\sup A+\sup B\) is an upper bound for \(S\).

Claim: If \(L\) is an upper bound for \(S\), then \(L\ge\sup A+\sup B\). Suppose toward a contradiction that \(L\) is an upper bound for \(S\) but \(L<\sup A+\sup B\). Set
\[
  \varepsilon=\frac12(\sup A+\sup B-L)>0.
\]
As \(\sup A-\varepsilon<\sup A\), then \(\sup A-\varepsilon\) cannot be an upper bound for \(A\), and so there exists \(a_0\in A\) such that
\[
  a_0>\sup A-\varepsilon.
\]
Similarly, there exists \(b_0\in B\) such that
\[
  b_0>\sup B-\varepsilon.
\]
Adding these inequalities and using the definition of \(\varepsilon\),
\[
  a_0+b_0>\sup A+\sup B-2\varepsilon=L.
\]
This contradicts that \(L\) is an upper bound for \(S\), so \(L\ge\sup A+\sup B\).

Therefore \(\sup S=\sup A+\sup B\).
\end{solutionbox}

\newpage
\reflectionpage{Homework 3, Problem 5}

% ---------------------------------------------------------------------------
% HW3 Problem 6
% ---------------------------------------------------------------------------
\newpage
\subsection{Problem 6}

\begin{problemstatementbox}{Homework 3, Problem 6}
Let \(A,B\subseteq\R\) be nonempty and bounded sets. Prove or disprove the following statements:
\begin{enumerate}[label=(\alph*)]
  \item If \(a\le b\) for all \(a\in A\) and for all \(b\in B\), then \(\sup A\le\inf B\).
  \item If \(a<b\) for all \(a\in A\) and for all \(b\in B\), then \(\sup A<\inf B\).
\end{enumerate}
\end{problemstatementbox}

\newpage
\firstdiagnosis{Homework 3, Problem 6}

\newpage
\recipecard{Homework 3, Problem 6}

\newpage
\begin{selfexplanationbox}{Homework 3, Problem 6}
\workquestion{1. In part (a), why is each \(b\in B\) an upper bound for \(A\)?}{0.95in}
\workquestion{2. Why does that imply \(\sup A\le b\) for all \(b\in B\)?}{0.95in}
\workquestion{3. Why is \(\sup A\) then a lower bound for \(B\)?}{0.9in}
\workquestion{4. In part (b), why might strict pointwise inequalities fail to create a strict gap between sup and inf?}{1.0in}
\workquestion{5. What counterexample has \(\sup A=\inf B\) but \(a<b\) for all choices?}{0.95in}
\end{selfexplanationbox}

\newpage
\begin{skeletonbox}{Homework 3, Problem 6}
\textbf{Part (a).}
\workquestion{Fix \(b\in B\). Use the hypothesis to show \(b\) is an upper bound for \(A\):}{0.9in}
\workquestion{Conclude \(\sup A\le b\) for all \(b\in B\):}{0.85in}
\workquestion{Use the definition of \(\inf B\) to get \(\sup A\le\inf B\):}{0.9in}

\textbf{Part (b).}
\workquestion{Give the counterexample \(A=\{1-\frac1n:n\in\N\}\), \(B=\{1\}\):}{0.95in}
\workquestion{Verify nonempty and bounded:}{0.75in}
\workquestion{Verify \(a<b\) for all \(a\in A,b\in B\), but \(\sup A=\inf B\):}{0.95in}
\end{skeletonbox}

\newpage
\fullproofpage{Homework 3, Problem 6}{Use this page for part (a).}

\newpage
\fullproofpage{Homework 3, Problem 6}{Use this page for part (b).}

\newpage
\begin{solutionbox}{Official Solution (Professor)}
(a) Prove: As \(A,B\subseteq\R\) are nonempty and bounded, both \(\sup A,\inf B\in\R\) exist.

Let \(b\in B\). As \(a\le b\) for all \(a\in A\), \(b\) is an upper bound for \(A\), and so
\[
  \sup A\le b.
\]
This is true for all \(b\in B\). Therefore \(\sup A\) is a lower bound for \(B\), and so
\[
  \sup A\le\inf B.
\]

(b) Disprove: Consider the sets
\[
  A=\left\{1-\frac1n:n\in\N\right\},
  \qquad
  B=\{1\}.
\]
Note that \(\inf B=\min B=1\), and recall from lecture that \(\sup A=1\).

Altogether:
\begin{itemize}
  \item \(A,B\) are nonempty.
  \item \(A,B\) are bounded: \(0\) is a lower bound and \(1\) is an upper bound for both \(A\) and \(B\).
  \item \(a<b\) for all \(a\in A\) and \(b\in B\), since \(1-\frac1n<1\) for all \(n\ge1\).
  \item \(\sup A\not<\inf B\), since \(\sup A=1=\inf B\).
\end{itemize}
\end{solutionbox}

\newpage
\reflectionpage{Homework 3, Problem 6}

\section{Homework 4}

% ---------------------------------------------------------------------------
% HW4 Problem 1
% ---------------------------------------------------------------------------
\newpage
\subsection{Problem 1}

\begin{problemstatementbox}{Homework 4, Problem 1}
Let \(\{a_n\}_{n\ge1}\) be a sequence.
\begin{enumerate}[label=(\alph*)]
  \item Prove that if \(\{a_n\}_{n\ge1}\) converges to \(a\), then the sequence \(\{|a_n|\}_{n\ge1}\) converges to \(|a|\). (Hint: \(\bigl||x|-|y|\bigr|\le |x-y|\) for any \(x,y\in\R\).)
  \item Provide an example where \(\{|a_n|\}_{n\ge1}\) converges, but \(\{a_n\}_{n\ge1}\) diverges.
\end{enumerate}
\end{problemstatementbox}

\newpage
\firstdiagnosis{Homework 4, Problem 1}

\newpage
\recipecard{Homework 4, Problem 1}

\newpage
\begin{selfexplanationbox}{Homework 4, Problem 1}
\workquestion{1. What does \(a_n\to a\) mean in \(\varepsilon\)-\(N\) language?}{0.9in}
\workquestion{2. How does the reverse triangle inequality compare \(\bigl||a_n|-|a|\bigr|\) with \(|a_n-a|\)?}{0.95in}
\workquestion{3. Why does the same \(N\) work for \(|a_n|\to|a|\)?}{0.9in}
\workquestion{4. What is a simple sign-oscillating sequence whose absolute values are constant?}{0.95in}
\workquestion{5. Why does that sequence fail to converge?}{0.9in}
\end{selfexplanationbox}

\newpage
\begin{skeletonbox}{Homework 4, Problem 1}
\textbf{Part (a).}
\workquestion{Fix \(\varepsilon>0\). Since \(a_n\to a\), choose \(N\) with:}{0.85in}
\workquestion{Apply \(\bigl||x|-|y|\bigr|\le |x-y|\) with \(x=a_n\), \(y=a\):}{0.9in}
\workquestion{Conclude \(|a_n|\to|a|\):}{0.75in}

\textbf{Part (b).}
\workquestion{Choose \(a_n=(-1)^n\):}{0.75in}
\workquestion{Show \(|a_n|=1\) converges, while \(a_n\) diverges:}{0.9in}
\end{skeletonbox}

\newpage
\fullproofpage{Homework 4, Problem 1}{Use this page for parts (a) and (b).}

\newpage
\begin{solutionbox}{Official Solution (Professor)}
(a) Fix \(\varepsilon>0\). As \(\lim_{n\to\infty}a_n=a\), there exists \(N\in\N\) so that
\[
  |a_n-a|<\varepsilon \qquad \text{for all }n\ge N.
\]
By the reverse triangle inequality, we have
\[
  \bigl||a_n|-|a|\bigr|\le |a_n-a|<\varepsilon
  \qquad \text{for all }n\ge N.
\]
So \(\lim_{n\to\infty}|a_n|=|a|\).

(b) By Lecture 10, the sequence \(\{a_n\}_{n\ge1}\) where \(a_n=(-1)^n\) diverges. On the other hand, \(|a_n|=1\) for all \(n\ge1\), and so \(\{|a_n|\}_{n\ge1}\) converges to \(1\).
\end{solutionbox}

\newpage
\reflectionpage{Homework 4, Problem 1}

% ---------------------------------------------------------------------------
% HW4 Problem 2
% ---------------------------------------------------------------------------
\newpage
\subsection{Problem 2}

\begin{problemstatementbox}{Homework 4, Problem 2}
(Squeeze theorem). Prove that if \(\{a_n\}_{n\ge1}\), \(\{b_n\}_{n\ge1}\), and \(\{c_n\}_{n\ge1}\) are sequences such that
\[
  a_n\le b_n\le c_n\quad\text{for all }n\ge1
\]
and
\[
  \lim_{n\to\infty}a_n=L=\lim_{n\to\infty}c_n,
\]
then \(\{b_n\}_{n\ge1}\) converges and \(\lim_{n\to\infty}b_n=L\).
\end{problemstatementbox}

\newpage
\firstdiagnosis{Homework 4, Problem 2}

\newpage
\recipecard{Homework 4, Problem 2}

\newpage
\begin{selfexplanationbox}{Homework 4, Problem 2}
\workquestion{1. What does convergence of \(a_n\) to \(L\) give for a fixed \(\varepsilon>0\)?}{0.9in}
\workquestion{2. What does convergence of \(c_n\) to \(L\) give?}{0.9in}
\workquestion{3. Why should \(N=\max\{N_1,N_2\}\) be used?}{0.85in}
\workquestion{4. How do \(a_n\le b_n\le c_n\) and the two convergence estimates trap \(b_n\)?}{1.0in}
\workquestion{5. Why is \(L-\varepsilon<b_n<L+\varepsilon\) equivalent to \(|b_n-L|<\varepsilon\)?}{0.95in}
\end{selfexplanationbox}

\newpage
\begin{skeletonbox}{Homework 4, Problem 2}
\workquestion{Fix \(\varepsilon>0\). Choose \(N_1\) so \(|a_n-L|<\varepsilon\) for \(n\ge N_1\):}{0.9in}
\workquestion{Choose \(N_2\) so \(|c_n-L|<\varepsilon\) for \(n\ge N_2\):}{0.9in}
\workquestion{Let \(N=\max\{N_1,N_2\}\), then for \(n\ge N\) write the full chain:}{1.0in}
\workquestion{Conclude \(|b_n-L|<\varepsilon\):}{0.8in}
\workquestion{State \(b_n\to L\):}{0.75in}
\end{skeletonbox}

\newpage
\fullproofpage{Homework 4, Problem 2}{Use this page for the complete proof.}

\newpage
\begin{solutionbox}{Official Solution (Professor)}
Fix \(\varepsilon>0\). As \(\lim_{n\to\infty}a_n=L\), there exists \(N_1\in\N\) so that
\[
  |a_n-L|<\varepsilon \qquad \text{for all } n\ge N_1.
\]
Similarly, as \(\lim_{n\to\infty}c_n=L\), there exists \(N_2\in\N\) so that
\[
  |c_n-L|<\varepsilon \qquad \text{for all } n\ge N_2.
\]
Set \(N=\max\{N_1,N_2\}\). Then, for \(n\ge N\),
\[
  b_n-L\le c_n-L<\varepsilon
  \qquad \text{and} \qquad
  b_n-L\ge a_n-L>-\varepsilon.
\]
Together, we conclude
\[
  |b_n-L|<\varepsilon \qquad \text{for all } n\ge N.
\]
So \(\lim_{n\to\infty}b_n=L\), as desired.
\end{solutionbox}

\newpage
\reflectionpage{Homework 4, Problem 2}

% ---------------------------------------------------------------------------
% HW4 Problem 3
% ---------------------------------------------------------------------------
\newpage
\subsection{Problem 3}

\begin{problemstatementbox}{Homework 4, Problem 3}
Suppose that the sequence \(\{a_n\}_{n\ge1}\) converges.
\begin{enumerate}[label=(\alph*)]
  \item Prove that if \(a_n\ge m\) for all but finitely many \(n\), then \(\lim_{n\to\infty}a_n\ge m\).
  \item Prove that if \(a_n\le M\) for all but finitely many \(n\), then \(\lim_{n\to\infty}a_n\le M\).
\end{enumerate}
\end{problemstatementbox}

\newpage
\firstdiagnosis{Homework 4, Problem 3}

\newpage
\recipecard{Homework 4, Problem 3}

\newpage
\begin{selfexplanationbox}{Homework 4, Problem 3}
\workquestion{1. What does "for all but finitely many \(n\)" mean in terms of an \(N_1\)?}{0.95in}
\workquestion{2. Why is it useful to evaluate the convergent sequence at \(N=\max\{N_1,N_2\}\)?}{0.95in}
\workquestion{3. In part (a), how does \(a-a_N\) get controlled by \(\varepsilon\)?}{0.95in}
\workquestion{4. Why does \(a-m>-\varepsilon\) for all \(\varepsilon>0\) imply \(a\ge m\)?}{1.0in}
\workquestion{5. How is part (b) the symmetric argument?}{0.9in}
\end{selfexplanationbox}

\newpage
\begin{skeletonbox}{Homework 4, Problem 3}
\textbf{Part (a).}
\workquestion{Let \(a=\lim a_n\). Choose \(N_1\) so \(a_n\ge m\) for all \(n\ge N_1\):}{0.9in}
\workquestion{For \(\varepsilon>0\), choose \(N_2\) so \(|a_n-a|<\varepsilon\) for all \(n\ge N_2\):}{0.9in}
\workquestion{Set \(N=\max\{N_1,N_2\}\) and show \(a-m=(a-a_N)+(a_N-m)>-\varepsilon\):}{1.0in}
\workquestion{Conclude \(a\ge m\):}{0.75in}

\textbf{Part (b).}
\workquestion{Repeat with \(a_n\le M\) to get \(a-M<\varepsilon\) for all \(\varepsilon>0\):}{0.95in}
\end{skeletonbox}

\newpage
\fullproofpage{Homework 4, Problem 3}{Use this page for both parts.}

\newpage
\begin{solutionbox}{Official Solution (Professor)}
(a) Set \(a=\lim_{n\to\infty}a_n\). Let \(N_1\in\N\) be larger than the largest index \(n\) for which \(a_n\not\ge m\). Then
\[
  a_n\ge m \qquad \text{for all } n\ge N_1.
\]
Fix \(\varepsilon>0\). As \(\lim_{n\to\infty}a_n=a\), there exists \(N_2\in\N\) so that
\[
  |a_n-a|<\varepsilon \qquad \text{for all }n\ge N_2.
\]
Set \(N=\max\{N_1,N_2\}\). Then
\[
  a-m=(a-a_N)+(a_N-m)>-\varepsilon+0.
\]
As \(\varepsilon>0\) was arbitrary, we must have \(a-m\ge0\). Otherwise, picking \(\varepsilon=m-a>0\) would yield a contradiction.

(b) Set \(a=\lim_{n\to\infty}a_n\). Let \(N_1\in\N\) be larger than the largest index \(n\) for which \(a_n\not\le M\). Then
\[
  a_n\le M \qquad \text{for all }n\ge N_1.
\]
Fix \(\varepsilon>0\). As \(\lim_{n\to\infty}a_n=a\), there exists \(N_2\in\N\) so that
\[
  |a_n-a|<\varepsilon \qquad \text{for all }n\ge N_2.
\]
Set \(N=\max\{N_1,N_2\}\). Then
\[
  a-M=(a-a_N)+(a_N-M)<\varepsilon+0.
\]
As \(\varepsilon>0\) was arbitrary, we must have \(a-M\le0\).
\end{solutionbox}

\newpage
\reflectionpage{Homework 4, Problem 3}

% ---------------------------------------------------------------------------
% HW4 Problem 4
% ---------------------------------------------------------------------------
\newpage
\subsection{Problem 4}

\begin{problemstatementbox}{Homework 4, Problem 4}
Suppose that the sequence \(\{a_n\}_{n\ge1}\) converges. Prove or disprove the following statements:
\begin{enumerate}[label=(\alph*)]
  \item If \(\lim_{n\to\infty}a_n\ge m\), then there exists \(N\in\N\) so that \(a_n\ge m\) for all \(n\ge N\).
  \item If \(\lim_{n\to\infty}a_n>m\), then there exists \(N\in\N\) so that \(a_n>m\) for all \(n\ge N\).
\end{enumerate}
\end{problemstatementbox}

\newpage
\firstdiagnosis{Homework 4, Problem 4}

\newpage
\recipecard{Homework 4, Problem 4}

\newpage
\begin{selfexplanationbox}{Homework 4, Problem 4}
\workquestion{1. Why is the weak inequality in part (a) suspicious?}{0.95in}
\workquestion{2. What counterexample converges to \(0\) from below?}{0.9in}
\workquestion{3. Why does that counterexample never satisfy \(a_n\ge0\)?}{0.85in}
\workquestion{4. In part (b), why does \(\lim a_n>m\) give a positive gap?}{0.9in}
\workquestion{5. What choice of \(\varepsilon\) keeps \(a_n\) eventually above \(m\)?}{0.95in}
\end{selfexplanationbox}

\newpage
\begin{skeletonbox}{Homework 4, Problem 4}
\textbf{Part (a).}
\workquestion{Disprove using \(a_n=-1/n\) and \(m=0\):}{0.9in}
\workquestion{Verify \(\lim a_n=0\ge0\) but \(a_n<0\) for every \(n\):}{0.9in}

\textbf{Part (b).}
\workquestion{Let \(a=\lim a_n\), assume \(a>m\), and set \(\varepsilon=a-m\):}{0.85in}
\workquestion{Use convergence to get \(|a_n-a|<\varepsilon\) eventually:}{0.85in}
\workquestion{Show \(a_n-m=(a_n-a)+(a-m)>0\):}{0.9in}
\end{skeletonbox}

\newpage
\fullproofpage{Homework 4, Problem 4}{Use this page for both parts.}

\newpage
\begin{solutionbox}{Official Solution (Professor)}
(a) Disprove. Consider the sequence \(\{a_n\}_{n\ge1}\) with
\[
  a_n=-\frac1n.
\]
From Lecture 10, \(\lim_{n\to\infty}\frac1n=0\), and so \(a=\lim_{n\to\infty}a_n=0\) by the limit law for multiplication by a constant. So \(m=0\) satisfies
\[
  a=0\ge0=m.
\]
On the other hand,
\[
  a_n=-\frac1n\not\ge0=m
\]
for all \(n\in\N\), and so no such \(N\) can exist.

(b) Prove. Suppose \(a>m\) and set
\[
  \varepsilon=a-m.
\]
Then \(\varepsilon>0\), and so there exists \(N\in\N\) such that
\[
  |a_n-a|<\varepsilon \qquad \text{for all } n\ge N.
\]
Then, for \(n\ge N\),
\[
  a_n-m=(a_n-a)+(a-m)>-\varepsilon+(a-m)=0.
\]
Therefore \(a_n>m\) for all \(n\ge N\), as desired.
\end{solutionbox}

\newpage
\reflectionpage{Homework 4, Problem 4}

% ---------------------------------------------------------------------------
% HW4 Problem 5
% ---------------------------------------------------------------------------
\newpage
\subsection{Problem 5}

\begin{problemstatementbox}{Homework 4, Problem 5}
Prove that if \(\{a_n\}_{n\ge1}\) is a sequence of positive numbers that converges to \(a>0\), then the sequence \(\{\sqrt{a_n}\}_{n\ge1}\) converges to \(\sqrt a\).
\end{problemstatementbox}

\newpage
\firstdiagnosis{Homework 4, Problem 5}

\newpage
\recipecard{Homework 4, Problem 5}

\newpage
\begin{selfexplanationbox}{Homework 4, Problem 5}
\workquestion{1. What algebraic identity rationalizes \(\sqrt{a_n}-\sqrt a\)?}{0.95in}
\workquestion{2. Why is the denominator \(\sqrt{a_n}+\sqrt a\) bounded below by \(\sqrt a\)?}{0.95in}
\workquestion{3. What tolerance for \(|a_n-a|\) should be used to make \(|\sqrt{a_n}-\sqrt a|<\varepsilon\)?}{1.0in}
\workquestion{4. Where is the assumption \(a>0\) used?}{0.9in}
\workquestion{5. Where is the assumption \(a_n>0\) used?}{0.9in}
\end{selfexplanationbox}

\newpage
\begin{skeletonbox}{Homework 4, Problem 5}
\workquestion{Fix \(\varepsilon>0\). Since \(a_n\to a\), choose \(N\) so \(|a_n-a|<\sqrt a\,\varepsilon\):}{0.95in}
\workquestion{For \(n\ge N\), write \(|\sqrt{a_n}-\sqrt a|=\frac{|a_n-a|}{\sqrt{a_n}+\sqrt a}\):}{1.0in}
\workquestion{Use \(\sqrt{a_n}+\sqrt a\ge\sqrt a\):}{0.85in}
\workquestion{Conclude \(|\sqrt{a_n}-\sqrt a|<\varepsilon\):}{0.8in}
\end{skeletonbox}

\newpage
\fullproofpage{Homework 4, Problem 5}{Use this page for the complete proof.}

\newpage
\begin{solutionbox}{Official Solution (Professor)}
Fix \(\varepsilon>0\). As \(\lim_{n\to\infty}a_n=a\) and \(a>0\), there exists \(N\in\N\) so that
\[
  |a_n-a|<\sqrt a\,\varepsilon
  \qquad \text{for all } n\ge N.
\]
Then, for \(n\ge N\),
\[
  |\sqrt{a_n}-\sqrt a|
  =
  \frac{|(\sqrt{a_n}-\sqrt a)(\sqrt{a_n}+\sqrt a)|}{\sqrt{a_n}+\sqrt a}
  =
  \frac{|a_n-a|}{\sqrt{a_n}+\sqrt a}
  \le
  \frac{|a_n-a|}{\sqrt a}
  <
  \frac{\sqrt a\,\varepsilon}{\sqrt a}
  =
  \varepsilon.
\]
In the inequality above, we used \(\sqrt{a_n}+\sqrt a\ge\sqrt a\) since \(a_n>0\). Therefore \(\lim_{n\to\infty}\sqrt{a_n}=\sqrt a\), as desired.
\end{solutionbox}

\newpage
\reflectionpage{Homework 4, Problem 5}

% ---------------------------------------------------------------------------
% HW4 Problem 6
% ---------------------------------------------------------------------------
\newpage
\subsection{Problem 6}

\begin{problemstatementbox}{Homework 4, Problem 6}
Consider the following sequence:
\[
  a_1=\sqrt2
  \qquad\text{and}\qquad
  a_{n+1}=\sqrt{2+a_n}
  \quad\text{for all }n\ge1.
\]
\begin{enumerate}[label=(\alph*)]
  \item Prove that the sequence \(\{a_n\}_{n\ge1}\) is bounded.
  \item Prove that the sequence is increasing.
  \item Conclude that the sequence converges, and find its limit.
\end{enumerate}
\end{problemstatementbox}

\newpage
\firstdiagnosis{Homework 4, Problem 6}

\newpage
\recipecard{Homework 4, Problem 6}

\newpage
\begin{selfexplanationbox}{Homework 4, Problem 6}
\workquestion{1. What upper and lower bounds should be proved by induction?}{0.95in}
\workquestion{2. Why does \(a_n<2\) imply \(a_{n+1}<2\)?}{0.9in}
\workquestion{3. Why is comparing \(a_n\le a_{n+1}\) equivalent to comparing squares?}{0.95in}
\workquestion{4. How does \((a_n-2)(a_n+1)<0\) prove monotonicity?}{0.95in}
\workquestion{5. After convergence, what equation does the limit satisfy?}{0.95in}
\end{selfexplanationbox}

\newpage
\begin{skeletonbox}{Homework 4, Problem 6}
\textbf{Part (a).}
\workquestion{Prove \(\sqrt2\le a_n<2\) by induction: base case:}{0.85in}
\workquestion{Inductive step: from \(\sqrt2\le a_n<2\), show \(\sqrt2\le a_{n+1}<2\):}{0.95in}

\textbf{Part (b).}
\workquestion{Use \(a_n>0\) and compare \(a_n\le\sqrt{2+a_n}\) by squaring:}{0.95in}
\workquestion{Factor \(a_n^2-a_n-2=(a_n-2)(a_n+1)\):}{0.85in}

\textbf{Part (c).}
\workquestion{Invoke monotone convergence. Let \(a=\lim a_n\):}{0.85in}
\workquestion{Use Problem 5 and limit laws to pass to the limit in \(a_{n+1}=\sqrt{2+a_n}\):}{0.95in}
\workquestion{Solve \(a=\sqrt{2+a}\) and choose the valid root:}{0.9in}
\end{skeletonbox}

\newpage
\fullproofpage{Homework 4, Problem 6}{Use this page for parts (a) and (b).}

\newpage
\fullproofpage{Homework 4, Problem 6}{Use this page for part (c).}

\newpage
\begin{solutionbox}{Official Solution (Professor)}
(a) We will prove \(\sqrt2\le a_n<2\) for all \(n\in\N\) by induction.

Base case: For \(n=1\), the statement \(\sqrt2\le\sqrt2<2\) is true.

Inductive step: Suppose \(\sqrt2\le a_n<2\) for some \(n\in\N\). Then
\[
  a_{n+1}=\sqrt{2+a_n}>\sqrt{2+0}=\sqrt2,
\]
and
\[
  a_{n+1}=\sqrt{2+a_n}<\sqrt{2+2}=2.
\]
Together, \(\sqrt2\le a_{n+1}<2\). By induction, \(\sqrt2\le a_n<2\) for all \(n\in\N\).

(b) Fix \(n\in\N\). Then
\[
  a_n\le a_{n+1}=\sqrt{2+a_n}
  \Longleftrightarrow
  a_n^2\le 2+a_n
  \Longleftrightarrow
  0\ge a_n^2-a_n-2=(a_n-2)(a_n+1),
\]
where the first equivalence uses \(a_n>0\) from part (a). By part (a), \(a_n-2<0\) and \(a_n+1>0\), so \((a_n-2)(a_n+1)<0\). Working backwards, \(a_n\le a_{n+1}\).

(c) As \(\{a_n\}_{n\ge1}\) is increasing and bounded, it converges by Lecture 11. By part (a) and Problem 3,
\[
  a=\lim_{n\to\infty}a_n
  \quad\text{satisfies}\quad
  \sqrt2\le a\le2.
\]
By the limit law for addition and Problem 5,
\[
  \lim_{n\to\infty}(2+a_n)=2+a>0,
  \qquad
  \lim_{n\to\infty}\sqrt{2+a_n}=\sqrt{2+a}.
\]
Taking the limit \(n\to\infty\) in
\[
  a_{n+1}=\sqrt{2+a_n},
\]
we see that \(a\) satisfies
\[
  a=\sqrt{2+a}.
\]
Thus
\[
  a^2=2+a
  \Longrightarrow
  0=a^2-a-2=(a-2)(a+1),
\]
so \(a=2\) or \(a=-1\). The case \(a=-1\) is not possible, since \(a\ge\sqrt2\). Therefore \(a=2\).
\end{solutionbox}

\newpage
\reflectionpage{Homework 4, Problem 6}

\section{Homework 5}

% ---------------------------------------------------------------------------
% HW5 Problem 1
% ---------------------------------------------------------------------------
\newpage
\subsection{Problem 1}

\begin{problemstatementbox}{Homework 5, Problem 1}
Let \(\{a_n\}_{n\ge1}\) be a sequence of positive real numbers. Prove that \(\lim_{n\to\infty} a_n=+\infty\) if and only if \(\lim_{n\to\infty}\frac{1}{a_n}=0\).
\end{problemstatementbox}

\newpage
\firstdiagnosis{Homework 5, Problem 1}

\newpage
\recipecard{Homework 5, Problem 1}

\newpage
\begin{selfexplanationbox}{Homework 5, Problem 1}
\workquestion{1. What is the definition of \(a_n\to+\infty\)?}{1.0in}
\workquestion{2. Why is positivity of \(a_n\) essential before taking reciprocals?}{0.95in}
\workquestion{3. In the forward direction, what should \(M\) be in terms of \(\varepsilon\)?}{0.95in}
\workquestion{4. In the reverse direction, what should \(\varepsilon\) be in terms of \(M\)?}{0.95in}
\workquestion{5. Where does the inequality reversal happen?}{0.95in}
\end{selfexplanationbox}

\newpage
\begin{skeletonbox}{Homework 5, Problem 1}
\textbf{Forward direction.}
\workquestion{Fix \(\varepsilon>0\). Use \(a_n\to+\infty\) with \(M=\)}{0.85in}
\workquestion{For all \(n\ge N\), take reciprocals to obtain:}{0.85in}
\workquestion{Conclude \(\left|\frac1{a_n}-0\right|<\varepsilon\):}{0.8in}

\textbf{Reverse direction.}
\workquestion{Fix \(M>0\). Use \(\frac1{a_n}\to0\) with \(\varepsilon=\)}{0.85in}
\workquestion{Use positivity to convert \(\frac1{a_n}<\frac1M\) into:}{0.85in}
\workquestion{State the definition this verifies:}{0.75in}
\end{skeletonbox}

\newpage
\fullproofpage{Homework 5, Problem 1}{Use this page for both directions.}

\newpage
\begin{solutionbox}{Official Solution (Professor)}
\(\Rightarrow\) Fix \(\varepsilon>0\). As \(\lim a_n=+\infty\), then there exists \(N\) so that
\[
  a_n>\frac1\varepsilon \qquad \text{for all } n\ge N.
\]
So, for \(n\ge N\) we have
\[
  -\varepsilon<0<\frac1{a_n}<\varepsilon
  \qquad \Longrightarrow \qquad
  \left|\frac1{a_n}-0\right|<\varepsilon.
\]
Therefore \(\lim \frac1{a_n}=0\).

\(\Leftarrow\) Fix \(M>0\). As \(\lim \frac1{a_n}=0\), then there exists \(N\) so that
\[
  \left|\frac1{a_n}-0\right|<\frac1M
  \qquad \text{for all } n\ge N.
\]
Note that \(\left|\frac1{a_n}-0\right|=\frac1{a_n}>0\). So
\[
  a_n>M \qquad \text{for all } n\ge N.
\]
Therefore \(\lim a_n=+\infty\).
\end{solutionbox}

\newpage
\reflectionpage{Homework 5, Problem 1}

% ---------------------------------------------------------------------------
% HW5 Problem 2
% ---------------------------------------------------------------------------
\newpage
\subsection{Problem 2}

\begin{problemstatementbox}{Homework 5, Problem 2}
Let \(\{a_n\}_{n\ge1}\) be a sequence. Prove that if \(\{a_{k_n}\}_{n\ge1}\) is a subsequence of \(\{a_n\}_{n\ge1}\) that converges to \(L\in\R\), then
\[
  \liminf_{n\to\infty} a_n \le L \le \limsup_{n\to\infty} a_n.
\]

(Hint: As \(\{a_n\}_{n\ge1}\) is not necessarily bounded, we have to consider the cases \(\liminf_{n\to\infty}a_n=\pm\infty\) and \(\limsup_{n\to\infty}a_n=\pm\infty\). However, the proof is much easier in these cases.)
\end{problemstatementbox}

\newpage
\firstdiagnosis{Homework 5, Problem 2}

\newpage
\recipecard{Homework 5, Problem 2}

\newpage
\begin{selfexplanationbox}{Homework 5, Problem 2}
\workquestion{1. What do \(\limsup a_n\) and \(\liminf a_n\) measure using tail suprema/infima?}{1.0in}
\workquestion{2. Why does \(k_n\ge n\) imply that a fixed tail eventually contains the subsequence terms?}{1.0in}
\workquestion{3. Why is the case \(\limsup a_n=+\infty\) immediate?}{0.9in}
\workquestion{4. Why is the case \(\limsup a_n=-\infty\) impossible when \(L\in\R\)?}{1.05in}
\workquestion{5. How do you get \(a_{k_n}\le \sup\{a_m:m\ge N\}\) for all large \(n\)?}{0.95in}
\end{selfexplanationbox}

\newpage
\begin{skeletonbox}{Homework 5, Problem 2}
\textbf{Prove \(L\le\limsup a_n\).}
\workquestion{Separate the easy cases \(\limsup a_n=+\infty\) and \(\limsup a_n=-\infty\):}{1.0in}
\workquestion{Assume \(\limsup a_n\in\R\). For fixed \(N\), compare \(a_{k_n}\) to the tail supremum for all large \(n\):}{1.0in}
\workquestion{Pass \(n\to\infty\) to get \(L\le\sup\{a_m:m\ge N\}\):}{0.85in}
\workquestion{Pass \(N\to\infty\) to obtain \(L\le\limsup a_n\):}{0.85in}

\textbf{Prove \(\liminf a_n\le L\).}
\workquestion{Repeat the same argument using tail infima and reverse inequalities:}{1.0in}
\end{skeletonbox}

\newpage
\fullproofpage{Homework 5, Problem 2}{Use this page for the \(\limsup\) inequality.}

\newpage
\fullproofpage{Homework 5, Problem 2}{Use this page for the \(\liminf\) inequality and endpoint cases.}

\newpage
\begin{solutionbox}{Official Solution (Professor)}
Let's start with the second inequality:
\[
  L\le \limsup_{n\to\infty} a_n.
\]
Case: \(\limsup a_n=+\infty\). Then the statement is automatically true, since every real number \(L\) satisfies \(L<+\infty\).

Case: \(\limsup a_n=-\infty\). As \(\liminf a_n\le \limsup a_n\), then we must have \(\liminf a_n=-\infty=\limsup a_n\). This means that \(\lim a_n=-\infty\) by Lecture 14.

Fix \(m<0\). As \(\lim a_n=-\infty\), there exists \(N\) so that
\[
  a_n<m \qquad \text{for all } n\ge N.
\]
As \(k_n\ge n\) for all \(n\), this implies
\[
  a_{k_n}<m \qquad \text{for all } n\ge N.
\]
Therefore \(L=\lim a_{k_n}=-\infty\). This contradicts that \(L\in\R\), and so this case is impossible.

Case: \(\limsup a_n\in\R\). For \(N\ge1\), notice that
\[
  a_{k_n}\le \sup\{a_n:n\ge N\}
  \qquad \text{for all } n\ge N.
\]
We know that the sequence \(\{a_{k_n}\}\) converges to \(L\). So, by HW4\#3, we have
\[
  L=\lim a_{k_n}\le \sup\{a_n:n\ge N\}.
\]
On the other hand, \(v_N=\sup\{a_n:n\ge N\}\) is a sequence that converges to \(\limsup a_n\). So, using HW4\#3 again, we conclude that
\[
  L\le \lim_{N\to\infty}\sup\{a_n:n\ge N\}=\limsup a_n.
\]
Exercise: Using similar steps, prove that
\[
  \liminf_{n\to\infty} a_n\le L.
\]
\end{solutionbox}

\newpage
\reflectionpage{Homework 5, Problem 2}

% ---------------------------------------------------------------------------
% HW5 Problem 3
% ---------------------------------------------------------------------------
\newpage
\subsection{Problem 3}

\begin{problemstatementbox}{Homework 5, Problem 3}
Let \(\{a_n\}_{n\ge1}\) and \(\{b_n\}_{n\ge1}\) be two bounded sequences. Prove that
\[
  \limsup_{n\to\infty}(a_n+b_n)
  \le
  \limsup_{n\to\infty}a_n+\limsup_{n\to\infty}b_n.
\]
\end{problemstatementbox}

\newpage
\firstdiagnosis{Homework 5, Problem 3}

\newpage
\recipecard{Homework 5, Problem 3}

\newpage
\begin{selfexplanationbox}{Homework 5, Problem 3}
\workquestion{1. Why does boundedness guarantee finite tail suprema for \(a_n\), \(b_n\), and \(a_n+b_n\)?}{1.0in}
\workquestion{2. For a fixed \(N\), what upper bounds do the tail suprema give for \(a_n\) and \(b_n\)?}{1.0in}
\workquestion{3. Why does adding two pointwise upper bounds produce an upper bound for \(a_n+b_n\)?}{0.95in}
\workquestion{4. How does an upper bound for every tail term become an inequality of suprema?}{1.0in}
\workquestion{5. Which limit law finishes the proof after passing \(N\to\infty\)?}{0.9in}
\end{selfexplanationbox}

\newpage
\begin{skeletonbox}{Homework 5, Problem 3}
\workquestion{Fix \(N\ge1\). Write the two tail upper bounds for all \(n\ge N\):}{1.0in}
\workquestion{Add them to get an upper bound for \(a_n+b_n\):}{0.9in}
\workquestion{Take the supremum over \(n\ge N\) to obtain:}{0.9in}
\workquestion{Let \(N\to\infty\) and identify each term with a limsup:}{1.0in}
\workquestion{State the final inequality:}{0.75in}
\end{skeletonbox}

\newpage
\fullproofpage{Homework 5, Problem 3}{Use this page for the complete proof.}

\newpage
\begin{solutionbox}{Official Solution (Professor)}
Fix \(N\ge1\). For \(n\ge N\) we have
\[
  a_n\le \sup\{a_n:n\ge N\}, \qquad
  b_n\le \sup\{b_n:n\ge N\},
\]
and so
\[
  a_n+b_n\le \sup\{a_n:n\ge N\}+\sup\{b_n:n\ge N\}.
\]
In other words, \(\sup\{a_n:n\ge N\}+\sup\{b_n:n\ge N\}\) is an upper bound for the set \(\{a_n+b_n:n\ge N\}\). Therefore
\[
  \sup\{a_n+b_n:n\ge N\}
  \le
  \sup\{a_n:n\ge N\}+\sup\{b_n:n\ge N\}.
\]

Altogether, we have shown that for any \(N\ge1\) we have
\[
  \sup\{a_n+b_n:n\ge N\}
  \le
  \sup\{a_n:n\ge N\}+\sup\{b_n:n\ge N\}.
\]
As \(\{a_n\}\) is a bounded sequence, we know from lecture that the sequence \(v_N=\sup\{a_n:n\ge N\}\) converges to \(v=\limsup a_n\in\R\). Similarly, as \(\{b_n\}\) is bounded, the sequence \(\sup\{b_n:n\ge N\}\) converges to \(\limsup b_n\) as \(N\to\infty\). By the inequality above, we see that the sequence \(\sup\{a_n+b_n:n\ge N\}\) is also bounded above. Therefore \(\limsup(a_n+b_n)\) is a finite number, and \(\sup\{a_n+b_n:n\ge N\}\) converges to \(\limsup(a_n+b_n)\) as \(N\to\infty\).

By the limit laws, the sequence
\[
  \sup\{a_n+b_n:n\ge N\}
  -\sup\{a_n:n\ge N\}
  -\sup\{b_n:n\ge N\}
\]
converges as \(N\to\infty\). We also know that it is \(\le0\) for all \(N\ge1\). Therefore, by HW4\#3, the limit is also \(\le0\):
\[
  \limsup(a_n+b_n)-\limsup a_n-\limsup b_n\le0
\]
and hence
\[
  \limsup_{n\to\infty}(a_n+b_n)
  \le
  \limsup_{n\to\infty}a_n+\limsup_{n\to\infty}b_n.
\]
\end{solutionbox}

\newpage
\reflectionpage{Homework 5, Problem 3}

% ---------------------------------------------------------------------------
% HW5 Problem 4
% ---------------------------------------------------------------------------
\newpage
\subsection{Problem 4}

\begin{problemstatementbox}{Homework 5, Problem 4}
Let \(\{a_n\}_{n\ge1}\) be a sequence. Prove that \(\limsup_{n\to\infty}|a_n|=0\) if and only if \(\lim_{n\to\infty}a_n=0\).
\end{problemstatementbox}

\newpage
\firstdiagnosis{Homework 5, Problem 4}

\newpage
\recipecard{Homework 5, Problem 4}

\newpage
\begin{selfexplanationbox}{Homework 5, Problem 4}
\workquestion{1. Why is \(\liminf |a_n|\ge0\) automatic?}{0.9in}
\workquestion{2. How does \(\liminf |a_n|=\limsup |a_n|=0\) imply \(|a_n|\to0\)?}{1.0in}
\workquestion{3. Why is \(|a_n|\to0\) exactly the same as \(a_n\to0\)?}{0.95in}
\workquestion{4. In the reverse direction, which earlier homework result gives \(|a_n|\to0\)?}{0.95in}
\workquestion{5. How does Lecture 14 connect ordinary limits with liminf and limsup?}{0.95in}
\end{selfexplanationbox}

\newpage
\begin{skeletonbox}{Homework 5, Problem 4}
\textbf{Forward direction.}
\workquestion{Assume \(\limsup |a_n|=0\). Prove \(\liminf |a_n|\ge0\):}{0.95in}
\workquestion{Use \(\liminf |a_n|\le\limsup |a_n|\) to conclude:}{0.85in}
\workquestion{Invoke the liminf/limsup convergence criterion to get \(|a_n|\to0\):}{0.85in}
\workquestion{Translate \(|a_n|\to0\) into \(a_n\to0\):}{0.85in}

\textbf{Reverse direction.}
\workquestion{Assume \(a_n\to0\). Use HW4\#1 to get:}{0.8in}
\workquestion{Use Lecture 14 to identify \(\limsup |a_n|\):}{0.8in}
\end{skeletonbox}

\newpage
\fullproofpage{Homework 5, Problem 4}{Use this page for both directions.}

\newpage
\begin{solutionbox}{Official Solution (Professor)}
\(\Rightarrow\) Suppose \(\limsup |a_n|=0\).

Claim: \(\liminf |a_n|\ge0\). As \(|a_n|\ge0\) for all \(n\), we have
\[
  \inf_{n\ge N}|a_n|\ge0
  \qquad \text{for all } N\ge1.
\]
Therefore
\[
  \lim_{N\to\infty}\inf_{n\ge N}|a_n|\ge0
\]
by HW4\#3.

Altogether, we have
\[
  0\le \liminf |a_n|\le \limsup |a_n|=0.
\]
Therefore
\[
  \liminf |a_n|=0=\limsup |a_n|.
\]
By Lecture 10, this means that
\[
  \lim |a_n|=0.
\]
In other words: for all \(\varepsilon>0\) there exists \(N\) so that
\[
  ||a_n|-0|<\varepsilon \qquad \text{for all } n\ge N.
\]
This is the same as: for all \(\varepsilon>0\) there exists \(N\) so that
\[
  |a_n-0|<\varepsilon \qquad \text{for all } n\ge N.
\]
So \(\lim a_n=0\).

\(\Leftarrow\) Suppose \(\lim a_n=0\). By HW4\#1, this implies
\[
  \lim |a_n|=0.
\]
By Lecture 14, we have
\[
  \liminf |a_n|=0=\limsup |a_n|.
\]
So \(\limsup |a_n|=0\), as desired.
\end{solutionbox}

\newpage
\reflectionpage{Homework 5, Problem 4}

\section{Homework 6}

% ---------------------------------------------------------------------------
% HW6 Problem 1
% ---------------------------------------------------------------------------
\newpage
\subsection{Problem 1}

\begin{problemstatementbox}{Homework 6, Problem 1}
Prove that the series
\[
  \sum_{n=1}^{\infty}\frac{1}{(3n-2)(3n+1)}
\]
converges, and find its value. (Hint: You computed the partial sums on Homework 1.)
\end{problemstatementbox}

\newpage
\firstdiagnosis{Homework 6, Problem 1}

\newpage
\recipecard{Homework 6, Problem 1}

\newpage
\begin{selfexplanationbox}{Homework 6, Problem 1}
\workquestion{1. What does it mean for a series to converge in terms of partial sums?}{0.95in}
\workquestion{2. What partial fraction decomposition should be used?}{0.95in}
\workquestion{3. Which terms cancel in the \(N\)th partial sum?}{1.0in}
\workquestion{4. What explicit formula for \(S_N\) should remain after cancellation?}{0.95in}
\workquestion{5. What limit gives the value of the series?}{0.85in}
\end{selfexplanationbox}

\newpage
\begin{skeletonbox}{Homework 6, Problem 1}
\workquestion{Define the \(N\)th partial sum \(S_N=\)}{0.8in}
\workquestion{Rewrite the summand using partial fractions:}{0.9in}
\workquestion{Write the telescoping sum and identify the surviving first and last terms:}{1.05in}
\workquestion{Compute \(\lim_{N\to\infty}S_N\):}{0.9in}
\workquestion{State convergence and the value of the series:}{0.75in}
\end{skeletonbox}

\newpage
\fullproofpage{Homework 6, Problem 1}{Use this page for the complete proof.}

\newpage
\begin{solutionbox}{Student-Derived Solution (No official solution available)}
Let
\[
  a_n=\frac{1}{(3n-2)(3n+1)}
\]
and let \(S_N=\sum_{n=1}^N a_n\). Using partial fractions,
\[
  \frac{1}{(3n-2)(3n+1)}
  =
  \frac13\left(\frac1{3n-2}-\frac1{3n+1}\right).
\]
Hence
\[
  S_N
  =
  \frac13\sum_{n=1}^N
  \left(\frac1{3n-2}-\frac1{3n+1}\right).
\]
Writing out the terms gives
\[
  S_N
  =
  \frac13\left[
    \left(1-\frac14\right)
    +
    \left(\frac14-\frac17\right)
    +\cdots+
    \left(\frac1{3N-2}-\frac1{3N+1}\right)
  \right].
\]
All intermediate terms cancel, so
\[
  S_N=\frac13\left(1-\frac1{3N+1}\right)=\frac{N}{3N+1}.
\]
Therefore
\[
  \lim_{N\to\infty}S_N
  =
  \lim_{N\to\infty}\frac{N}{3N+1}
  =
  \frac13.
\]
Since the sequence of partial sums converges, the series converges, and its value is
\[
  \sum_{n=1}^{\infty}\frac{1}{(3n-2)(3n+1)}=\frac13.
\]
\end{solutionbox}

\newpage
\reflectionpage{Homework 6, Problem 1}

% ---------------------------------------------------------------------------
% HW6 Problem 2
% ---------------------------------------------------------------------------
\newpage
\subsection{Problem 2}

\begin{problemstatementbox}{Homework 6, Problem 2}
Let \(\sum_{n=1}^{\infty} a_n\) and \(\sum_{n=1}^{\infty} b_n\) be convergent series.

\begin{enumerate}[label=(\alph*)]
  \item Prove that for any \(k\in\R\), the series \(\sum_{n=1}^{\infty}ka_n\) converges and
  \[
    \sum_{n=1}^{\infty}ka_n
    =
    k\sum_{n=1}^{\infty}a_n.
  \]
  \item Prove that the series \(\sum_{n=1}^{\infty}(a_n+b_n)\) converges and
  \[
    \sum_{n=1}^{\infty}(a_n+b_n)
    =
    \sum_{n=1}^{\infty}a_n+\sum_{n=1}^{\infty}b_n.
  \]
\end{enumerate}
\end{problemstatementbox}

\newpage
\firstdiagnosis{Homework 6, Problem 2}

\newpage
\recipecard{Homework 6, Problem 2}

\newpage
\begin{selfexplanationbox}{Homework 6, Problem 2}
\workquestion{1. How is the sum of a convergent series defined using partial sums?}{0.95in}
\workquestion{2. In part (a), what are the partial sums of \(\sum ka_n\) in terms of the partial sums of \(\sum a_n\)?}{1.0in}
\workquestion{3. Which sequence limit law proves convergence in part (a)?}{0.9in}
\workquestion{4. In part (b), how do the partial sums of \(\sum(a_n+b_n)\) split?}{1.0in}
\workquestion{5. Which limit law gives the final identity for the sum?}{0.9in}
\end{selfexplanationbox}

\newpage
\begin{skeletonbox}{Homework 6, Problem 2}
\textbf{Part (a).}
\workquestion{Let \(s_n=\sum_{j=1}^n a_j\). The partial sums of \(\sum ka_n\) are:}{0.95in}
\workquestion{Use \(s_n\to \sum a_n\) and the constant multiple limit law:}{0.95in}
\workquestion{Translate the limit of partial sums into the series identity:}{0.85in}

\textbf{Part (b).}
\workquestion{Let \(s_n=\sum_{j=1}^n a_j\) and \(t_n=\sum_{j=1}^n b_j\). The partial sums are:}{0.95in}
\workquestion{Use the addition limit law and state the resulting sum:}{0.95in}
\end{skeletonbox}

\newpage
\fullproofpage{Homework 6, Problem 2}{Use this page for part (a).}

\newpage
\fullproofpage{Homework 6, Problem 2}{Use this page for part (b).}

\newpage
\begin{solutionbox}{Official Solution (Professor)}
(a) Let \(s_n=\sum_{j=1}^n a_j\) denote the \(n\)th partial sum. As \(\sum a_n\) converges, we know that the sequence \(\{s_n\}\) converges.

Fix \(k\in\R\). Then the partial sums \(\sum_{j=1}^n ka_j\) are equal to \(ks_n\). By the limit law for multiplication by a constant, we conclude that the sequence \(\{ks_n\}_{n\ge1}\) converges and \(\lim(ks_n)=k\lim s_n\). By definition, this is exactly what it means for the series \(\sum ka_n\) to converge to the number \(k\sum a_n\).

(b) Let \(s_n=\sum_{j=1}^n a_j\) and \(t_n=\sum_{j=1}^n b_j\) denote the \(n\)th partial sums. As \(\sum a_n\) and \(\sum b_n\) converge, we know that the sequences \(\{s_n\}\) and \(\{t_n\}\) converge.

Note that the partial sums \(\sum_{j=1}^n(a_j+b_j)\) are equal to \(s_n+t_n\). By the limit law for addition, we conclude that the sequence \(\{s_n+t_n\}_{n\ge1}\) converges and \(\lim(s_n+t_n)=\lim s_n+\lim t_n\). By definition, this is exactly what it means for the series \(\sum(a_n+b_n)\) to converge to the number \(\sum a_n+\sum b_n\).
\end{solutionbox}

\newpage
\reflectionpage{Homework 6, Problem 2}

% ---------------------------------------------------------------------------
% HW6 Problem 3
% ---------------------------------------------------------------------------
\newpage
\subsection{Problem 3}

\begin{problemstatementbox}{Homework 6, Problem 3}
\begin{enumerate}[label=(\alph*)]
  \item Prove that
  \[
    \sum_{n=1}^{\infty}\frac{n-1}{2^{n+1}}=\frac12.
  \]
  (Hint: \(\frac{k-1}{2^{k+1}}=\frac{k}{2^k}-\frac{k+1}{2^{k+1}}\).)
  \item Use part (a) to compute
  \[
    \sum_{n=1}^{\infty}\frac{n}{2^n}.
  \]
\end{enumerate}
\end{problemstatementbox}

\newpage
\firstdiagnosis{Homework 6, Problem 3}

\newpage
\recipecard{Homework 6, Problem 3}

\newpage
\begin{selfexplanationbox}{Homework 6, Problem 3}
\workquestion{1. What identity turns the summand in part (a) into a telescoping difference?}{0.95in}
\workquestion{2. What formula should the \(n\)th partial sum satisfy?}{0.95in}
\workquestion{3. Why does \(\frac{n}{2^n}\to0\)?}{0.95in}
\workquestion{4. In part (b), how do you rewrite \(\frac{n}{2^n}\) using \(\frac{n-1}{2^{n+1}}\) and \(\frac1{2^n}\)?}{1.0in}
\workquestion{5. Which geometric series value is needed?}{0.85in}
\end{selfexplanationbox}

\newpage
\begin{skeletonbox}{Homework 6, Problem 3}
\textbf{Part (a).}
\workquestion{Define \(s_n=\sum_{k=1}^n \frac{k-1}{2^{k+1}}\). Prove by induction that \(s_n=\)}{0.95in}
\workquestion{Base case \(n=1\):}{0.75in}
\workquestion{Inductive step using \(\frac{k-1}{2^{k+1}}=\frac{k}{2^k}-\frac{k+1}{2^{k+1}}\):}{1.0in}
\workquestion{Pass to the limit using \(\frac{n}{2^n}\to0\):}{0.9in}

\textbf{Part (b).}
\workquestion{Rewrite \(\sum_{k=1}^n\frac{k}{2^k}\) as \(2\sum\frac{k-1}{2^{k+1}}+\sum\frac1{2^k}\):}{0.95in}
\workquestion{Use part (a) and the geometric series to compute the value:}{0.85in}
\end{skeletonbox}

\newpage
\fullproofpage{Homework 6, Problem 3}{Use this page for part (a).}

\newpage
\fullproofpage{Homework 6, Problem 3}{Use this page for part (b).}

\newpage
\begin{solutionbox}{Official Solution (Professor)}
(a) Let \(s_n=\sum_{k=1}^n\frac{k-1}{2^{k+1}}\).

We will prove that
\[
  s_n=\frac12-\frac{n+1}{2^{n+1}}
  \qquad \text{for } n\ge1
\]
by induction. For the base case, we note that
\[
  s_1=\frac04=\frac12-\frac24.
\]
For the inductive step, assume that \(s_n=\frac12-\frac{n+1}{2^{n+1}}\) for some \(n\ge1\). Then
\[
  s_{n+1}
  =
  s_n+\frac{n}{2^{n+2}}
  =
  \frac12-\frac{n+1}{2^{n+1}}+\frac{n+1}{2^{n+1}}-\frac{n+2}{2^{n+2}}
  =
  \frac12-\frac{n+2}{2^{n+2}}.
\]
Together, we conclude that our formula for \(s_n\) holds for all \(n\ge1\) by induction.

Next, we claim that \(\lim\frac{n}{2^n}=0\). Notice that \(n\le(3/2)^n\) for all \(n\ge1\), and so
\[
  0\le \frac{n}{2^n}\le \frac{(3/2)^n}{2^n}=\left(\frac34\right)^n.
\]
The right-hand side converges to zero as \(n\to\infty\). Therefore, \(\frac{n}{2^n}\) also converges to zero by the squeeze theorem.

Altogether, we conclude
\[
  \lim s_n
  =
  \lim\left(\frac12-\frac{n+1}{2^{n+1}}\right)
  =
  \frac12.
\]
(In the last equality, we noted that \(\frac{n+1}{2^{n+1}}\) is a subsequence of \(\frac{n}{2^n}\), and therefore also converges to zero.) By definition of a convergent series, we conclude that
\[
  \sum_{n=1}^{\infty}\frac{n-1}{2^{n+1}}=\frac12.
\]

(b) Note that
\[
  \sum_{k=1}^n\frac{k}{2^k}
  =
  2\sum_{k=1}^n\frac{k-1}{2^{k+1}}
  +
  \sum_{k=1}^n\frac1{2^k}.
\]
By part (a), \(\sum\frac{k-1}{2^{k+1}}\) converges to \(\frac12\), and thus \(2\sum\frac{k-1}{2^{k+1}}\) converges to \(1\) by problem 2. From lecture, we know that the geometric series \(\sum\frac1{2^k}\) converges to \(\frac1{1-(1/2)}-1=1\). (Notice that we are starting at \(k=1\), not \(k=0\).) Together, we conclude \(\sum_{k=1}^{\infty}\frac{k}{2^k}\) converges to \(1+1=2\).
\end{solutionbox}

\newpage
\reflectionpage{Homework 6, Problem 3}

% ---------------------------------------------------------------------------
% HW6 Problem 4
% ---------------------------------------------------------------------------
\newpage
\subsection{Problem 4}

\begin{problemstatementbox}{Homework 6, Problem 4}
Prove that if \(\sum_{n=1}^{\infty}a_n\) converges absolutely and \(\{b_n\}_{n\ge1}\) is a bounded sequence, then the series \(\sum_{n=1}^{\infty}a_nb_n\) also converges. (Hint: Use the Cauchy criterion.)
\end{problemstatementbox}

\newpage
\firstdiagnosis{Homework 6, Problem 4}

\newpage
\recipecard{Homework 6, Problem 4}

\newpage
\begin{selfexplanationbox}{Homework 6, Problem 4}
\workquestion{1. What does absolute convergence of \(\sum a_n\) say about \(\sum |a_n|\)?}{0.95in}
\workquestion{2. What bound does boundedness of \(\{b_n\}\) provide?}{0.9in}
\workquestion{3. What is the Cauchy criterion for convergence of a series?}{1.0in}
\workquestion{4. How can \(\left|\sum_{k=m}^n a_kb_k\right|\) be controlled by \(\sum_{k=m}^n |a_k|\)?}{1.0in}
\workquestion{5. Why is \(M+1\) used in the denominator of the tail estimate?}{0.9in}
\end{selfexplanationbox}

\newpage
\begin{skeletonbox}{Homework 6, Problem 4}
\workquestion{Choose \(M\ge0\) such that \(|b_n|\le M\) for all \(n\):}{0.8in}
\workquestion{Use the Cauchy criterion for \(\sum|a_n|\) with tolerance \(\varepsilon/(M+1)\):}{1.0in}
\workquestion{For \(n\ge m\ge N\), estimate \(\left|\sum_{k=m}^n a_kb_k\right|\):}{1.05in}
\workquestion{Conclude the Cauchy criterion for \(\sum a_nb_n\):}{0.85in}
\end{skeletonbox}

\newpage
\fullproofpage{Homework 6, Problem 4}{Use this page for the complete proof.}

\newpage
\begin{solutionbox}{Student-Derived Solution (No official solution available)}
Since \(\sum_{n=1}^{\infty}a_n\) converges absolutely, the series
\[
  \sum_{n=1}^{\infty}|a_n|
\]
converges. By the Cauchy criterion for series, given \(\varepsilon>0\), there exists \(N\in\N\) such that
\[
  \sum_{k=m}^{n}|a_k|<\frac{\varepsilon}{M+1}
  \qquad \text{for all } n\ge m\ge N,
\]
where \(M\ge0\) is chosen so that
\[
  |b_n|\le M \qquad \text{for all } n\in\N.
\]
Such an \(M\) exists because \(\{b_n\}\) is bounded.

Now let \(n\ge m\ge N\). Then
\[
  \left|\sum_{k=m}^{n}a_kb_k\right|
  \le
  \sum_{k=m}^{n}|a_kb_k|
  =
  \sum_{k=m}^{n}|a_k||b_k|
  \le
  M\sum_{k=m}^{n}|a_k|.
\]
Therefore
\[
  \left|\sum_{k=m}^{n}a_kb_k\right|
  <
  M\cdot\frac{\varepsilon}{M+1}
  <
  \varepsilon.
\]
Thus, for every \(\varepsilon>0\), there exists \(N\in\N\) such that
\[
  \left|\sum_{k=m}^{n}a_kb_k\right|<\varepsilon
  \qquad \text{for all } n\ge m\ge N.
\]
By the Cauchy criterion for series, \(\sum_{n=1}^{\infty}a_nb_n\) converges.
\end{solutionbox}

\newpage
\reflectionpage{Homework 6, Problem 4}

% ---------------------------------------------------------------------------
% HW6 Problem 5
% ---------------------------------------------------------------------------
\newpage
\subsection{Problem 5}

\begin{problemstatementbox}{Homework 6, Problem 5}
\begin{enumerate}[label=(\alph*)]
  \item Prove that if \(\sum_{n=1}^{\infty}a_n\) converges absolutely, then the series \(\sum_{n=1}^{\infty}a_n^2\) also converges.
  \item Provide an example where \(\sum_{n=1}^{\infty}a_n\) diverges and \(\sum_{n=1}^{\infty}a_n^2\) converges.
\end{enumerate}
\end{problemstatementbox}

\newpage
\firstdiagnosis{Homework 6, Problem 5}

\newpage
\recipecard{Homework 6, Problem 5}

\newpage
\begin{selfexplanationbox}{Homework 6, Problem 5}
\workquestion{1. Why does absolute convergence imply \(a_n\to0\)?}{0.95in}
\workquestion{2. Why does \(a_n\to0\) imply \(|a_n|<1\) eventually?}{0.9in}
\workquestion{3. For \(|x|<1\), why is \(x^2\le |x|\)?}{0.95in}
\workquestion{4. Which comparison proves the tail of \(\sum a_n^2\) converges?}{0.95in}
\workquestion{5. What standard example makes \(\sum a_n\) diverge but \(\sum a_n^2\) converge?}{0.95in}
\end{selfexplanationbox}

\newpage
\begin{skeletonbox}{Homework 6, Problem 5}
\textbf{Part (a).}
\workquestion{From absolute convergence, write the convergent positive series:}{0.85in}
\workquestion{Use convergence of \(\sum a_n\) to get \(a_n\to0\), then choose \(N\) with \(|a_n|<1\):}{0.95in}
\workquestion{For \(n\ge N\), compare \(a_n^2\) with \(|a_n|\):}{0.9in}
\workquestion{Use the comparison test and finite initial terms:}{0.9in}

\textbf{Part (b).}
\workquestion{Choose \(a_n=\) and cite the harmonic and \(p\)-series facts:}{0.95in}
\end{skeletonbox}

\newpage
\fullproofpage{Homework 6, Problem 5}{Use this page for part (a).}

\newpage
\fullproofpage{Homework 6, Problem 5}{Use this page for part (b).}

\newpage
\begin{solutionbox}{Official Solution (Professor)}
(a) Note that \(x^2\le |x|\) for any \(x\in[-1,1]\). (I recommend that you draw the graphs of \(x^2\) and \(|x|\) to check that this is true.)

As \(\sum a_n\) converges, then we have \(\lim a_n=0\), and so there exists an \(N\) so that \(|a_n|<1\) for all \(n\ge N\). By the previous paragraph, this means that \(a_n^2\le |a_n|\) for all \(n\ge N\). As \(\sum |a_n|\) converges, then \(\sum_{n=N}^{\infty}a_n^2\) also converges by the comparison test. Adding the first \(N-1\) terms back in, we conclude \(\sum_{n=1}^{\infty}a_n^2\) converges as well.

(b) Consider \(a_n=\frac1n\). In lecture, we proved that \(\sum\frac1n\) diverges and \(\sum\frac1{n^2}\) converges.
\end{solutionbox}

\newpage
\reflectionpage{Homework 6, Problem 5}

% ---------------------------------------------------------------------------
% HW6 Problem 6
% ---------------------------------------------------------------------------
\newpage
\subsection{Problem 6}

\begin{problemstatementbox}{Homework 6, Problem 6}
For each of the following series, determine if it converges and prove your answer.
\[
  \text{(a) }\sum_{n=2}^{\infty}\frac{1}{[n+(-1)^n]^2}
  \qquad
  \text{(b) }\sum_{n=1}^{\infty}\left(\sqrt{n+1}-\sqrt n\right)
  \qquad
  \text{(c) }\sum_{n=1}^{\infty}(-1)^n
\]
\end{problemstatementbox}

\newpage
\firstdiagnosis{Homework 6, Problem 6}

\newpage
\recipecard{Homework 6, Problem 6}

\newpage
\begin{selfexplanationbox}{Homework 6, Problem 6}
\workquestion{1. In part (a), what lower bound can you put on \(n+(-1)^n\) for \(n\ge2\)?}{0.95in}
\workquestion{2. Which convergent \(p\)-series controls part (a)?}{0.85in}
\workquestion{3. In part (b), what is the \(N\)th partial sum after telescoping?}{0.95in}
\workquestion{4. Why does that partial sum sequence diverge to \(+\infty\)?}{0.95in}
\workquestion{5. In part (c), what necessary condition for convergence fails?}{0.9in}
\end{selfexplanationbox}

\newpage
\begin{skeletonbox}{Homework 6, Problem 6}
\textbf{Part (a).}
\workquestion{For \(n\ge2\), show \(n+(-1)^n\ge n-1\), then compare:}{0.95in}
\workquestion{Use the comparison test with \(\sum 1/(n-1)^2\):}{0.85in}

\textbf{Part (b).}
\workquestion{Let \(s_n=\sum_{k=1}^n(\sqrt{k+1}-\sqrt{k})\). Prove by induction/telescoping that \(s_n=\)}{0.95in}
\workquestion{Show \(s_n\to+\infty\), so the series diverges:}{0.85in}

\textbf{Part (c).}
\workquestion{Suppose the series converged. Then its terms would satisfy:}{0.8in}
\workquestion{Explain why \((-1)^n\) does not satisfy this:}{0.8in}
\end{skeletonbox}

\newpage
\fullproofpage{Homework 6, Problem 6}{Use this page for parts (a) and (b).}

\newpage
\fullproofpage{Homework 6, Problem 6}{Use this page for part (c) and final conclusions.}

\newpage
\begin{solutionbox}{Official Solution (Professor)}
(a) We will prove that this series converges. For \(n\ge2\), we have
\[
  \frac{1}{[n+(-1)^n]^2}\le \frac{1}{(n-1)^2}.
\]
From lecture, we know that the series
\[
  \sum_{n=2}^{\infty}\frac{1}{(n-1)^2}
  =
  \sum_{k=1}^{\infty}\frac1{k^2}
\]
converges. Therefore our series \(\sum \frac1{[n+(-1)^n]^2}\) also converges by the comparison test.

(b) We will prove that this series diverges. Let
\[
  s_n=\sum_{k=1}^n(\sqrt{k+1}-\sqrt{k})
\]
denote the \(n\)th partial sum.

Claim: \(s_n=\sqrt{n+1}-1\) for all \(n\ge1\). We will argue using induction.

\begin{itemize}
  \item Base case: For \(n=1\), we have
  \[
    s_1=\sqrt2-\sqrt1=\sqrt2-1.
  \]
  \item Inductive step: Suppose that \(s_n=\sqrt{n+1}-1\) for some \(n\ge1\). Then the formula for \(s_{n+1}\) holds, since
  \[
    s_{n+1}
    =
    s_n+(\sqrt{n+2}-\sqrt{n+1})
    =
    (\sqrt{n+1}-1)+(\sqrt{n+2}-\sqrt{n+1})
    =
    \sqrt{n+2}-1.
  \]
\end{itemize}
By induction, we conclude that \(s_n=\sqrt{n+1}-1\) holds for all \(n\ge1\).

Notice that the sequence \(s_n=\sqrt{n+1}-1\) diverges to \(+\infty\). Indeed, given \(M>0\), pick \(N\in\N\) so that \(N>(M+1)^2\), and then for any \(n\ge N\) we have
\[
  \sqrt{n+1}-1\ge \sqrt N-1>\sqrt{(M+1)^2}-1=M.
\]
As the sequence \(\{s_n\}\) diverges, then the series \(\sum(\sqrt{n+1}-\sqrt n)\) diverges.

(c) We will prove that this series diverges. Suppose towards a contradiction that \(\sum(-1)^n\) converges. Then \(\lim(-1)^n=0\) by Lecture 16. However, we know that the sequence \(\{(-1)^n\}\) diverges by Lecture 10, and so this is a contradiction.
\end{solutionbox}

\newpage
\reflectionpage{Homework 6, Problem 6}

\section{Homework 7}

% ---------------------------------------------------------------------------
% HW7 Problem 1
% ---------------------------------------------------------------------------
\newpage
\subsection{Problem 1}

\begin{problemstatementbox}{Homework 7, Problem 1}
For each of the following series, determine if it converges and prove your answer.
\[
  \text{(a) }\sum_{n=2}^{\infty}\frac1{\log n}
  \qquad
  \text{(b) }\sum_{n=1}^{\infty}\frac{n-1}{n^2}
  \qquad
  \text{(c) }\sum_{n=1}^{\infty}\frac{n^2}{n!}
\]
\end{problemstatementbox}

\newpage
\firstdiagnosis{Homework 7, Problem 1}

\newpage
\recipecard{Homework 7, Problem 1}

\newpage
\begin{selfexplanationbox}{Homework 7, Problem 1}
\workquestion{1. Which convergence test is adapted to \(\sum 1/\log n\), and why?}{1.0in}
\workquestion{2. In part (a), why do the dyadic terms fail the term test?}{0.95in}
\workquestion{3. In part (b), how does \((n-1)/n^2\) compare to \(1/(2n)\)?}{0.95in}
\workquestion{4. In part (c), why is the ratio test natural for \(n^2/n!\)?}{0.9in}
\workquestion{5. Which parts converge and which diverge?}{0.9in}
\end{selfexplanationbox}

\newpage
\begin{skeletonbox}{Homework 7, Problem 1}
\textbf{Part (a).}
\workquestion{Reindex with \(a_k=1/\log(k+1)\) and state the dyadic criterion:}{0.95in}
\workquestion{Estimate \(2^n a_{2^n}\) from below and explain why the terms do not go to zero:}{1.0in}

\textbf{Part (b).}
\workquestion{For \(n\ge2\), prove \(\frac{n-1}{n^2}\ge\frac1{2n}\):}{0.95in}
\workquestion{Use comparison with the harmonic series:}{0.85in}

\textbf{Part (c).}
\workquestion{Set \(a_n=\frac{n^2}{n!}\). Compute \(a_{n+1}/a_n\):}{0.9in}
\workquestion{Apply the ratio test:}{0.8in}
\end{skeletonbox}

\newpage
\fullproofpage{Homework 7, Problem 1}{Use this page for parts (a) and (b).}

\newpage
\fullproofpage{Homework 7, Problem 1}{Use this page for part (c) and final conclusions.}

\newpage
\begin{solutionbox}{Official Solution (Professor)}
(a) We'll use the dyadic criterion. Consider the series \(\sum_{k=1}^{\infty}a_k\) with \(a_k=\frac1{\log(k+1)}\) for \(k\ge1\). (Strictly speaking, in order to apply the dyadic criterion we need a series that starts at \(n=1\), so we reindexed.)

Then \(\{a_k\}_{k\ge1}\) is a decreasing sequence of nonnegative numbers. Therefore, by the dyadic criterion, \(\sum_{k=1}^{\infty}a_k\) converges if and only if the series \(\sum_{n=0}^{\infty}2^n a_{2^n}\) converges. Note that
\[
  2^n a_{2^n}
  =
  \frac{2^n}{\log(2^n+1)}
  \ge
  \frac{2^n}{\log(2^{n+1})}
  =
  \frac1{\log2}\cdot\frac{2^n}{n+1}.
\]
Now, the sequence \(\frac{2^n}{n+1}\) diverges to \(+\infty\), and so it cannot converge to \(0\). Multiplying this by the constant \(\frac1{\log2}\), we see that \(2^n a_{2^n}\) also cannot converge to \(0\). Therefore the series \(\sum 2^n a_{2^n}\) diverges, and thus \(\sum a_k\) diverges as well.

(b) We'll use the comparison test. For \(n\ge2\), we have
\[
  n-1\ge n-\frac n2=\frac n2
  \qquad \Longrightarrow \qquad
  \frac{n-1}{n^2}\ge \frac{n/2}{n^2}=\frac1{2n}.
\]
From lecture we know that the series \(\sum \frac1n\) diverges, and thus \(\sum_{n=2}^{\infty}\frac1{2n}\) diverges as well. Therefore, by the comparison test, \(\sum_{n=2}^{\infty}\frac{n-1}{n^2}\) also diverges. Note that \(\frac{n-1}{n^2}\) is equal to zero for \(n=1\), so \(\sum_{n=2}^{\infty}\frac{n-1}{n^2}\) and \(\sum_{n=1}^{\infty}\frac{n-1}{n^2}\) are the same series.

(c) We'll use the ratio test. Set \(a_n=\frac{n^2}{n!}\). Then
\[
  \frac{a_{n+1}}{a_n}
  =
  \frac{(n+1)^2\cdot n!}{n^2\cdot (n+1)!}
  =
  \frac{n+1}{n^2}
  \longrightarrow 0
  \qquad \text{as } n\to\infty.
\]
Therefore
\[
  \limsup\left|\frac{a_{n+1}}{a_n}\right|
  =
  \lim\frac{a_{n+1}}{a_n}
  =
  0<1,
\]
and so the series converges by the ratio test.
\end{solutionbox}

\newpage
\reflectionpage{Homework 7, Problem 1}

% ---------------------------------------------------------------------------
% HW7 Problem 2
% ---------------------------------------------------------------------------
\newpage
\subsection{Problem 2}

\begin{problemstatementbox}{Homework 7, Problem 2}
For each of the following series, determine if it converges and prove your answer.
\[
  \text{(a) }\sum_{n=1}^{\infty}\frac{n^4}{4^n}
  \qquad
  \text{(b) }\sum_{n=1}^{\infty}\frac{n!}{n^4+3}
  \qquad
  \text{(c) }\sum_{n=1}^{\infty}\frac{2^n}{n!}
\]
\end{problemstatementbox}

\newpage
\firstdiagnosis{Homework 7, Problem 2}

\newpage
\recipecard{Homework 7, Problem 2}

\newpage
\begin{selfexplanationbox}{Homework 7, Problem 2}
\workquestion{1. Why does a polynomial-over-exponential series suggest the ratio test?}{0.95in}
\workquestion{2. In part (b), what necessary condition for convergence should be checked?}{0.95in}
\workquestion{3. How does factorial growth compare to \(n^4+3\)?}{0.95in}
\workquestion{4. In part (c), what ratio do successive terms have?}{0.9in}
\workquestion{5. Which of these three series fail because the terms do not go to \(0\)?}{0.9in}
\end{selfexplanationbox}

\newpage
\begin{skeletonbox}{Homework 7, Problem 2}
\textbf{Part (a).}
\workquestion{Set \(a_n=\frac{n^4}{4^n}\). Compute \(a_{n+1}/a_n\) and its limit:}{1.0in}
\workquestion{Apply the ratio test:}{0.75in}

\textbf{Part (b).}
\workquestion{Show \(\frac{n^4+3}{n!}\to0\), for example using \(n!\ge n(n-1)(n-2)(n-3)(n-4)\):}{1.1in}
\workquestion{Use HW5\#1 to conclude \(\frac{n!}{n^4+3}\to+\infty\), hence the series diverges:}{0.95in}

\textbf{Part (c).}
\workquestion{Set \(a_n=\frac{2^n}{n!}\). Compute \(a_{n+1}/a_n\) and apply the ratio test:}{0.95in}
\end{skeletonbox}

\newpage
\fullproofpage{Homework 7, Problem 2}{Use this page for parts (a) and (b).}

\newpage
\fullproofpage{Homework 7, Problem 2}{Use this page for part (c) and final conclusions.}

\newpage
\begin{solutionbox}{Official Solution (Professor)}
(a) We'll use the ratio test. Set \(a_n=\frac{n^4}{4^n}\). Then
\[
  \frac{a_{n+1}}{a_n}
  =
  \frac{(n+1)^4\cdot 4^n}{n^4\cdot 4^{n+1}}
  =
  \frac{(n+1)^4}{4n^4}
  \longrightarrow \frac14
  \qquad \text{as } n\to\infty.
\]
Therefore \(\limsup\left|\frac{a_{n+1}}{a_n}\right|=\frac14<1\), and so the series converges by the ratio test.

(b) This series diverges. In brief, the sequence \(\frac{n!}{n^4+3}\) diverges to \(+\infty\), and so it cannot converge to zero. More concretely, notice that for \(n\ge5\) we have
\[
  n!\ge n(n-1)(n-2)(n-3)(n-4),
\]
and so
\[
  \frac{n^4+3}{n!}
  \le
  \frac{n^4+3}{n(n-1)(n-2)(n-3)(n-4)}
  \longrightarrow 0
  \qquad \text{as } n\to\infty.
\]
Therefore the sequence \(\frac{n!}{n^4+3}\to+\infty\) by HW5\#1. In particular, there must exist an \(N\) so that \(\frac{n!}{n^4+3}\ge1\) for all \(n\ge N\), and so \(\frac{n!}{n^4+3}\) cannot converge to \(0\).

(c) We'll use the ratio test. Set \(a_n=\frac{2^n}{n!}\). Then
\[
  \frac{a_{n+1}}{a_n}
  =
  \frac{2^{n+1}\cdot n!}{(n+1)!\cdot 2^n}
  =
  \frac2{n+1}
  \longrightarrow 0
  \qquad \text{as } n\to\infty.
\]
Therefore \(\limsup\left|\frac{a_{n+1}}{a_n}\right|=0<1\), and so the series converges by the ratio test.
\end{solutionbox}

\newpage
\reflectionpage{Homework 7, Problem 2}

% ---------------------------------------------------------------------------
% HW7 Problem 3
% ---------------------------------------------------------------------------
\newpage
\subsection{Problem 3}

\begin{problemstatementbox}{Homework 7, Problem 3}
For each of the following series, determine if it converges and prove your answer.
\[
  \text{(a) }\sum_{n=1}^{\infty}\frac1{n^n}
  \qquad
  \text{(b) }\sum_{n=1}^{\infty}\frac{n!}{n^n}
  \qquad
  \text{(c) }\sum_{n=1}^{\infty}\frac{(-1)^n n!}{2^n}
\]
\end{problemstatementbox}

\newpage
\firstdiagnosis{Homework 7, Problem 3}

\newpage
\recipecard{Homework 7, Problem 3}

\newpage
\begin{selfexplanationbox}{Homework 7, Problem 3}
\workquestion{1. Why does \(\frac1{n^n}\) suggest the root test?}{0.9in}
\workquestion{2. What lecture result gives \((n!)^{1/n}/n\to 1/e\)?}{0.95in}
\workquestion{3. Why does \(1/e<1\) prove convergence in part (b)?}{0.9in}
\workquestion{4. In part (c), why should you check whether the terms go to \(0\)?}{0.95in}
\workquestion{5. How does HW5\#1 convert \(2^n/n!\to0\) into \(n!/2^n\to+\infty\)?}{1.0in}
\end{selfexplanationbox}

\newpage
\begin{skeletonbox}{Homework 7, Problem 3}
\textbf{Part (a).}
\workquestion{Set \(a_n=1/n^n\). Compute \(|a_n|^{1/n}\):}{0.9in}
\workquestion{Apply the root test:}{0.75in}

\textbf{Part (b).}
\workquestion{Set \(a_n=n!/n^n\). Compute \(|a_n|^{1/n}\):}{0.9in}
\workquestion{Use Lecture 15 to identify the limit and apply the root test:}{0.95in}

\textbf{Part (c).}
\workquestion{Use Problem 2(c) to know \(\frac{2^n}{n!}\to0\):}{0.85in}
\workquestion{Apply HW5\#1 to show \(\frac{n!}{2^n}\to+\infty\), then conclude the terms do not go to \(0\):}{1.0in}
\end{skeletonbox}

\newpage
\fullproofpage{Homework 7, Problem 3}{Use this page for parts (a) and (b).}

\newpage
\fullproofpage{Homework 7, Problem 3}{Use this page for part (c) and final conclusions.}

\newpage
\begin{solutionbox}{Official Solution (Professor)}
(a) We'll use the root test. Set \(a_n=\frac1{n^n}\). Then
\[
  |a_n|^{1/n}=\frac1n \longrightarrow 0
  \qquad \text{as } n\to\infty.
\]
Therefore \(\limsup |a_n|^{1/n}=0<1\), and so the series converges by the root test.

(b) We'll use the root test. Set \(a_n=\frac{n!}{n^n}\), and consider the sequence
\[
  |a_n|^{1/n}
  =
  \left(\frac{n!}{n^n}\right)^{1/n}
  =
  \frac{(n!)^{1/n}}{n}.
\]
In Lecture 15 we showed that this sequence converges to \(1/e\). Therefore
\[
  \limsup |a_n|^{1/n}=1/e<1,
\]
and so the series converges by the root test.

(c) This series diverges, because the sequence \(\frac{(-1)^n n!}{2^n}\) does not converge to \(0\). Specifically, we know that the sequence \(\frac{2^n}{n!}\) converges to \(0\). (Indeed, we proved in problem 2c that the series \(\sum\frac{2^n}{n!}\) converges, and so its terms must converge to \(0\).) So, by HW5\#1, the sequence \(\frac{n!}{2^n}\) diverges to \(+\infty\). In particular, there exists an \(N\) so that \(\frac{n!}{2^n}\ge1\) for all \(n\ge N\). Therefore \(\left|\frac{(-1)^n n!}{2^n}\right|\ge1\) for all \(n\ge N\), and so the terms \(\frac{(-1)^n n!}{2^n}\) in our series cannot converge to \(0\).
\end{solutionbox}

\newpage
\reflectionpage{Homework 7, Problem 3}

% ---------------------------------------------------------------------------
% HW7 Problem 4
% ---------------------------------------------------------------------------
\newpage
\subsection{Problem 4}

\begin{problemstatementbox}{Homework 7, Problem 4}
Consider the series
\[
  \sum_{n=1}^{\infty}2^{(-1)^n-n}
  =
  \frac14+\frac12+\frac1{16}+\frac18+\frac1{64}+\frac1{32}+\cdots.
\]
\begin{enumerate}[label=(\alph*)]
  \item Show that the ratio test is inconclusive.
  \item Prove that the series converges using the root test.
\end{enumerate}
\end{problemstatementbox}

\newpage
\firstdiagnosis{Homework 7, Problem 4}

\newpage
\recipecard{Homework 7, Problem 4}

\newpage
\begin{selfexplanationbox}{Homework 7, Problem 4}
\workquestion{1. What are the first few terms of \(a_{n+1}/a_n\)?}{0.95in}
\workquestion{2. Why does the ratio sequence have two subsequential limits?}{0.95in}
\workquestion{3. What does it mean for the ratio test to be inconclusive?}{0.9in}
\workquestion{4. Why does the root test handle \(2^{(-1)^n-n}\) better than the ratio test?}{0.95in}
\workquestion{5. What is \(\lim |a_n|^{1/n}\)?}{0.9in}
\end{selfexplanationbox}

\newpage
\begin{skeletonbox}{Homework 7, Problem 4}
\textbf{Part (a).}
\workquestion{Set \(a_n=2^{(-1)^n-n}\). Compute the first few ratios \(|a_{n+1}/a_n|\):}{1.0in}
\workquestion{Identify the liminf and limsup of the ratio sequence:}{0.9in}
\workquestion{State why this makes the ratio test inconclusive:}{0.75in}

\textbf{Part (b).}
\workquestion{Compute \(|a_n|^{1/n}=2^{((-1)^n/n)-1}\):}{0.9in}
\workquestion{Use \((-1)^n/n\to0\) to find the root-test limit:}{0.9in}
\workquestion{Conclude convergence by the root test:}{0.75in}
\end{skeletonbox}

\newpage
\fullproofpage{Homework 7, Problem 4}{Use this page for part (a).}

\newpage
\fullproofpage{Homework 7, Problem 4}{Use this page for part (b).}

\newpage
\begin{solutionbox}{Official Solution (Professor)}
(a) Set \(a_n=2^{(-1)^n-n}\). Then the sequence \(\left|\frac{a_{n+1}}{a_n}\right|\) is
\[
  \frac{1/2}{1/4}=2,\qquad
  \frac{1/16}{1/2}=\frac18,\qquad
  \frac{1/8}{1/16}=2,\qquad
  \frac{1/64}{1/8}=\frac18,\ldots
\]
In other words, the even terms of this sequence are \(\frac18\) and the odd terms are \(2\). Therefore
\[
  \liminf\left|\frac{a_{n+1}}{a_n}\right|=\frac18,
  \qquad
  \limsup\left|\frac{a_{n+1}}{a_n}\right|=2,
\]
and so
\[
  \liminf\left|\frac{a_{n+1}}{a_n}\right|
  \le 1
  \le
  \limsup\left|\frac{a_{n+1}}{a_n}\right|.
\]

(b) Set \(a_n=2^{(-1)^n-n}\). Then
\[
  |a_n|^{1/n}
  =
  2^{\frac{(-1)^n}{n}-1}.
\]
Since \(\frac{(-1)^n}{n}\to0\), we get
\[
  \lim_{n\to\infty}|a_n|^{1/n}=\frac12<1.
\]
Therefore the series converges by the root test.
\end{solutionbox}

\newpage
\reflectionpage{Homework 7, Problem 4}

% ---------------------------------------------------------------------------
% HW7 Problem 5
% ---------------------------------------------------------------------------
\newpage
\subsection{Problem 5}

\begin{problemstatementbox}{Homework 7, Problem 5}
Below are two series \(\sum_{n=1}^{\infty}a_n\) and the numbers \(s\in\R\) to which they converge. For each of these series, what is the set of all numbers \(\widetilde{s}\in\R\) such that there exists a rearrangement that satisfies \(\sum_{n=1}^{\infty}\widetilde{a}_n=\widetilde{s}\)? Prove your answer.
\[
  \text{(a) }\sum_{n=1}^{\infty}\frac{(-1)^n}{n}=-\log2
  \qquad
  \text{(b) }\sum_{n=1}^{\infty}\frac{(-1)^n}{n^2}=-\frac{\pi^2}{12}
\]
(Hint: Your proof should be quite short, because we already did most of the work in lecture.)
\end{problemstatementbox}

\newpage
\firstdiagnosis{Homework 7, Problem 5}

\newpage
\recipecard{Homework 7, Problem 5}

\newpage
\begin{selfexplanationbox}{Homework 7, Problem 5}
\workquestion{1. What is the difference between conditional and absolute convergence for rearrangements?}{1.0in}
\workquestion{2. Why is the alternating harmonic series conditionally convergent?}{0.95in}
\workquestion{3. What does Riemann's rearrangement theorem say about possible sums?}{0.95in}
\workquestion{4. Why does \(\sum (-1)^n/n^2\) converge absolutely?}{0.9in}
\workquestion{5. What theorem says every rearrangement of an absolutely convergent series has the same sum?}{0.95in}
\end{selfexplanationbox}

\newpage
\begin{skeletonbox}{Homework 7, Problem 5}
\textbf{Part (a).}
\workquestion{Show \(\sum (-1)^n/n\) converges but not absolutely:}{1.0in}
\workquestion{Apply Riemann's rearrangement result with \(\alpha=\widetilde{s}=\beta\):}{0.95in}
\workquestion{Conclude the set of possible rearranged sums is:}{0.75in}

\textbf{Part (b).}
\workquestion{Show \(\sum (-1)^n/n^2\) converges absolutely:}{0.9in}
\workquestion{Apply the theorem on rearrangements of absolutely convergent series:}{0.95in}
\workquestion{Conclude the set of possible rearranged sums is:}{0.75in}
\end{skeletonbox}

\newpage
\fullproofpage{Homework 7, Problem 5}{Use this page for part (a).}

\newpage
\fullproofpage{Homework 7, Problem 5}{Use this page for part (b).}

\newpage
\begin{solutionbox}{Official Solution (Professor)}
(a) We claim that the set is \(\R\). Note that:
\begin{itemize}
  \item \(\sum \frac{(-1)^n}{n}\) converges: by the alternating series test, since \(\{\frac1n\}\) is a decreasing sequence that converges to \(0\).
  \item \(\sum \frac{(-1)^n}{n}\) does not converge absolutely: since \(\left|\frac{(-1)^n}{n}\right|=\frac1n\), and we know \(\sum\frac1n\) diverges by Lecture 11.
\end{itemize}
Fix \(\widetilde{s}\in\R\). Applying the Riemann rearrangement result (with \(\alpha=\widetilde{s}=\beta\)), we deduce that there exists a rearrangement \(\sum \widetilde{a}_n\) that converges to \(\widetilde{s}\). This works for any \(\widetilde{s}\in\R\), and so \(\R\) is the desired set.

(b) We claim that the set is \(\{-\frac{\pi^2}{12}\}\). Note that \(\sum\frac{(-1)^n}{n^2}\) converges absolutely, since \(\left|\frac{(-1)^n}{n^2}\right|=\frac1{n^2}\) and we know \(\sum\frac1{n^2}\) converges by Lecture 16. Suppose that \(\sum \widetilde{a}_n\) is a rearrangement that converges to some \(\widetilde{s}\in\R\). Applying the second theorem from Lecture 20, we deduce that \(\widetilde{s}=-\frac{\pi^2}{12}\). This works for any rearrangement, so \(-\frac{\pi^2}{12}\) is the only value that can be attained.
\end{solutionbox}

\newpage
\reflectionpage{Homework 7, Problem 5}

% ---------------------------------------------------------------------------
% HW7 Problem 6
% ---------------------------------------------------------------------------
\newpage
\subsection{Problem 6}

\begin{problemstatementbox}{Homework 7, Problem 6}
Prove or disprove the following statements:
\begin{enumerate}[label=(\alph*)]
  \item If \(\sum_{n=1}^{\infty}a_n\) converges, then \(\sum_{n=1}^{\infty}a_{2n}\) converges.
  \item If \(\sum_{n=1}^{\infty}a_n\) converges absolutely, then \(\sum_{n=1}^{\infty}a_{2n}\) converges absolutely.
\end{enumerate}
\end{problemstatementbox}

\newpage
\firstdiagnosis{Homework 7, Problem 6}

\newpage
\recipecard{Homework 7, Problem 6}

\newpage
\begin{selfexplanationbox}{Homework 7, Problem 6}
\workquestion{1. Why can ordinary convergence depend on cancellation between odd and even terms?}{1.0in}
\workquestion{2. What standard alternating example disproves part (a)?}{0.95in}
\workquestion{3. What is \(a_{2n}\) for that example?}{0.95in}
\workquestion{4. Why does absolute convergence survive passing to a subseries?}{0.95in}
\workquestion{5. How does the Cauchy criterion prove part (b)?}{0.95in}
\end{selfexplanationbox}

\newpage
\begin{skeletonbox}{Homework 7, Problem 6}
\textbf{Part (a).}
\workquestion{Choose \(a_n=\) so that \(\sum a_n\) is alternating harmonic:}{0.85in}
\workquestion{Compute \(a_{2n}\) and identify the divergent series:}{0.95in}
\workquestion{State the counterexample conclusion:}{0.75in}

\textbf{Part (b).}
\workquestion{Assume \(\sum|a_n|\) converges. Use the Cauchy criterion to choose \(N\):}{0.95in}
\workquestion{For \(n\ge m\ge N\), compare \(\sum_{k=m}^n |a_{2k}|\) to a tail of \(\sum |a_j|\):}{1.05in}
\workquestion{Conclude \(\sum |a_{2n}|\) converges:}{0.8in}
\end{skeletonbox}

\newpage
\fullproofpage{Homework 7, Problem 6}{Use this page for part (a).}

\newpage
\fullproofpage{Homework 7, Problem 6}{Use this page for part (b).}

\newpage
\begin{solutionbox}{Official Solution (Professor)}
(a) Disprove: Consider the series \(a_n=\frac{(-1)^n}{n}\). Then:
\begin{itemize}
  \item \(\sum a_n\) converges: by the alternating series test, since \(\frac1n\) is a decreasing sequence that converges to \(0\).
  \item \(\sum a_{2n}=\sum \frac1{2n}\) diverges: This is \(\frac12\) times the series \(\sum\frac1n\), which we know diverges by lecture.
\end{itemize}

(b) Prove: We will use the Cauchy criterion. Fix \(\varepsilon>0\). As \(\sum |a_n|\) converges, then there exists an \(N\) so that
\[
  \sum_{k=m}^{n}|a_k|<\varepsilon
  \qquad \text{for all } n\ge m\ge N.
\]
Then for \(n\ge m\ge N\), we have
\[
  \sum_{k=m}^{n}|a_{2k}|
  \le
  \sum_{j=2m}^{2n}|a_j|
  <
  \varepsilon.
\]
(Note that first sum above consists of only even terms \(|a_{2m}|+|a_{2m+2}|+\cdots+|a_{2n}|\), which differs from the second sum \(|a_{2m}|+|a_{2m+1}|+|a_{2m+2}|+\cdots+|a_{2n}|\) by only the odd terms, each of which is nonnegative.) Therefore \(\sum |a_{2n}|\) converges by the Cauchy criterion.
\end{solutionbox}

\newpage
\reflectionpage{Homework 7, Problem 6}

\section{Homework 8}

% ---------------------------------------------------------------------------
% HW8 Problem 1
% ---------------------------------------------------------------------------
\newpage
\subsection{Problem 1}

\begin{problemstatementbox}{Homework 8, Problem 1}
Consider the metric space \((\R,d)\) with the usual metric \(d(x,y)=|x-y|\). For each of the following sets \(A\subseteq\R\), find its interior and prove your answer.
\begin{enumerate}[label=(\alph*)]
  \item \(A=[0,1)\)
  \item \(A=\Z\)
  \item \(A=\{x\in\Q:0<x<1\}\)
\end{enumerate}
\end{problemstatementbox}

\newpage
\firstdiagnosis{Homework 8, Problem 1}

\newpage
\recipecard{Homework 8, Problem 1}

\newpage
\begin{selfexplanationbox}{Homework 8, Problem 1}
\workquestion{1. What does it mean for \(x\) to be an interior point of \(A\)?}{0.95in}
\workquestion{2. Why is every point of \((0,1)\) interior to \([0,1)\)?}{0.95in}
\workquestion{3. Why is \(0\) not an interior point of \([0,1)\)?}{0.9in}
\workquestion{4. Why can no open interval in \(\R\) be contained in \(\Z\)?}{0.95in}
\workquestion{5. How does density of irrational numbers show \(\{x\in\Q:0<x<1\}^\circ=\varnothing\)?}{1.0in}
\end{selfexplanationbox}

\newpage
\begin{skeletonbox}{Homework 8, Problem 1}
\textbf{Part (a).}
\workquestion{Show \((0,1)\subseteq A^\circ\) by choosing \(r=\frac12\min\{x,1-x\}\):}{1.0in}
\workquestion{Show \(0\notin A^\circ\) by finding a point in \((-r,r)\setminus[0,1)\):}{0.95in}
\workquestion{Conclude \(A^\circ=(0,1)\):}{0.75in}

\textbf{Part (b).}
\workquestion{Fix \(n\in\Z\) and \(r>0\). Find a noninteger point in \((n-r,n+r)\):}{0.95in}
\workquestion{Conclude \(\Z^\circ=\varnothing\):}{0.75in}

\textbf{Part (c).}
\workquestion{Fix \(a\in A\) and \(r>0\). Use density of \(\R\setminus\Q\) to find \(y\in(a-r,a+r)\setminus A\):}{0.95in}
\end{skeletonbox}

\newpage
\fullproofpage{Homework 8, Problem 1}{Use this page for parts (a) and (b).}

\newpage
\fullproofpage{Homework 8, Problem 1}{Use this page for part (c) and final conclusions.}

\newpage
\begin{solutionbox}{Official Solution (Professor)}
(a) Claim: \(A^\circ=(0,1)\). As \(A^\circ\) is the largest open subset of \(A\), and \((0,1)\) is one open subset, we must have \((0,1)\subseteq A^\circ\subseteq[0,1)\). Note that \(0\) is not an interior point: given \(r>0\) we have \((-r,r)\not\subseteq[0,1)\), since \(-\frac r2\) is a point in \((-r,r)\) that is not in \([0,1)\).

(b) Claim: \(A^\circ=\varnothing\). Fix \(n\in\Z\); we will show that \(n\) is not an interior point. Given \(r>0\) we have \((n-r,n+r)\not\subseteq\Z\), since \(\min\{n+\frac r2,n+\frac12\}\) is in \((n-r,n+r)\) but not in \(\Z\).

(c) Claim: \(A^\circ=\varnothing\). Fix \(a\in A\); we will show that \(a\) is not an interior point. As the irrational numbers are dense in \(\R\), then for any \(r>0\) we can find a point \(y\in\Q^c\) between \(a-r\) and \(a\). This \(y\) is a point in \((a-r,a+r)\) that is not in \(A\), and so \((a-r,a+r)\not\subseteq A\).
\end{solutionbox}

\newpage
\reflectionpage{Homework 8, Problem 1}

% ---------------------------------------------------------------------------
% HW8 Problem 2
% ---------------------------------------------------------------------------
\newpage
\subsection{Problem 2}

\begin{problemstatementbox}{Homework 8, Problem 2}
Recall that \(d_1\) and \(d_\infty\) are both metrics on \(\R^2\).
\begin{enumerate}[label=(\alph*)]
  \item Prove that \(d_\infty(x,y)\le d_1(x,y)\le 2d_\infty(x,y)\) for all \(x,y\in\R^2\).
  \item Prove that for any \(x\in\R^2\) and \(r>0\) we have
  \[
    \{y\in\R^2:d_\infty(x,y)<r/2\}
    \subseteq
    \{y\in\R^2:d_1(x,y)<r\}
    \subseteq
    \{y\in\R^2:d_\infty(x,y)<r\}.
  \]
  I recommend that you draw a picture of these three sets.
  \item Prove that a set \(A\subseteq\R^2\) is open in the metric space \((\R^2,d_1)\) if and only if it is open in the metric space \((\R^2,d_\infty)\).
\end{enumerate}
\end{problemstatementbox}

\newpage
\firstdiagnosis{Homework 8, Problem 2}

\newpage
\recipecard{Homework 8, Problem 2}

\newpage
\begin{selfexplanationbox}{Homework 8, Problem 2}
\workquestion{1. What are the formulas for \(d_1(x,y)\) and \(d_\infty(x,y)\) in coordinates?}{0.95in}
\workquestion{2. Why is the maximum of two nonnegative numbers at most their sum?}{0.9in}
\workquestion{3. Why is the sum of two nonnegative numbers at most twice their maximum?}{0.95in}
\workquestion{4. How do metric inequalities translate into ball containments?}{1.0in}
\workquestion{5. How do the ball containments prove the two metrics define the same open sets?}{1.0in}
\end{selfexplanationbox}

\newpage
\begin{skeletonbox}{Homework 8, Problem 2}
\textbf{Part (a).}
\workquestion{Write \(x=(x_1,x_2)\), \(y=(y_1,y_2)\), and expand \(d_1,d_\infty\):}{0.95in}
\workquestion{Prove \(d_\infty\le d_1\):}{0.85in}
\workquestion{Prove \(d_1\le2d_\infty\):}{0.85in}

\textbf{Part (b).}
\workquestion{If \(d_\infty(x,y)<r/2\), use part (a) to prove \(d_1(x,y)<r\):}{0.9in}
\workquestion{If \(d_1(x,y)<r\), use part (a) to prove \(d_\infty(x,y)<r\):}{0.9in}

\textbf{Part (c).}
\workquestion{For \(d_\infty\)-open \(\Rightarrow\) \(d_1\)-open, use \(B_r^1(a)\subseteq B_r^\infty(a)\):}{0.9in}
\workquestion{For \(d_1\)-open \(\Rightarrow\) \(d_\infty\)-open, use \(B_{r/2}^\infty(a)\subseteq B_r^1(a)\):}{0.9in}
\end{skeletonbox}

\newpage
\fullproofpage{Homework 8, Problem 2}{Use this page for parts (a) and (b).}

\newpage
\fullproofpage{Homework 8, Problem 2}{Use this page for part (c).}

\newpage
\begin{solutionbox}{Official Solution (Professor)}
(a) Fix \(x=(x_1,x_2)\) and \(y=(y_1,y_2)\) in \(\R^2\). Then
\[
  |x_1-y_1|\le |x_1-y_1|+|x_2-y_2|=d_1(x,y),
\]
\[
  |x_2-y_2|\le |x_1-y_1|+|x_2-y_2|=d_1(x,y),
\]
and so
\[
  \max\{|x_1-y_1|,|x_2-y_2|\}\le d_1(x,y).
\]
This proves the first inequality: \(d_\infty(x,y)\le d_1(x,y)\). On the other hand, note that
\[
  |x_1-y_1|\le \max\{|x_1-y_1|,|x_2-y_2|\}=d_\infty(x,y),
\]
\[
  |x_2-y_2|\le \max\{|x_1-y_1|,|x_2-y_2|\}=d_\infty(x,y),
\]
and so
\[
  |x_1-y_1|+|x_2-y_2|\le 2d_\infty(x,y).
\]
This proves that \(d_1(x,y)\le2d_\infty(x,y)\).

(b) Let's start with the first containment:
\[
  \{y\in\R^2:d_\infty(x,y)<r/2\}
  \subseteq
  \{y\in\R^2:d_1(x,y)<r\}.
\]
Fix \(y\in\R^2\) so that \(d_\infty(x,y)<r/2\). Then by the second inequality in part (a) we have
\[
  d_1(x,y)\le2d_\infty(x,y)<2\cdot\frac r2=r,
\]
and so \(d_1(x,y)<r\).

Next, we turn to the second containment:
\[
  \{y\in\R^2:d_1(x,y)<r\}
  \subseteq
  \{y\in\R^2:d_\infty(x,y)<r\}.
\]
Fix \(y\in\R^2\) so that \(d_1(x,y)<r\). Then by the first inequality in part (a), we have
\[
  d_\infty(x,y)\le d_1(x,y)<r,
\]
and so \(d_\infty(x,y)<r\).

(c) \(\Leftarrow\) Suppose that \(A\) is open in \((\R^2,d_\infty)\). Fix \(a\in A\). Then there exists \(r>0\) so that
\[
  B_r^\infty(a)=\{x\in\R^2:d_\infty(x,a)<r\}\subseteq A.
\]
Note that \(B_r^1(a)\subseteq B_r^\infty(a)\) by part (b), and so \(B_r^1(a)\subseteq B_r^\infty(a)\subseteq A\). As \(a\in A\) was arbitrary, we conclude that \(A\) is open in \((\R^2,d_1)\).

\(\Rightarrow\) Suppose that \(A\) is open in \((\R^2,d_1)\). Fix \(a\in A\). Then there exists \(r>0\) so that \(B_r^1(a)\subseteq A\). Note that \(B_{r/2}^\infty(a)\subseteq B_r^1(a)\) by part (b), and so \(B_{r/2}^\infty(a)\subseteq B_r^1(a)\subseteq A\). As \(a\in A\) was arbitrary, we conclude that \(A\) is open in \((\R^2,d_\infty)\).
\end{solutionbox}

\newpage
\reflectionpage{Homework 8, Problem 2}

% ---------------------------------------------------------------------------
% HW8 Problem 3
% ---------------------------------------------------------------------------
\newpage
\subsection{Problem 3}

\begin{problemstatementbox}{Homework 8, Problem 3}
Let \(X\) be a nonempty set, and recall that the discrete metric \(d\) is a metric on \(X\). Prove that every set \(A\subseteq X\) is open in the metric space \((X,d)\).
\end{problemstatementbox}

\newpage
\firstdiagnosis{Homework 8, Problem 3}

\newpage
\recipecard{Homework 8, Problem 3}

\newpage
\begin{selfexplanationbox}{Homework 8, Problem 3}
\workquestion{1. What is the definition of the discrete metric?}{0.9in}
\workquestion{2. What is \(B_{1/2}(a)\) in the discrete metric?}{0.95in}
\workquestion{3. Why is \(\{a\}\subseteq A\) whenever \(a\in A\)?}{0.85in}
\workquestion{4. How does this prove every point of \(A\) is interior?}{0.9in}
\workquestion{5. Why does the proof also work when \(A=\varnothing\)?}{0.9in}
\end{selfexplanationbox}

\newpage
\begin{skeletonbox}{Homework 8, Problem 3}
\workquestion{Fix \(A\subseteq X\) and \(a\in A\):}{0.75in}
\workquestion{Compute \(B_{1/2}(a)=\{x\in X:d(x,a)<1/2\}\):}{0.95in}
\workquestion{Use \(B_{1/2}(a)=\{a\}\subseteq A\) to show \(a\) is interior:}{0.9in}
\workquestion{Conclude \(A\) is open because \(a\in A\) was arbitrary:}{0.8in}
\end{skeletonbox}

\newpage
\fullproofpage{Homework 8, Problem 3}{Use this page for the complete proof.}

\newpage
\begin{solutionbox}{Official Solution (Professor)}
Fix \(A\subseteq X\) and \(a\in A\). Notice that
\[
  B_{1/2}(a)=\{x\in X:d(x,a)<1/2\}=\{a\}.
\]
(Indeed, \(d(x,a)\) is either equal to \(0\) or \(1\), and so the only time that \(d(x,a)<1/2\) is when \(d(x,a)=0\).) Then \(a\in B_{1/2}(a)\subseteq A\), and so \(a\) is an interior point for \(A\). As \(a\in A\) was arbitrary, we conclude that \(A\) is open.
\end{solutionbox}

\newpage
\reflectionpage{Homework 8, Problem 3}

% ---------------------------------------------------------------------------
% HW8 Problem 4
% ---------------------------------------------------------------------------
\newpage
\subsection{Problem 4}

\begin{problemstatementbox}{Homework 8, Problem 4}
Let \((X,d)\) be any metric space. Define the function \(\widetilde{d}:X\times X\to\R\) by
\[
  \widetilde{d}(x,y)=\min\{d(x,y),1\}.
\]
\begin{enumerate}[label=(\alph*)]
  \item Prove that \(\widetilde{d}\) is also a metric on \(X\).
  \item Prove that a set \(A\subseteq X\) is open in \((X,d)\) if and only if it is open in \((X,\widetilde{d})\).
\end{enumerate}
\end{problemstatementbox}

\newpage
\firstdiagnosis{Homework 8, Problem 4}

\newpage
\recipecard{Homework 8, Problem 4}

\newpage
\begin{selfexplanationbox}{Homework 8, Problem 4}
\workquestion{1. Which metric axioms are immediate from the definition of \(\widetilde d\)?}{0.95in}
\workquestion{2. Why is definiteness preserved by taking \(\min\{d(x,y),1\}\)?}{0.95in}
\workquestion{3. What are the key cases for the triangle inequality?}{1.0in}
\workquestion{4. Why do \(d\) and \(\widetilde d\) agree on distances less than \(1\)?}{0.95in}
\workquestion{5. Why does shrinking a radius to \(\min\{r,1/2\}\) transfer openness?}{1.0in}
\end{selfexplanationbox}

\newpage
\begin{skeletonbox}{Homework 8, Problem 4}
\textbf{Part (a).}
\workquestion{Check nonnegativity, definiteness, and symmetry:}{1.0in}
\workquestion{For the triangle inequality, split into cases depending on whether \(d(x,y)<1\) and \(d(y,z)<1\):}{1.1in}
\workquestion{Conclude \(\widetilde d\) is a metric:}{0.75in}

\textbf{Part (b).}
\workquestion{Assume \(A\) is open in \((X,\widetilde d)\). Use \(\widetilde d(x,a)\le d(x,a)\) to get a \(d\)-ball inside \(A\):}{1.0in}
\workquestion{Assume \(A\) is open in \((X,d)\). Set \(r_2=\min\{r_1,1/2\}\) and prove the \(\widetilde d\)-ball is inside the \(d\)-ball:}{1.1in}
\end{skeletonbox}

\newpage
\fullproofpage{Homework 8, Problem 4}{Use this page for part (a).}

\newpage
\fullproofpage{Homework 8, Problem 4}{Use this page for part (b).}

\newpage
\begin{solutionbox}{Official Solution (Professor)}
(a) There are \(4\) properties to check.
\begin{enumerate}[label=\arabic*.]
  \item \(\widetilde d(x,y)\ge0\) for all \(x,y\in X\). Note that \(\widetilde d(x,y)\) is equal to \(1\) or \(d(x,y)\). Both of these are nonnegative; we know \(d(x,y)\ge0\) since \(d\) is a metric on \(X\).
  \item \(\widetilde d(x,y)=0\Longleftrightarrow x=y\). If \(x=y\), then \(d(x,y)=0\), so \(\widetilde d(x,y)=\min\{0,1\}=0\). Conversely, if \(\widetilde d(x,y)=0\), then \(\widetilde d(x,y)\ne1\) and so we must have \(d(x,y)=\widetilde d(x,y)=0\). As \(d\) is a metric, this can only mean that \(x=y\).
  \item \(\widetilde d(x,y)=\widetilde d(y,x)\) for all \(x,y\in X\). As \(d\) is a metric, we know that \(d(x,y)=d(y,x)\), and so
  \[
    \widetilde d(x,y)=\min\{d(x,y),1\}=\min\{d(y,x),1\}=\widetilde d(y,x).
  \]
  \item \(\widetilde d(x,z)\le \widetilde d(x,y)+\widetilde d(y,z)\) for all \(x,y,z\in X\). We distinguish \(4\) possible cases:
  \begin{itemize}
    \item \(d(x,y)<1\) and \(d(y,z)<1\): Then
    \[
      \widetilde d(x,z)\le d(x,z)\le d(x,y)+d(y,z)=\widetilde d(x,y)+\widetilde d(y,z).
    \]
    \item \(d(x,y)<1\) and \(d(y,z)\ge1\): Then
    \[
      \widetilde d(x,z)\le1\le d(x,y)+1=\widetilde d(x,y)+\widetilde d(y,z).
    \]
    \item \(d(x,y)\ge1\) and \(d(y,z)<1\): Then
    \[
      \widetilde d(x,z)\le1\le1+d(y,z)=\widetilde d(x,y)+\widetilde d(y,z).
    \]
    \item \(d(x,y)\ge1\) and \(d(y,z)\ge1\): Then
    \[
      \widetilde d(x,z)\le1\le1+1=\widetilde d(x,y)+\widetilde d(y,z).
    \]
  \end{itemize}
\end{enumerate}
In all cases, the triangle inequality holds.

(b) \(\Leftarrow\) Suppose that \(A\) is open in \((X,\widetilde d)\). Fix \(a\in A\). Then there exists \(r>0\) so that
\[
  B_r^{\widetilde d}(a)=\{x\in X:\widetilde d(x,a)<r\}\subseteq A.
\]
Note that \(B_r^d(a)\subseteq B_r^{\widetilde d}(a)\), since
\[
  x\in B_r^d(a)
  \Longrightarrow
  \widetilde d(x,a)\le d(x,a)<r
  \Longrightarrow
  x\in B_r^{\widetilde d}(a).
\]
So \(B_r^d(a)\subseteq B_r^{\widetilde d}(a)\subseteq A\). As \(a\in A\) was arbitrary, we conclude that \(A\) is open in \((X,d)\).

\(\Rightarrow\) Suppose that \(A\) is open in \((X,d)\). Fix \(a\in A\). Then there exists \(r_1>0\) so that \(B_{r_1}^d(a)\subseteq A\). Set \(r_2=\min\{r_1,\frac12\}\). Note that \(B_{r_2}^{\widetilde d}(a)\subseteq B_{r_1}^d(a)\), since
\[
  x\in B_{r_2}^{\widetilde d}(a)
  \Longrightarrow
  \widetilde d(x,a)=\min\{d(x,a),1\}<r_2\le\frac12
\]
\[
  \Longrightarrow
  \widetilde d(x,a)=d(x,a)
  \Longrightarrow
  d(x,a)<r_2\le r_1.
\]
So \(B_{r_2}^{\widetilde d}(a)\subseteq B_{r_1}^d(a)\subseteq A\). As \(a\in A\) was arbitrary, we conclude that \(A\) is open in \((X,\widetilde d)\).
\end{solutionbox}

\newpage
\reflectionpage{Homework 8, Problem 4}

% ---------------------------------------------------------------------------
% HW8 Problem 5
% ---------------------------------------------------------------------------
\newpage
\subsection{Problem 5}

\begin{problemstatementbox}{Homework 8, Problem 5}
Let \((X,d)\) be a metric space, \(A\subseteq X\), and \(x\in X\). Prove that \(x\) is in \(A^\circ\) if and only if there exists an open set \(U\) so that \(x\in U\subseteq A\).
\end{problemstatementbox}

\newpage
\firstdiagnosis{Homework 8, Problem 5}

\newpage
\recipecard{Homework 8, Problem 5}

\newpage
\begin{selfexplanationbox}{Homework 8, Problem 5}
\workquestion{1. What is the definition of \(x\in A^\circ\)?}{0.9in}
\workquestion{2. Why is an open ball an open set?}{0.9in}
\workquestion{3. In the forward direction, what should \(U\) be?}{0.85in}
\workquestion{4. In the reverse direction, why does openness of \(U\) give a ball around \(x\)?}{0.95in}
\workquestion{5. How does \(B_r(x)\subseteq U\subseteq A\) finish the proof?}{0.9in}
\end{selfexplanationbox}

\newpage
\begin{skeletonbox}{Homework 8, Problem 5}
\textbf{Forward direction.}
\workquestion{Assume \(x\in A^\circ\). Choose \(r>0\) with \(B_r(x)\subseteq A\):}{0.85in}
\workquestion{Set \(U=\) and verify open, \(x\in U\), and \(U\subseteq A\):}{0.95in}

\textbf{Reverse direction.}
\workquestion{Assume \(x\in U\subseteq A\) with \(U\) open. Choose \(r>0\) with:}{0.9in}
\workquestion{Use \(B_r(x)\subseteq U\subseteq A\) to conclude \(x\in A^\circ\):}{0.85in}
\end{skeletonbox}

\newpage
\fullproofpage{Homework 8, Problem 5}{Use this page for both directions.}

\newpage
\begin{solutionbox}{Official Solution (Professor)}
\(\Rightarrow\) Suppose that \(x\in A^\circ\). Then there exists \(r>0\) such that \(B_r(x)\subseteq A\). Notice that \(U=B_r(x)\) is an open set that satisfies \(x\in U\subseteq A\).

\(\Leftarrow\) Suppose that \(x\in U\subseteq A\) for some open set \(U\). As \(U\) is open, there exists \(r>0\) so that \(B_r(x)\subseteq U\). Together we have \(B_r(x)\subseteq U\subseteq A\), and so \(x\in A^\circ\).
\end{solutionbox}

\newpage
\reflectionpage{Homework 8, Problem 5}

% ---------------------------------------------------------------------------
% HW8 Problem 6
% ---------------------------------------------------------------------------
\newpage
\subsection{Problem 6}

\begin{problemstatementbox}{Homework 8, Problem 6}
Let \((X,d)\) be a metric space and let \(A\) be a nonempty subset of \(X\). Prove that \(A\) is open if and only if there exists a set \(I\), positive numbers \(r_i>0\), and points \(a_i\in X\) so that
\[
  A=\bigcup_{i\in I}B_{r_i}(a_i).
\]
\end{problemstatementbox}

\newpage
\firstdiagnosis{Homework 8, Problem 6}

\newpage
\recipecard{Homework 8, Problem 6}

\newpage
\begin{selfexplanationbox}{Homework 8, Problem 6}
\workquestion{1. Why should open sets be expressible as unions of balls?}{0.95in}
\workquestion{2. In the forward direction, what index set \(I\) is the natural choice?}{0.95in}
\workquestion{3. Why does every point \(a\in A\) belong to its chosen ball \(B_{r_a}(a)\)?}{0.9in}
\workquestion{4. Why is each ball \(B_{r_i}(a_i)\) open?}{0.9in}
\workquestion{5. Which theorem says arbitrary unions of open sets are open?}{0.95in}
\end{selfexplanationbox}

\newpage
\begin{skeletonbox}{Homework 8, Problem 6}
\textbf{Forward direction.}
\workquestion{Assume \(A\) is open. For each \(a\in A\), choose \(r_a>0\) with \(B_{r_a}(a)\subseteq A\):}{0.95in}
\workquestion{Use \(I=A\), and for each \(i=a\in A\), set the center \(a_i=a\) and radius \(r_i=r_a\):}{0.95in}
\workquestion{Prove both inclusions \(A\subseteq\bigcup_{a\in A}B_{r_a}(a)\) and \(\bigcup_{a\in A}B_{r_a}(a)\subseteq A\):}{1.05in}

\textbf{Reverse direction.}
\workquestion{Assume \(A=\bigcup_{i\in I}B_{r_i}(a_i)\):}{0.75in}
\workquestion{Each ball is open, and arbitrary unions of open sets are open. Conclude:}{0.9in}
\end{skeletonbox}

\newpage
\fullproofpage{Homework 8, Problem 6}{Use this page for the forward direction.}

\newpage
\fullproofpage{Homework 8, Problem 6}{Use this page for the reverse direction.}

\newpage
\begin{solutionbox}{Student-Derived Solution (No official solution available)}
We prove both directions.

Suppose first that \(A\) is open. Since \(A\) is nonempty, for each \(a\in A\) choose \(r_a>0\) such that
\[
  B_{r_a}(a)\subseteq A.
\]
Then
\[
  A=\bigcup_{a\in A}B_{r_a}(a).
\]
Indeed, if \(x\in A\), then \(x\in B_{r_x}(x)\), so \(x\in\bigcup_{a\in A}B_{r_a}(a)\). Conversely, each \(B_{r_a}(a)\subseteq A\), so the union is contained in \(A\). Thus \(A\) has the required form, with index set \(I=A\), and for each index \(i=a\in A\), center \(a_i=a\) and radius \(r_i=r_a\).

Conversely, suppose there exist a set \(I\), positive numbers \(r_i>0\), and points \(a_i\in X\) such that
\[
  A=\bigcup_{i\in I}B_{r_i}(a_i).
\]
Each ball \(B_{r_i}(a_i)\) is open in \((X,d)\), and an arbitrary union of open sets is open. Therefore \(A\) is open.
\end{solutionbox}

\newpage
\reflectionpage{Homework 8, Problem 6}

\section{Homework 9}

% ---------------------------------------------------------------------------
% HW9 Problem 1
% ---------------------------------------------------------------------------
\newpage
\subsection{Problem 1}

\begin{problemstatementbox}{Homework 9, Problem 1}
Consider the metric space \((\R,d)\) with the usual metric \(d(x,y)=|x-y|\). For each of the following sets \(A\subseteq\R\), find its closure, isolated points, and limit points. Prove your answer.

\begin{enumerate}[label=(\alph*)]
  \item \(A=[0,1)\)
  \item \(A=\Z\)
  \item \(A=\{x\in\Q:0<x<1\}\)
\end{enumerate}
\end{problemstatementbox}

\newpage
\firstdiagnosis{Homework 9, Problem 1}

\newpage
\recipecard{Homework 9, Problem 1}

\newpage
\begin{selfexplanationbox}{Homework 9, Problem 1}
\workquestion{1. For each set, which points are obvious members of the closure, and which boundary points require a ball argument?}{0.95in}
\workquestion{2. How do you prove that a point is not in \(\closure{\Z}\)? What radius separates it from all integers?}{0.95in}
\workquestion{3. What is the difference between proving \(a\in A\) is isolated and proving \(a\in A\) is a limit point?}{0.95in}
\workquestion{4. Where is density of \(\Q\) used in part (c), and why is it not enough just to say ``there are rationals nearby''?}{1.05in}
\workquestion{5. How does the identity \(\closure{A}=(\text{isolated points})\cup A'\) finish the limit-point computations?}{0.95in}
\end{selfexplanationbox}

\newpage
\begin{skeletonbox}{Homework 9, Problem 1}
\textbf{Part (a): \(A=[0,1)\)}
\workquestion{Closure: show \([0,1)\subseteq\closure{A}\subseteq[0,1]\), then prove \(1\in\closure{A}\) by choosing a point of \(A\) in \((1-r,1+r)\):}{1.0in}
\workquestion{Isolated points: fix \(a\in A\). Given \(r>0\), choose a second point of \(A\) inside \((a-r,a+r)\):}{0.85in}
\workquestion{Limit points: use the closure/isolated-point relationship to conclude:}{0.55in}

\textbf{Part (b): \(A=\Z\)}
\workquestion{Closure: for \(x\notin\Z\), choose \(n\in\Z\) with \(n<x<n+1\) and radius \(r=\)}{0.75in}
\workquestion{Isolated points and limit points: choose the ball around \(n\in\Z\):}{0.75in}
\end{skeletonbox}

\newpage
\begin{skeletonbox}{Homework 9, Problem 1 (continued)}
\textbf{Part (c): \(A=\{x\in\Q:0<x<1\}\)}
\workquestion{Closure: prove \(0\), points in \((0,1]\), and no points outside \([0,1]\) behave correctly:}{0.95in}
\workquestion{Isolated points and limit points: use density of \(\Q\) to find a second point:}{0.75in}
\end{skeletonbox}

\newpage
\fullproofpage{Homework 9, Problem 1}{Use this page for parts (a) and (b).}

\newpage
\fullproofpage{Homework 9, Problem 1}{Use this page for part (c) and final cleanup.}

\newpage
\begin{solutionbox}{Official Solution (Professor)}
\begin{enumerate}[label=(\alph*)]
  \item Closure: \(\closure{A}=[0,1]\). Since \(\closure{A}\) is the smallest closed set containing \(A\), and \([0,1]\) is one such closed set, we must have
  \[
    [0,1)\subseteq\closure{A}\subseteq[0,1].
  \]
  Also \(1\in\closure{A}\): given \(r>0\),
  \[
    (1-r,1+r)\cap[0,1)\neq\varnothing,
  \]
  since \(\max\{1-r/2,0\}\in(1-r,1+r)\cap[0,1)\).

  Isolated points: \(\varnothing\). Fix \(a\in A\). Given \(r>0\), the point
  \[
    \min\left\{a+\frac r2,\frac{a+1}{2}\right\}
  \]
  is in \((a-r,a+r)\cap A\) and is not \(a\).

  Limit points: \(A'=[0,1]\). By definition, \(\closure{A}\) is the union of isolated points and limit points, and there are no isolated points.

  \item Closure: \(\closure{A}=\Z\). We automatically have \(\Z\subseteq\closure{A}\). Fix \(x\in\Z^c\). Let \(n\in\Z\) satisfy \(n<x<n+1\). Then
  \[
    r=\min\{x-n,n+1-x\}
  \]
  is positive and \((x-r,x+r)\cap\Z=\varnothing\). Thus \(x\notin\closure{A}\).

  Isolated points: \(\Z\). For \(n\in\Z\),
  \[
    (n-1,n+1)\cap\Z=\{n\}.
  \]

  Limit points: \(A'=\varnothing\). Every point of \(A\) is isolated.

  \item Closure: \(\closure{A}=[0,1]\). First, \(A\subseteq[0,1]\), and \([0,1]\) is closed, so \(\closure{A}\subseteq[0,1]\). The point \(0\) is adherent to \(A\): given \(r>0\), choose \(n>\max\{1/r,1\}\), so \(1/n\in(-r,r)\cap A\). Now fix \(x\in(0,1]\). By density of \(\Q\) in \(\R\), for every \(r>0\) there is \(q\in\Q\) between \(\max\{x-r,0\}\) and \(x\), and this \(q\in(x-r,x+r)\cap A\).

  Isolated points: \(\varnothing\). Given \(x\in A\) and \(r>0\), choose \(q\in\Q\) as above. This \(q\) lies in \((x-r,x+r)\cap A\) and \(q\neq x\).

  Limit points: \(A'=[0,1]\). There are no isolated points.
\end{enumerate}
\end{solutionbox}

\newpage
\reflectionpage{Homework 9, Problem 1}

% ---------------------------------------------------------------------------
% HW9 Problem 2
% ---------------------------------------------------------------------------
\newpage
\subsection{Problem 2}

\begin{problemstatementbox}{Homework 9, Problem 2}
Let \((X,d)\) be a metric space, let \(A\) be a nonempty subset of \(X\), and let \(U\) be an open subset of \(X\).

\begin{enumerate}[label=(\alph*)]
  \item Prove that \(U\cap\closure{A}\subseteq\closure{U\cap A}\).
  \item Prove that \(\closure{U\cap\closure{A}}=\closure{U\cap A}\). (Hint: Use part (a) and properties of closure from lecture.)
  \item Conclude that if \(U\cap A=\varnothing\), then \(U\cap\closure{A}=\varnothing\).
\end{enumerate}
\end{problemstatementbox}

\newpage
\firstdiagnosis{Homework 9, Problem 2}

\newpage
\recipecard{Homework 9, Problem 2}

\newpage
\begin{selfexplanationbox}{Homework 9, Problem 2}
\workquestion{1. Why is the assumption that \(U\) is open essential in part (a)?}{1.05in}
\workquestion{2. After fixing \(x\in U\cap\closure{A}\) and \(r>0\), why do we introduce a second radius \(s>0\)?}{1.0in}
\workquestion{3. Why is \(\widetilde r=\min\{r,s\}\) the right radius to use?}{0.95in}
\workquestion{4. Which closure property turns part (a) into the first inclusion in part (b)?}{0.95in}
\workquestion{5. In part (c), why does \(\closure{U\cap\closure{A}}=\varnothing\) force \(U\cap\closure{A}=\varnothing\)?}{1.0in}
\end{selfexplanationbox}

\newpage
\begin{skeletonbox}{Homework 9, Problem 2}
\textbf{Part (a): prove \(U\cap\closure{A}\subseteq\closure{U\cap A}\).}
\workquestion{Fix \(x\in\underline{\hspace{1.6in}}\). To show \(x\in\closure{U\cap A}\), fix \(r>0\) and aim to find a point in:}{0.85in}
\workquestion{Use openness of \(U\): since \(x\in U\), choose \(s>0\) such that:}{0.75in}
\workquestion{Use \(x\in\closure{A}\) with \(\widetilde r=\min\{r,s\}\) to choose \(y\in\):}{0.75in}
\workquestion{Check the three memberships \(y\in A\), \(y\in U\), and \(y\in B_r(x)\):}{0.95in}
\end{skeletonbox}

\newpage
\begin{skeletonbox}{Homework 9, Problem 2 (continued)}
\textbf{Part (b): prove equality of closures.}
\workquestion{First inclusion from part (a) and monotonicity of closure:}{0.75in}
\workquestion{Second inclusion from \(A\subseteq\closure{A}\):}{0.75in}

\textbf{Part (c): empty-set conclusion.}
\workquestion{Start from \(U\cap A=\varnothing\), take closures, use part (b), and then use \(B\subseteq\closure{B}\):}{0.85in}
\end{skeletonbox}

\newpage
\fullproofpage{Homework 9, Problem 2}{Use this page for parts (a) and (b).}

\newpage
\fullproofpage{Homework 9, Problem 2}{Use this page for part (c) and final cleanup.}

\newpage
\begin{solutionbox}{Official Solution (Professor)}
\begin{enumerate}[label=(\alph*)]
  \item Fix \(x\in U\cap\closure{A}\). We show \(x\in\closure{U\cap A}\), i.e., for any \(r>0\), \(\ball{r}{x}\cap U\cap A\) is nonempty.

  Fix \(r>0\). Since \(x\in U\) and \(U\) is open, there exists \(s>0\) such that \(\ball{s}{x}\subseteq U\). Set \(\widetilde r=\min\{r,s\}\). Since \(x\in\closure{A}\), the set \(\ball{\widetilde r}{x}\cap A\) is nonempty, so choose \(y\in\ball{\widetilde r}{x}\cap A\).

  Then \(y\in A\). Also \(y\in\ball{r}{x}\), since \(d(y,x)<\widetilde r\le r\). Finally \(y\in U\), since \(d(y,x)<\widetilde r\le s\), so \(y\in\ball{s}{x}\subseteq U\). Hence \(y\in\ball{r}{x}\cap U\cap A\), as desired.

  \item Recall that \(B\subseteq C\) implies \(\closure{B}\subseteq\closure{C}\). Applying this to part (a),
  \[
    \closure{U\cap\closure{A}}\subseteq\closure{U\cap A}.
  \]
  On the other hand, \(A\subseteq\closure{A}\), so \(U\cap A\subseteq U\cap\closure{A}\). Therefore
  \[
    \closure{U\cap A}\subseteq\closure{U\cap\closure{A}}.
  \]
  Together,
  \[
    \closure{U\cap\closure{A}}=\closure{U\cap A}.
  \]

  \item For any set \(B\), \(B\subseteq\closure{B}\). If \(U\cap A=\varnothing\), then \(\closure{U\cap A}=\varnothing\). By part (b),
  \[
    \closure{U\cap\closure{A}}=\varnothing.
  \]
  Therefore \(U\cap\closure{A}=\varnothing\).
\end{enumerate}
\end{solutionbox}

\newpage
\reflectionpage{Homework 9, Problem 2}

% ---------------------------------------------------------------------------
% HW9 Problem 3
% ---------------------------------------------------------------------------
\newpage
\subsection{Problem 3}

\begin{problemstatementbox}{Homework 9, Problem 3}
Recall that \(d_2\) is a metric on \(\R^2\). Consider the subspace
\[
  Y=\{(x,0):x\in\R\}
\]
with the induced metric \(d_Y\).

\begin{enumerate}[label=(\alph*)]
  \item Provide an example of a set \(A\subseteq Y\) where \(A\) is open in \((Y,d_Y)\) but not open in \((\R^2,d_2)\).
  \item Prove that for any \(B\subseteq Y\), if \(B\) is closed in \((Y,d_Y)\) then \(B\) is closed in \((\R^2,d_2)\).
\end{enumerate}
\end{problemstatementbox}

\newpage
\firstdiagnosis{Homework 9, Problem 3}

\newpage
\recipecard{Homework 9, Problem 3}

\newpage
\begin{selfexplanationbox}{Homework 9, Problem 3}
\workquestion{1. What does a ball in the subspace \(Y\) look like compared with a ball in \(\R^2\)?}{1.05in}
\workquestion{2. Why is the line segment \((-1,1)\times\{0\}\) open in \(Y\)?}{1.0in}
\workquestion{3. Why does the point \((0,r/2)\) prove that this set is not open in \(\R^2\)?}{0.95in}
\workquestion{4. What lecture result lets us reduce part (b) to proving that \(Y\) is closed in \(\R^2\)?}{0.95in}
\workquestion{5. How does the vertical coordinate \(x_2\neq0\) give a ball contained in \(Y^c\)?}{1.0in}
\end{selfexplanationbox}

\newpage
\begin{skeletonbox}{Homework 9, Problem 3}
\textbf{Part (a): example open in \(Y\), not open in \(\R^2\).}
\workquestion{Choose \(A=\)}{0.65in}
\workquestion{To prove \(A\) is open in \(Y\), fix \(x=(x_1,0)\in A\) and choose \(r=\)}{0.8in}
\workquestion{If \(y\in B_r^Y(x)\), show \(|y_1|<1\), hence \(y\in A\):}{0.95in}
\workquestion{To prove \(A\) is not open in \(\R^2\), fix \(r>0\) around \((0,0)\) and find an off-axis point:}{0.85in}

\textbf{Part (b): closedness transfer.}
\workquestion{State the lecture corollary being used:}{0.85in}
\end{skeletonbox}

\newpage
\begin{skeletonbox}{Homework 9, Problem 3 (continued)}
\textbf{Part (b): closedness transfer, continued.}
\workquestion{Prove \(Y\) is closed by fixing \(x=(x_1,x_2)\in Y^c\), so \(x_2\neq0\), and choosing radius \(r=\)}{0.9in}
\workquestion{Show \(B_r^{\R^2}(x)\subseteq Y^c\) using the inequality involving \(|x_2-y_2|\):}{1.0in}
\end{skeletonbox}

\newpage
\fullproofpage{Homework 9, Problem 3}{Use this page for part (a).}

\newpage
\fullproofpage{Homework 9, Problem 3}{Use this page for part (b).}

\newpage
\begin{solutionbox}{Official Solution (Professor)}
\begin{enumerate}[label=(\alph*)]
  \item Consider
  \[
    A=\{(x_1,0):-1<x_1<1\}.
  \]
  First, \(A\) is open in \((Y,d_Y)\). Fix \(x=(x_1,0)\in A\), and set \(r=1-|x_1|\). Then \(r>0\), and
  \[
    B_r^Y(x)=\{y\in Y:d_Y(x,y)<r\}\subseteq A,
  \]
  since
  \[
    y\in B_r^Y(x)
    \Longrightarrow
    r>d_Y(x,y)=\sqrt{|x_1-y_1|^2+|0-0|^2}=|x_1-y_1|
  \]
  and therefore
  \[
    |y_1|\le |x_1|+|y_1-x_1|<|x_1|+r=1.
  \]

  Next, \(A\) is not open in \((\R^2,d_2)\), because \((0,0)\in A\) is not an interior point. For any \(r>0\), the point \((0,r/2)\) lies in \(B_r^{\R^2}((0,0))\) but not in \(A\).

  \item By the corollary from Lecture 24, it suffices to show that \(Y\) is closed in \((\R^2,d_2)\). Fix \(x=(x_1,x_2)\in Y^c\). Then \(x_2\neq0\). Set \(r=|x_2|\). If \(y\in B_r^{\R^2}(x)\), then
  \[
    r>d_2(x,y)=\sqrt{|x_1-y_1|^2+|x_2-y_2|^2}\ge |x_2-y_2|.
  \]
  Hence
  \[
    |y_2|\ge |x_2|-|x_2-y_2|>|x_2|-r=0.
  \]
  Thus \(y_2\neq0\), so \(y\in Y^c\). Therefore \(B_r^{\R^2}(x)\subseteq Y^c\). Since every point of \(Y^c\) is interior to \(Y^c\), \(Y^c\) is open, and \(Y\) is closed.
\end{enumerate}
\end{solutionbox}

\newpage
\reflectionpage{Homework 9, Problem 3}

% ---------------------------------------------------------------------------
% HW9 Problem 4
% ---------------------------------------------------------------------------
\newpage
\subsection{Problem 4}

\begin{problemstatementbox}{Homework 9, Problem 4}
Let \((X,d)\) be a metric space, \(\{x_n\}_{n\geq1}\) a sequence in \(X\), and \(x\in X\). Prove that the sequence \(\{x_n\}_{n\geq1}\) converges to \(x\) if and only if for every open set \(U\) that contains \(x\), we have \(x_n\in U\) for all but finitely many \(n\in\N\).
\end{problemstatementbox}

\newpage
\firstdiagnosis{Homework 9, Problem 4}

\newpage
\recipecard{Homework 9, Problem 4}

\newpage
\begin{selfexplanationbox}{Homework 9, Problem 4}
\workquestion{1. In the forward direction, how does openness of \(U\) connect an arbitrary neighborhood to a ball around \(x\)?}{1.05in}
\workquestion{2. What does ``for all but finitely many \(n\)'' mean in terms of choosing an \(N\)?}{1.0in}
\workquestion{3. In the reverse direction, why is \(B_r(x)\) the correct open set to test?}{1.0in}
\workquestion{4. Why is it enough to prove \(x_n\in B_r(x)\) eventually for each \(r>0\)?}{1.0in}
\workquestion{5. What is the main logical difference between the two directions?}{0.95in}
\end{selfexplanationbox}

\newpage
\begin{skeletonbox}{Homework 9, Problem 4}
\textbf{Forward direction: assume \(x_n\to x\).}
\workquestion{Fix an open set \(U\subseteq X\) with \(x\in U\). By openness, choose \(r>0\) such that:}{0.9in}
\workquestion{By convergence, choose \(N\) such that for all \(n\ge N\):}{0.85in}
\workquestion{Conclude \(x_n\in U\) for all but finitely many \(n\):}{0.75in}

\textbf{Reverse direction: assume eventual membership in every open \(U\ni x\).}
\workquestion{Fix \(r>0\). Apply the hypothesis to the open set:}{0.8in}
\workquestion{Convert ``all but finitely many'' into an \(N\) so that for all \(n\ge N\):}{0.9in}
\workquestion{Translate \(x_n\in B_r(x)\) into the metric convergence inequality:}{0.85in}
\workquestion{Conclude because \(r>0\) was arbitrary:}{0.65in}
\end{skeletonbox}

\newpage
\fullproofpage{Homework 9, Problem 4}{Use this page for both directions.}

\newpage
\begin{solutionbox}{Official Solution (Professor)}
\(\Rightarrow\) Suppose \(x_n\to x\) in \((X,d)\) as \(n\to\infty\). Fix an open set \(U\subseteq X\) such that \(x\in U\). Since \(U\) is open, there exists \(r>0\) such that \(\ball{r}{x}\subseteq U\). By definition of convergence, there exists \(N\) such that \(d(x_n,x)<r\) for all \(n\ge N\). Therefore \(x_n\in\ball{r}{x}\subseteq U\) for all \(n\ge N\).

\(\Leftarrow\) Suppose that for every open set \(U\) containing \(x\), we have \(x_n\in U\) for all but finitely many \(n\in\N\). Fix \(r>0\). Apply the assumption to the open set \(\ball{r}{x}\). Then \(x_n\in\ball{r}{x}\) for all but finitely many \(n\in\N\). Let \(N\) be larger than the largest index for which \(x_n\notin\ball{r}{x}\). Then for every \(n\ge N\), \(x_n\in\ball{r}{x}\), and hence \(d(x_n,x)<r\). Since \(r>0\) was arbitrary, \(\{x_n\}_{n\ge1}\) converges to \(x\) in \((X,d)\).
\end{solutionbox}

\newpage
\reflectionpage{Homework 9, Problem 4}

% ---------------------------------------------------------------------------
% HW9 Problem 5
% ---------------------------------------------------------------------------
\newpage
\subsection{Problem 5}

\begin{problemstatementbox}{Homework 9, Problem 5}
Recall that \(d_2\) is a metric on \(\R^2\). Suppose that \(x^n=(x_1^n,x_2^n)\) is a sequence of points in \(\R^2\).

\begin{enumerate}[label=(\alph*)]
  \item Prove that \(|x_1^n-x_1|\leq d_2(x^n,x)\) for any \(x=(x_1,x_2)\in\R^2\).
  \item Prove that the sequence \(\{x^n\}_{n\geq1}\) converges in \((\R^2,d_2)\) if and only if the sequences \(\{x_1^n\}_{n\geq1}\) and \(\{x_2^n\}_{n\geq1}\) of real numbers both converge.
\end{enumerate}
\end{problemstatementbox}

\newpage
\firstdiagnosis{Homework 9, Problem 5}

\newpage
\recipecard{Homework 9, Problem 5}

\newpage
\begin{selfexplanationbox}{Homework 9, Problem 5}
\workquestion{1. Why does the Euclidean distance \(d_2(x^n,x)\) control each coordinate difference?}{1.05in}
\workquestion{2. In the forward direction, what is squeezed between \(0\) and \(d_2(x^n,x)\)?}{0.95in}
\workquestion{3. Why do the coordinate sequences converge to the coordinates of the same point \(x\)?}{0.95in}
\workquestion{4. In the reverse direction, how do the real limit laws imply \(d_2(x^n,x)\to0\)?}{1.05in}
\workquestion{5. What would change if the problem were in \(\R^3\) or \(\R^m\)?}{0.95in}
\end{selfexplanationbox}

\newpage
\begin{skeletonbox}{Homework 9, Problem 5}
\textbf{Part (a): coordinate estimate.}
\workquestion{Write the formula for \(d_2(x^n,x)\):}{0.8in}
\workquestion{Use \(|x_2^n-x_2|^2\ge0\) to get the desired inequality:}{0.9in}

\textbf{Part (b), forward direction.}
\workquestion{Assume \(x^n\to x=(x_1,x_2)\) in \((\R^2,d_2)\). Then \(d_2(x^n,x)\to\)}{0.7in}
\workquestion{Use part (a) and squeeze to prove \(x_1^n\to x_1\):}{0.85in}
\workquestion{Repeat the same argument for the second coordinate:}{0.75in}
\end{skeletonbox}

\newpage
\begin{skeletonbox}{Homework 9, Problem 5 (continued)}
\textbf{Part (b), reverse direction.}
\workquestion{Assume \(x_1^n\to x_1\) and \(x_2^n\to x_2\). Set \(x=\)}{0.65in}
\workquestion{Use the formula for \(d_2(x^n,x)\) and the limit laws to conclude:}{0.95in}
\end{skeletonbox}

\newpage
\fullproofpage{Homework 9, Problem 5}{Use this page for part (a) and the forward direction of part (b).}

\newpage
\fullproofpage{Homework 9, Problem 5}{Use this page for the reverse direction of part (b).}

\newpage
\begin{solutionbox}{Official Solution (Professor)}
\begin{enumerate}[label=(\alph*)]
  \item Recall that
  \[
    d_2(x^n,x)=\left(|x_1^n-x_1|^2+|x_2^n-x_2|^2\right)^{1/2}.
  \]
  Since \(|x_2^n-x_2|^2\ge0\),
  \[
    d_2(x^n,x)\ge \left(|x_1^n-x_1|^2+0\right)^{1/2}\ge |x_1^n-x_1|.
  \]

  \item \(\Rightarrow\) Suppose \(\{x^n\}_{n\ge1}\) converges to \(x\) in \((\R^2,d_2)\). Then
  \[
    d_2(x^n,x)\to0.
  \]
  By part (a),
  \[
    0\le |x_1^n-x_1|\le d_2(x^n,x).
  \]
  By the squeeze theorem, \(|x_1^n-x_1|\to0\), so \(\{x_1^n\}_{n\ge1}\) converges to \(x_1\). Similarly,
  \[
    0\le |x_2^n-x_2|\le d_2(x^n,x),
  \]
  so \(\{x_2^n\}_{n\ge1}\) converges to \(x_2\).

  \(\Leftarrow\) Suppose \(\{x_1^n\}_{n\ge1}\) converges to \(x_1\) and \(\{x_2^n\}_{n\ge1}\) converges to \(x_2\). Then
  \[
    |x_1^n-x_1|\to0
    \qquad\text{and}\qquad
    |x_2^n-x_2|\to0.
  \]
  By the limit laws,
  \[
    d_2(x^n,x)
    =
    \left(|x_1^n-x_1|^2+|x_2^n-x_2|^2\right)^{1/2}
    \to0.
  \]
  Thus \(\{x^n\}_{n\ge1}\) converges to \(x\) in \((\R^2,d_2)\).
\end{enumerate}
\end{solutionbox}

\newpage
\reflectionpage{Homework 9, Problem 5}

% ---------------------------------------------------------------------------
% HW9 Problem 6
% ---------------------------------------------------------------------------
\newpage
\subsection{Problem 6}

\begin{problemstatementbox}{Homework 9, Problem 6}
Let \((X,d)\) be a metric space and \(\{x_n\}_{n\geq1}\) a sequence in \(X\). Prove that if \(\{x_n\}_{n\geq1}\) converges to \(x\in X\), then for any \(y\in X\), the sequence \(\{d(x_n,y)\}_{n\geq1}\) of real numbers converges to \(d(x,y)\).
\end{problemstatementbox}

\newpage
\firstdiagnosis{Homework 9, Problem 6}

\newpage
\recipecard{Homework 9, Problem 6}

\newpage
\begin{selfexplanationbox}{Homework 9, Problem 6}
\workquestion{1. What real-valued sequence are we trying to prove converges?}{0.95in}
\workquestion{2. Why should the triangle inequality compare \(d(x_n,y)\) and \(d(x,y)\)?}{1.05in}
\workquestion{3. How do two applications of the triangle inequality produce an absolute value estimate?}{1.05in}
\workquestion{4. What quantity goes to \(0\), and why does it squeeze the distance error?}{1.0in}
\workquestion{5. How is this problem a sequence version of continuity of the map \(z\mapsto d(z,y)\)?}{1.0in}
\end{selfexplanationbox}

\newpage
\begin{skeletonbox}{Homework 9, Problem 6}
\textbf{Setup.}
\workquestion{Fix \(y\in X\). Since \(x_n\to x\), we know \(d(x_n,x)\to\)}{0.75in}

\textbf{Triangle inequality estimates.}
\workquestion{First apply triangle inequality to bound \(d(x_n,y)\) in terms of \(d(x_n,x)\) and \(d(x,y)\):}{0.9in}
\workquestion{Rearrange to get \(d(x_n,y)-d(x,y)\le\)}{0.7in}
\workquestion{Apply triangle inequality again to bound \(d(x,y)\) in terms of \(d(x,x_n)\) and \(d(x_n,y)\):}{0.9in}
\workquestion{Rearrange to get \(d(x,y)-d(x_n,y)\le\)}{0.7in}
\end{skeletonbox}

\newpage
\begin{skeletonbox}{Homework 9, Problem 6 (continued)}
\textbf{Finish.}
\workquestion{Combine the two inequalities into an absolute value estimate:}{0.8in}
\workquestion{Use squeeze theorem to conclude:}{0.75in}
\end{skeletonbox}

\newpage
\fullproofpage{Homework 9, Problem 6}{Use this page for the complete proof.}

\newpage
\begin{solutionbox}{Official Solution (Professor)}
Fix \(y\in X\). By the triangle inequality,
\[
\begin{aligned}
  d(x_n,y)&\le d(x_n,x)+d(x,y)
  &&\Longrightarrow&
  d(x_n,y)-d(x,y)&\le d(x_n,x),\\
  d(x,y)&\le d(x,x_n)+d(x_n,y)
  &&\Longrightarrow&
  d(x,y)-d(x_n,y)&\le d(x_n,x).
\end{aligned}
\]
Together,
\[
  |d(x_n,y)-d(x,y)|\le d(x_n,x).
\]
Since \(\{x_n\}_{n\ge1}\) converges to \(x\) in \((X,d)\), the right-hand side converges to \(0\). Therefore, by the squeeze theorem,
\[
  |d(x_n,y)-d(x,y)|\to0,
\]
and thus \(d(x_n,y)\) converges to \(d(x,y)\).
\end{solutionbox}

\newpage
\reflectionpage{Homework 9, Problem 6}

\newpage
\section{Homework 10}

% ---------------------------------------------------------------------------
% HW10 Problem 1
% ---------------------------------------------------------------------------
\newpage
\subsection{Problem 1}

\begin{problemstatementbox}{Homework 10, Problem 1}
Let \((X,d)\) be a metric space and let \(Y\subseteq X\). We call the set \(Y\) complete if the subspace \(Y\) with the induced metric is complete. Prove that if \(Y\) is complete, then \(Y\) is closed in \(X\). (Hint: It suffices to show \(Y\subseteq\closure{Y}\). Use sequences.)
\end{problemstatementbox}

\newpage
\firstdiagnosis{Homework 10, Problem 1}

\newpage
\recipecard{Homework 10, Problem 1}

\newpage
\begin{selfexplanationbox}{Homework 10, Problem 1}
\workquestion{1. What does it mean for \(Y\) to be complete as a subspace of \(X\)?}{1.0in}
\workquestion{2. What sequential characterization of closure should be used after fixing \(x\in\closure{Y}\)?}{1.0in}
\workquestion{3. Why is a sequence converging in \(X\) automatically Cauchy in \(X\)?}{0.95in}
\workquestion{4. Why is the same sequence Cauchy in the induced metric on \(Y\)?}{0.95in}
\workquestion{5. Where is uniqueness of limits in a metric space used?}{1.05in}
\end{selfexplanationbox}

\newpage
\begin{skeletonbox}{Homework 10, Problem 1}
\textbf{Goal: prove \(Y\) is closed.}
\workquestion{It suffices to show \(\closure{Y}\subseteq Y\). Fix:}{0.85in}
\workquestion{Use the sequential characterization of closure to get a sequence \(\{y_n\}\subseteq Y\) such that:}{0.95in}
\workquestion{Explain why \(\{y_n\}\) is Cauchy in \(Y\):}{0.95in}
\workquestion{Use completeness of \(Y\) to get \(y\in Y\) with:}{0.85in}
\workquestion{Compare convergence in \(Y\) with convergence in \(X\), then use uniqueness of limits:}{1.0in}
\workquestion{Conclude \(\closure{Y}\subseteq Y\), hence \(Y\) is closed:}{0.8in}
\end{skeletonbox}

\newpage
\fullproofpage{Homework 10, Problem 1}{Use this page for the complete proof.}

\newpage
\begin{solutionbox}{Student-Derived Solution (No official solution available)}
We show that \(\closure{Y}\subseteq Y\). Since always \(Y\subseteq\closure{Y}\), this implies \(\closure{Y}=Y\), and hence \(Y\) is closed in \(X\).

Let \(x\in\closure{Y}\). By the sequential characterization of closure, there exists a sequence \(\{y_n\}_{n\ge1}\) in \(Y\) such that
\[
  y_n\to x
  \qquad\text{in }(X,d).
\]
Every convergent sequence in a metric space is Cauchy, so \(\{y_n\}\) is Cauchy in \((X,d)\). Since the induced metric on \(Y\) is the restriction of \(d\), the same sequence is Cauchy in \((Y,d_Y)\).

Because \(Y\) is complete, there is some \(y\in Y\) such that
\[
  y_n\to y
  \qquad\text{in }(Y,d_Y).
\]
Convergence in the subspace metric is the same distance condition as convergence in \(X\), so \(y_n\to y\) in \((X,d)\). Thus the sequence \(\{y_n\}\) converges to both \(x\) and \(y\) in \(X\). By uniqueness of limits in metric spaces, \(x=y\). Since \(y\in Y\), we have \(x\in Y\).

Therefore \(\closure{Y}\subseteq Y\), so \(Y\) is closed in \(X\).
\end{solutionbox}

\newpage
\reflectionpage{Homework 10, Problem 1}

% ---------------------------------------------------------------------------
% HW10 Problem 2
% ---------------------------------------------------------------------------
\newpage
\subsection{Problem 2}

\begin{problemstatementbox}{Homework 10, Problem 2}
Let \((X,d)\) be a metric space. Prove that if \(A,B\subseteq X\) are both compact then \(A\cup B\) is compact.
\end{problemstatementbox}

\newpage
\firstdiagnosis{Homework 10, Problem 2}

\newpage
\recipecard{Homework 10, Problem 2}

\newpage
\begin{selfexplanationbox}{Homework 10, Problem 2}
\workquestion{1. What definition of compactness is most direct for a finite union?}{0.95in}
\workquestion{2. Why does an open cover of \(A\cup B\) automatically cover \(A\) and \(B\) separately?}{1.0in}
\workquestion{3. What finite subcover do you get from compactness of \(A\)?}{0.95in}
\workquestion{4. What finite subcover do you get from compactness of \(B\)?}{0.95in}
\workquestion{5. Why is the union of the two finite index sets still finite and enough to cover \(A\cup B\)?}{1.05in}
\end{selfexplanationbox}

\newpage
\begin{skeletonbox}{Homework 10, Problem 2}
\textbf{Open-cover setup.}
\workquestion{Let \(\{U_i\}_{i\in I}\) be an arbitrary open cover of \(A\cup B\):}{0.85in}
\workquestion{Restrict the same cover to \(A\), and use compactness of \(A\) to choose a finite \(J_A\subseteq I\):}{1.0in}
\workquestion{Restrict the same cover to \(B\), and use compactness of \(B\) to choose a finite \(J_B\subseteq I\):}{1.0in}
\workquestion{Set \(J=\)}{0.6in}
\workquestion{Check \(J\) is finite and \(\{U_i\}_{i\in J}\) covers \(A\cup B\):}{1.0in}
\workquestion{Conclude compactness because the original open cover was arbitrary:}{0.75in}
\end{skeletonbox}

\newpage
\fullproofpage{Homework 10, Problem 2}{Use this page for the complete proof.}

\newpage
\begin{solutionbox}{Student-Derived Solution (No official solution available)}
Let \(\{U_i\}_{i\in I}\) be an open cover of \(A\cup B\). Thus each \(U_i\) is open in \(X\), and
\[
  A\cup B\subseteq \bigcup_{i\in I}U_i.
\]
Since \(A\subseteq A\cup B\), the same family is an open cover of \(A\). Because \(A\) is compact, there is a finite set \(J_A\subseteq I\) such that
\[
  A\subseteq \bigcup_{i\in J_A}U_i.
\]
Similarly, \(\{U_i\}_{i\in I}\) is an open cover of \(B\), so compactness of \(B\) gives a finite set \(J_B\subseteq I\) such that
\[
  B\subseteq \bigcup_{i\in J_B}U_i.
\]
Set \(J=J_A\cup J_B\). Then \(J\) is finite, and
\[
  A\cup B
  \subseteq
  \left(\bigcup_{i\in J_A}U_i\right)\cup
  \left(\bigcup_{i\in J_B}U_i\right)
  =
  \bigcup_{i\in J}U_i.
\]
Thus every open cover of \(A\cup B\) has a finite subcover. Therefore \(A\cup B\) is compact.
\end{solutionbox}

\newpage
\reflectionpage{Homework 10, Problem 2}

% ---------------------------------------------------------------------------
% HW10 Problem 3
% ---------------------------------------------------------------------------
\newpage
\subsection{Problem 3}

\begin{problemstatementbox}{Homework 10, Problem 3}
Let \((X,d)\) be a metric space.

\begin{enumerate}[label=(\alph*)]
  \item Prove that if \(A\subseteq B\subseteq X\), \(A\) is closed in \(X\), and \(B\) is compact, then \(A\) is compact. (Hint: Given an open cover of \(A\), we can add \(A^c\) to get an open cover of \(B\).)
  \item Prove that if \(C,D\subseteq X\) are both compact, then \(C\cap D\) is compact. (Hint: Use (a).)
\end{enumerate}
\end{problemstatementbox}

\newpage
\firstdiagnosis{Homework 10, Problem 3}

\newpage
\recipecard{Homework 10, Problem 3}

\newpage
\begin{selfexplanationbox}{Homework 10, Problem 3}
\workquestion{1. Why does adding \(A^c\) to an open cover of \(A\) help cover all of \(B\)?}{1.05in}
\workquestion{2. Where is the assumption that \(A\) is closed used?}{0.95in}
\workquestion{3. After compactness of \(B\) gives a finite subcover, why can \(A^c\) be discarded when covering \(A\)?}{1.05in}
\workquestion{4. Which lecture result says compact subsets of a metric space are closed?}{0.9in}
\workquestion{5. In part (b), which compact set plays the role of \(B\) in part (a)?}{0.95in}
\end{selfexplanationbox}

\newpage
\begin{skeletonbox}{Homework 10, Problem 3}
\textbf{Part (a): closed subset of a compact set.}
\workquestion{Let \(\{U_i\}_{i\in I}\) be an open cover of \(A\). Since \(A\) is closed, \(A^c\) is:}{0.85in}
\workquestion{Show \(\{A^c\}\cup\{U_i\}_{i\in I}\) covers \(B\):}{1.0in}
\workquestion{Use compactness of \(B\) to extract a finite subcover:}{0.95in}
\workquestion{Restrict back to \(A\) and remove the \(A^c\) part:}{1.0in}

\textbf{Part (b): compact intersection.}
\workquestion{Use compactness \(\Rightarrow\) closed to show \(C\cap D\) is closed in \(X\):}{0.9in}
\workquestion{Apply part (a) with \(A=C\cap D\) and \(B=\)}{0.75in}
\end{skeletonbox}

\newpage
\fullproofpage{Homework 10, Problem 3}{Use this page for part (a).}

\newpage
\fullproofpage{Homework 10, Problem 3}{Use this page for part (b).}

\newpage
\begin{solutionbox}{Official Solution (Professor)}
\begin{enumerate}[label=(\alph*)]
  \item Let \(\{U_i\}_{i\in I}\) be an open cover of \(A\). Since \(A\) is closed, \(A^c\) is open. Thus
  \[
    \{A^c\}\cup\{U_i\}_{i\in I}
  \]
  is an open cover of \(B\), because
  \[
    B\subseteq X=A^c\cup A\subseteq A^c\cup\bigcup_{i\in I}U_i.
  \]
  Since \(B\) is compact, there is a finite subcover for \(B\). We may write it as
  \[
    \{A^c\}\cup\{U_i\}_{i\in J}
  \]
  for some finite \(J\subseteq I\). Then
  \[
    B\subseteq A^c\cup\bigcup_{i\in J}U_i.
  \]
  Since \(A\subseteq B\) and \(A\cap A^c=\varnothing\), it follows that
  \[
    A\subseteq \bigcup_{i\in J}U_i.
  \]
  Hence \(\{U_i\}_{i\in J}\) is a finite subcover of \(A\). Therefore \(A\) is compact.

  \item Since \(C\) and \(D\) are compact, they are closed by Lecture 27. Hence \(C\cap D\) is closed. Also \(C\cap D\subseteq C\), and \(C\) is compact. By part (a), \(C\cap D\) is compact.
\end{enumerate}
\end{solutionbox}

\newpage
\reflectionpage{Homework 10, Problem 3}

% ---------------------------------------------------------------------------
% HW10 Problem 4
% ---------------------------------------------------------------------------
\newpage
\subsection{Problem 4}

\begin{problemstatementbox}{Homework 10, Problem 4}
Let \((X,d)\) be a metric space and let \(Y\subseteq X\). Prove that if \(Y\) is compact, then \(Y\) is complete.
\end{problemstatementbox}

\newpage
\firstdiagnosis{Homework 10, Problem 4}

\newpage
\recipecard{Homework 10, Problem 4}

\newpage
\begin{selfexplanationbox}{Homework 10, Problem 4}
\workquestion{1. To prove completeness, what kind of sequence in \(Y\) should you start with?}{0.95in}
\workquestion{2. Which compactness theorem gives a convergent subsequence in \(Y\)?}{0.95in}
\workquestion{3. Why is a convergent subsequence not enough by itself to prove completeness?}{1.0in}
\workquestion{4. How does the Cauchy property transfer convergence from a subsequence to the full sequence?}{1.1in}
\workquestion{5. Why does the final limit lie in \(Y\)?}{0.9in}
\end{selfexplanationbox}

\newpage
\begin{skeletonbox}{Homework 10, Problem 4}
\textbf{Setup: prove every Cauchy sequence in \(Y\) converges in \(Y\).}
\workquestion{Let \(\{y_n\}\) be a Cauchy sequence in \(Y\):}{0.8in}
\workquestion{Use compactness and Lecture 26 to obtain a subsequence \(\{y_{k_n}\}\) and point \(y\in Y\) such that:}{0.95in}
\workquestion{Fix \(r>0\). From subsequence convergence choose \(N_1\) such that:}{0.85in}
\workquestion{From the Cauchy property choose \(N_2\) such that:}{0.85in}
\workquestion{Use triangle inequality with \(y_n\), \(y_{k_n}\), and \(y\):}{1.0in}
\workquestion{Conclude the whole sequence converges to \(y\in Y\):}{0.75in}
\end{skeletonbox}

\newpage
\fullproofpage{Homework 10, Problem 4}{Use this page for the complete proof.}

\newpage
\begin{solutionbox}{Official Solution (Professor)}
Let \(\{y_n\}_{n\ge1}\) be a Cauchy sequence in \(Y\). We show that it converges to some point of \(Y\).

Since \(Y\) is compact, by Lecture 26 there exists a subsequence \(\{y_{k_n}\}_{n\ge1}\) that converges to some point \(y\in Y\).

We prove the whole sequence converges to \(y\). Fix \(r>0\). Since \(y_{k_n}\to y\), there exists \(N_1\) such that
\[
  d(y_{k_n},y)<\frac r2
  \qquad\text{for all }n\ge N_1.
\]
Since \(\{y_n\}\) is Cauchy, there exists \(N_2\) such that
\[
  d(y_n,y_m)<\frac r2
  \qquad\text{for all }n,m\ge N_2.
\]
Set \(N=\max\{N_1,N_2\}\). For \(n\ge N\), we have \(k_n\ge n\ge N\), and so
\[
  d(y_n,y)\le d(y_n,y_{k_n})+d(y_{k_n},y)<\frac r2+\frac r2=r.
\]
Since \(r>0\) was arbitrary, \(y_n\to y\). Because \(y\in Y\), every Cauchy sequence in \(Y\) converges in \(Y\). Therefore \(Y\) is complete.
\end{solutionbox}

\newpage
\reflectionpage{Homework 10, Problem 4}

\newpage
\section{Homework 11}

% ---------------------------------------------------------------------------
% HW11 Problem 1
% ---------------------------------------------------------------------------
\newpage
\subsection{Problem 1}

\begin{problemstatementbox}{Homework 11, Problem 1}
Consider the metric space \((\R,d)\) with the usual metric \(d(x,y)=|x-y|\). Note that the set \(\Q\) is not an interval, and so from lecture we know that it must be disconnected.

\begin{enumerate}[label=(\alph*)]
  \item Find two sets \(A,B\subseteq\R\) that are nonempty and separated so that \(\Q=A\cup B\), and prove your answer.
  \item Find a set \(\varnothing\neq C\subsetneq\Q\) that is both open and closed in \(\Q\), and prove your answer.
\end{enumerate}
\end{problemstatementbox}

\newpage
\firstdiagnosis{Homework 11, Problem 1}

\newpage
\recipecard{Homework 11, Problem 1}

\newpage
\begin{selfexplanationbox}{Homework 11, Problem 1}
\workquestion{1. Why is cutting \(\Q\) at an irrational number better than cutting it at a rational number?}{1.05in}
\workquestion{2. What does it mean for two subsets of \(\R\) to be separated?}{0.95in}
\workquestion{3. What are the closures in \(\R\) of \((-\infty,\sqrt2)\cap\Q\) and \((\sqrt2,\infty)\cap\Q\)?}{1.05in}
\workquestion{4. Why does \(\sqrt2\notin\Q\) make the separation work?}{0.95in}
\workquestion{5. How does the subspace characterization of open and closed sets in \(\Q\) prove part (b)?}{1.05in}
\end{selfexplanationbox}

\newpage
\begin{skeletonbox}{Homework 11, Problem 1}
\textbf{Part (a): separated decomposition of \(\Q\).}
\workquestion{Choose an irrational cut point and define \(A=\)}{0.75in}
\workquestion{Define \(B=\)}{0.65in}
\workquestion{Check \(A\) and \(B\) are nonempty and \(A\cup B=\Q\):}{0.95in}
\workquestion{Compute \(\closure{A}\) and \(\closure{B}\) in \(\R\):}{0.95in}
\workquestion{Use those closures to prove \(\closure{A}\cap B=\varnothing\) and \(A\cap\closure{B}=\varnothing\):}{1.0in}

\textbf{Part (b): clopen subset of \(\Q\).}
\workquestion{Choose \(C=\)}{0.65in}
\workquestion{Prove \(C\) is nonempty and proper:}{0.75in}
\workquestion{Show \(C\) is open in \(\Q\) and closed in \(\Q\) using subspace results:}{1.0in}
\end{skeletonbox}

\newpage
\fullproofpage{Homework 11, Problem 1}{Use this page for part (a).}

\newpage
\fullproofpage{Homework 11, Problem 1}{Use this page for part (b).}

\newpage
\begin{solutionbox}{Student-Derived Solution (No official solution available)}
\begin{enumerate}[label=(\alph*)]
  \item Define
  \[
    A=(-\infty,\sqrt2)\cap\Q,
    \qquad
    B=(\sqrt2,\infty)\cap\Q.
  \]
  Both sets are nonempty, for instance \(1\in A\) and \(2\in B\). Also \(A\cup B=\Q\), since every rational number is either less than \(\sqrt2\) or greater than \(\sqrt2\), and \(\sqrt2\notin\Q\).

  In \(\R\),
  \[
    \closure{A}=(-\infty,\sqrt2],
    \qquad
    \closure{B}=[\sqrt2,\infty).
  \]
  Hence
  \[
    \closure{A}\cap B
    =
    (-\infty,\sqrt2]\cap\big((\sqrt2,\infty)\cap\Q\big)
    =
    \varnothing,
  \]
  and similarly
  \[
    A\cap\closure{B}
    =
    \big((-\infty,\sqrt2)\cap\Q\big)\cap[\sqrt2,\infty)
    =
    \varnothing.
  \]
  Therefore \(A\) and \(B\) are separated, nonempty, and satisfy \(\Q=A\cup B\).

  \item Let
  \[
    C=(-\infty,\sqrt2)\cap\Q.
  \]
  Then \(C\neq\varnothing\), since \(1\in C\), and \(C\subsetneq\Q\), since \(2\in\Q\setminus C\).

  The set \((-\infty,\sqrt2)\) is open in \(\R\), so \(C=(-\infty,\sqrt2)\cap\Q\) is open in the subspace \(\Q\). Also \((-\infty,\sqrt2]\) is closed in \(\R\), and since \(\sqrt2\notin\Q\),
  \[
    (-\infty,\sqrt2]\cap\Q=(-\infty,\sqrt2)\cap\Q=C.
  \]
  Thus \(C\) is closed in \(\Q\). Therefore \(C\) is a nonempty proper subset of \(\Q\) that is both open and closed in \(\Q\).
\end{enumerate}
\end{solutionbox}

\newpage
\reflectionpage{Homework 11, Problem 1}

% ---------------------------------------------------------------------------
% HW11 Problem 2
% ---------------------------------------------------------------------------
\newpage
\subsection{Problem 2}

\begin{problemstatementbox}{Homework 11, Problem 2}
Let \((X,d)\) be a metric space.

\begin{enumerate}[label=(\alph*)]
  \item Prove that if \(F\subseteq X\) is disconnected and closed in \(X\), then there exist \(A,B\subseteq X\) that are nonempty, separated, and closed in \(X\) so that \(F=A\cup B\).
  \item Prove that if \(U\subseteq X\) is disconnected and open in \(X\), then there exist \(A,B\subseteq X\) that are nonempty, separated, and open in \(X\) so that \(U=A\cup B\).
\end{enumerate}
\end{problemstatementbox}

\newpage
\firstdiagnosis{Homework 11, Problem 2}

\newpage
\recipecard{Homework 11, Problem 2}

\newpage
\begin{selfexplanationbox}{Homework 11, Problem 2}
\workquestion{1. Which Lecture 29 characterization of disconnectedness gives a nontrivial clopen subset of \(F\) or \(U\)?}{1.05in}
\workquestion{2. In part (a), why does closed in \(F\) plus \(F\) closed in \(X\) imply closed in \(X\)?}{1.0in}
\workquestion{3. In part (b), why does open in \(U\) plus \(U\) open in \(X\) imply open in \(X\)?}{1.0in}
\workquestion{4. Why are the two pieces nonempty?}{0.85in}
\workquestion{5. How do closedness or openness of the pieces imply they are separated?}{1.05in}
\end{selfexplanationbox}

\newpage
\begin{skeletonbox}{Homework 11, Problem 2}
\textbf{Part (a): closed disconnected set.}
\workquestion{Use disconnectedness of \(F\) to choose \(A\neq\varnothing,F\) that is both open and closed in \(F\). Set \(B=\)}{0.9in}
\workquestion{Explain why \(A\) and \(B\) are nonempty and \(F=A\cup B\):}{0.8in}
\workquestion{Use \(F\) closed in \(X\) to prove \(A\) and \(B\) are closed in \(X\):}{1.0in}
\workquestion{Use closedness and disjointness to prove separatedness:}{0.95in}

\textbf{Part (b): open disconnected set.}
\workquestion{Use disconnectedness of \(U\) to choose \(A\neq\varnothing,U\) clopen in \(U\), and set \(B=\)}{0.9in}
\workquestion{Use \(U\) open in \(X\) to prove \(A\) and \(B\) are open in \(X\):}{1.0in}
\workquestion{Use complements of open sets to prove separatedness:}{0.95in}
\end{skeletonbox}

\newpage
\fullproofpage{Homework 11, Problem 2}{Use this page for part (a).}

\newpage
\fullproofpage{Homework 11, Problem 2}{Use this page for part (b).}

\newpage
\begin{solutionbox}{Official Solution (Professor)}
\begin{enumerate}[label=(\alph*)]
  \item Since \(F\) is disconnected, by Lecture 29 there exists a set \(A\neq\varnothing,F\) that is both open and closed in \(F\). Set \(B=F\setminus A\).

  The sets \(A\) and \(B\) are nonempty because \(A\neq\varnothing\) and \(A\neq F\). Since \(A\) is closed in \(F\) and \(F\) is closed in \(X\), \(A\) is closed in \(X\) by Lecture 24. Also \(B=F\setminus A\) is closed in \(F\), because \(A\) is open in \(F\); hence \(B\) is closed in \(X\).

  Finally, \(A\cap B=\varnothing\). Since \(A\) and \(B\) are closed in \(X\), \(\closure{A}=A\) and \(\closure{B}=B\). Thus
  \[
    \closure{A}\cap B=\varnothing,
    \qquad
    A\cap\closure{B}=\varnothing.
  \]
  Therefore \(A\) and \(B\) are nonempty, separated, closed in \(X\), and \(F=A\cup B\).

  \item Since \(U\) is disconnected, by Lecture 29 there exists a set \(A\neq\varnothing,U\) that is both open and closed in \(U\). Set \(B=U\setminus A\).

  The sets \(A\) and \(B\) are nonempty. Since \(A\) is open in \(U\) and \(U\) is open in \(X\), \(A\) is open in \(X\) by Lecture 24. Since \(A\) is closed in \(U\), \(B=U\setminus A\) is open in \(U\); because \(U\) is open in \(X\), \(B\) is open in \(X\).

  Since \(A\cap B=\varnothing\), we have \(A\subseteq B^c\). Because \(B\) is open, \(B^c\) is closed, so \(\closure{A}\subseteq B^c\), and hence \(\closure{A}\cap B=\varnothing\). Similarly, \(\closure{B}\cap A=\varnothing\). Therefore \(A\) and \(B\) are separated, open in \(X\), nonempty, and \(U=A\cup B\).
\end{enumerate}
\end{solutionbox}

\newpage
\reflectionpage{Homework 11, Problem 2}

% ---------------------------------------------------------------------------
% HW11 Problem 3
% ---------------------------------------------------------------------------
\newpage
\subsection{Problem 3}

\begin{problemstatementbox}{Homework 11, Problem 3}
Suppose that \((X,d)\) is a metric space such that for any pair of points \(a,b\in X\), there exists a connected set \(C\subseteq X\) so that \(a,b\in C\). Prove that \(X\) is connected.
\end{problemstatementbox}

\newpage
\firstdiagnosis{Homework 11, Problem 3}

\newpage
\recipecard{Homework 11, Problem 3}

\newpage
\begin{selfexplanationbox}{Homework 11, Problem 3}
\workquestion{1. Why is contradiction the natural proof method here?}{0.95in}
\workquestion{2. What does a separation of \(X\) give you?}{0.95in}
\workquestion{3. Why can you choose one point \(a\) from one side and one point \(b\) from the other?}{0.9in}
\workquestion{4. Which Lecture 29 result says a connected set inside a union of separated sets lies in one side?}{1.05in}
\workquestion{5. Why does the connected set containing both \(a\) and \(b\) contradict the separation?}{1.05in}
\end{selfexplanationbox}

\newpage
\begin{skeletonbox}{Homework 11, Problem 3}
\textbf{Contradiction setup.}
\workquestion{Assume \(X\) is disconnected. Then choose nonempty separated sets \(A,B\subseteq X\) such that:}{0.95in}
\workquestion{Choose \(a\in A\) and \(b\in B\). By hypothesis, choose connected \(C\subseteq X\) such that:}{0.95in}
\workquestion{Since \(C\subseteq X=A\cup B\), apply the connected-set/separated-union result:}{1.0in}
\workquestion{Use \(a\in A\cap C\) and \(b\in B\cap C\) to rule out both alternatives:}{1.05in}
\workquestion{Conclude \(X\) is connected:}{0.75in}
\end{skeletonbox}

\newpage
\fullproofpage{Homework 11, Problem 3}{Use this page for the complete proof.}

\newpage
\begin{solutionbox}{Official Solution (Professor)}
Suppose toward a contradiction that \(X\) is disconnected. Then there exist sets \(A,B\subseteq X\) such that
\[
  X=A\cup B,
  \qquad
  A\neq\varnothing,
  \qquad
  B\neq\varnothing,
\]
and \(A\) and \(B\) are separated.

Choose \(a\in A\) and \(b\in B\). By the hypothesis, there exists a connected set \(C\subseteq X\) such that \(a,b\in C\). Since \(C\subseteq X=A\cup B\), and \(A\) and \(B\) are separated, Lecture 29 implies that either \(C\subseteq A\) or \(C\subseteq B\).

Because \(a\in A\cap C\), the only possible case is \(C\subseteq A\). But then \(B\cap C=\varnothing\), contradicting \(b\in B\cap C\). Therefore \(X\) cannot be disconnected, so \(X\) is connected.
\end{solutionbox}

\newpage
\reflectionpage{Homework 11, Problem 3}

% ---------------------------------------------------------------------------
% HW11 Problem 4
% ---------------------------------------------------------------------------
\newpage
\subsection{Problem 4}

\begin{problemstatementbox}{Homework 11, Problem 4}
Let \((X,d)\) be a metric space and \(A\subseteq X\). Prove that if \(A\) is connected, then \(\closure{A}\) is connected.
\end{problemstatementbox}

\newpage
\firstdiagnosis{Homework 11, Problem 4}

\newpage
\recipecard{Homework 11, Problem 4}

\newpage
\begin{selfexplanationbox}{Homework 11, Problem 4}
\workquestion{1. Why should the proof begin by assuming \(\closure{A}\) is disconnected?}{0.95in}
\workquestion{2. If \(\closure{A}=B\cup C\) with \(B,C\) separated, what does connectedness of \(A\subseteq B\cup C\) imply?}{1.05in}
\workquestion{3. Why can one assume \(A\subseteq B\) after relabeling?}{0.9in}
\workquestion{4. How does \(A\subseteq B\) imply \(\closure{A}\subseteq\closure{B}\)?}{0.95in}
\workquestion{5. Why does separatedness then force the other side to be empty?}{1.0in}
\end{selfexplanationbox}

\newpage
\begin{skeletonbox}{Homework 11, Problem 4}
\textbf{Contradiction setup.}
\workquestion{Assume \(\closure{A}\) is disconnected, so \(\closure{A}=B\cup C\) with \(B,C\) nonempty and separated:}{0.95in}
\workquestion{Apply connectedness of \(A\) to the separated union \(B\cup C\):}{0.95in}
\workquestion{After relabeling, assume \(A\subseteq B\). Then conclude \(\closure{A}\subseteq\)}{0.85in}
\workquestion{Use separatedness \( \closure{B}\cap C=\varnothing\) to show \(\closure{A}\cap C=\varnothing\):}{0.95in}
\workquestion{Explain why this contradicts \(C\neq\varnothing\) and \(C\subseteq\closure{A}\):}{0.95in}
\end{skeletonbox}

\newpage
\fullproofpage{Homework 11, Problem 4}{Use this page for the complete proof.}

\newpage
\begin{solutionbox}{Official Solution (Professor)}
Suppose toward a contradiction that \(\closure{A}\) is disconnected. Then there exist sets \(B,C\subseteq X\) such that
\[
  \closure{A}=B\cup C,
  \qquad
  B\neq\varnothing,
  \qquad
  C\neq\varnothing,
\]
and \(B\) and \(C\) are separated.

Since \(A\) is connected and \(A\subseteq\closure{A}=B\cup C\), Lecture 29 implies that either \(A\subseteq B\) or \(A\subseteq C\). After swapping \(B\) and \(C\) if necessary, assume \(A\subseteq B\). Then
\[
  \closure{A}\subseteq\closure{B}.
\]
Because \(B\) and \(C\) are separated, \(\closure{B}\cap C=\varnothing\). Hence
\[
  \closure{A}\cap C=\varnothing.
\]
But \(C\subseteq\closure{A}\), since \(\closure{A}=B\cup C\). Therefore \(C=\varnothing\), a contradiction. Thus \(\closure{A}\) is connected.
\end{solutionbox}

\newpage
\reflectionpage{Homework 11, Problem 4}

% ---------------------------------------------------------------------------
% HW11 Problem 5
% ---------------------------------------------------------------------------
\newpage
\subsection{Problem 5}

\begin{problemstatementbox}{Homework 11, Problem 5}
Let \((X,d)\) be a metric space and \(A,B\subseteq X\). Prove that if \(A\) and \(B\) are both connected and \(A\cap B\neq\varnothing\), then \(A\cup B\) is connected.
\end{problemstatementbox}

\newpage
\firstdiagnosis{Homework 11, Problem 5}

\newpage
\recipecard{Homework 11, Problem 5}

\newpage
\begin{selfexplanationbox}{Homework 11, Problem 5}
\workquestion{1. Why does the assumption \(A\cap B\neq\varnothing\) prevent \(A\) and \(B\) from lying in different sides of a separation?}{1.1in}
\workquestion{2. If \(A\cup B=C\cup D\) is a separation, what does connectedness of \(A\) imply?}{1.0in}
\workquestion{3. What does connectedness of \(B\) imply?}{0.95in}
\workquestion{4. How does a shared point force \(A\) and \(B\) into the same separated side?}{1.05in}
\workquestion{5. Why does this make the other side empty?}{0.9in}
\end{selfexplanationbox}

\newpage
\begin{skeletonbox}{Homework 11, Problem 5}
\textbf{Contradiction setup.}
\workquestion{Assume \(A\cup B\) is disconnected. Choose nonempty separated \(C,D\subseteq X\) such that:}{0.95in}
\workquestion{Use connectedness of \(A\subseteq C\cup D\) to conclude:}{0.85in}
\workquestion{After relabeling, assume \(A\subseteq C\). Use \(A\cap B\neq\varnothing\) to show \(B\) has a point in \(C\):}{1.0in}
\workquestion{Use connectedness of \(B\) to conclude \(B\subseteq C\):}{0.9in}
\workquestion{Then \(A\cup B\subseteq C\), contradicting:}{0.85in}
\end{skeletonbox}

\newpage
\fullproofpage{Homework 11, Problem 5}{Use this page for the complete proof.}

\newpage
\begin{solutionbox}{Official Solution (Professor)}
Suppose toward a contradiction that \(A\cup B\) is disconnected. Then there exist sets \(C,D\subseteq X\) such that
\[
  A\cup B=C\cup D,
  \qquad
  C,D\neq\varnothing,
\]
and \(C\) and \(D\) are separated.

Since \(A\subseteq C\cup D\) and \(A\) is connected, Lecture 29 implies that either \(A\subseteq C\) or \(A\subseteq D\). After swapping \(C\) and \(D\), assume \(A\subseteq C\).

Similarly, since \(B\subseteq C\cup D\) and \(B\) is connected, either \(B\subseteq C\) or \(B\subseteq D\). But \(A\subseteq C\) and \(A\cap B\neq\varnothing\), so at least one point of \(B\) lies in \(C\). Therefore \(B\) cannot be contained in \(D\), and we must have \(B\subseteq C\).

Thus \(A\cup B\subseteq C\). Since \(A\cup B=C\cup D\), this forces \(D=\varnothing\), contradicting the separation. Therefore \(A\cup B\) is connected.
\end{solutionbox}

\newpage
\reflectionpage{Homework 11, Problem 5}

% ---------------------------------------------------------------------------
% HW11 Problem 6
% ---------------------------------------------------------------------------
\newpage
\subsection{Problem 6}

\begin{problemstatementbox}{Homework 11, Problem 6}
Let \((X,d)\) be a metric space.

\begin{enumerate}[label=(\alph*)]
  \item Let \(A,B,C\subseteq X\). Prove that if the sets \(A\) and \(B\) are separated and the sets \(A\) and \(C\) are separated, then the sets \(A\) and \(B\cup C\) are separated.
  \item Prove that if \(F,G\subseteq X\) are closed and \(F\cup G\) and \(F\cap G\) are connected, then \(F\) is connected.
\end{enumerate}
\end{problemstatementbox}

\newpage
\firstdiagnosis{Homework 11, Problem 6}

\newpage
\recipecard{Homework 11, Problem 6}

\newpage
\begin{selfexplanationbox}{Homework 11, Problem 6}
\workquestion{1. What two closure-intersection conditions define separated sets?}{0.95in}
\workquestion{2. In part (a), how do unions distribute over intersections with \(\closure{A}\)?}{1.0in}
\workquestion{3. For part (b), why is contradiction a natural strategy?}{0.9in}
\workquestion{4. If \(F=A\cup B\) is a separation, how does connectedness of \(F\cap G\) force \(F\cap G\) into one side?}{1.1in}
\workquestion{5. Why does proving \(B\) is separated from \(A\cup G\) contradict connectedness of \(F\cup G\)?}{1.05in}
\end{selfexplanationbox}

\newpage
\begin{skeletonbox}{Homework 11, Problem 6}
\textbf{Part (a): separated from a union.}
\workquestion{Write what it means that \(A\) is separated from \(B\) and from \(C\):}{0.95in}
\workquestion{Compute \(\closure{A}\cap(B\cup C)\):}{0.85in}
\workquestion{Use \(\closure{B\cup C}=\closure{B}\cup\closure{C}\) to compute \(A\cap\closure{B\cup C}\):}{0.95in}
\workquestion{Conclude \(A\) and \(B\cup C\) are separated:}{0.65in}

\textbf{Part (b): contradiction.}
\workquestion{Assume \(F\) is disconnected and choose \(F=A\cup B\) with \(A,B\) nonempty separated:}{0.85in}
\workquestion{Use connectedness of \(F\cap G\) to assume, after relabeling, \(F\cap G\subseteq A\):}{0.95in}
\workquestion{Show \(B\) and \(G\) are separated, using closedness of \(G\):}{0.95in}
\workquestion{Apply part (a) to show \(B\) and \(A\cup G\) are separated, yielding a separation of \(F\cup G\):}{1.0in}
\end{skeletonbox}

\newpage
\fullproofpage{Homework 11, Problem 6}{Use this page for part (a).}

\newpage
\fullproofpage{Homework 11, Problem 6}{Use this page for part (b).}

\newpage
\begin{solutionbox}{Official Solution (Professor)}
\begin{enumerate}[label=(\alph*)]
  \item Since \(A\) and \(B\) are separated,
  \[
    \closure{A}\cap B=\varnothing
    \qquad\text{and}\qquad
    A\cap\closure{B}=\varnothing.
  \]
  Since \(A\) and \(C\) are separated,
  \[
    \closure{A}\cap C=\varnothing
    \qquad\text{and}\qquad
    A\cap\closure{C}=\varnothing.
  \]
  Therefore
  \[
    \closure{A}\cap(B\cup C)
    =
    (\closure{A}\cap B)\cup(\closure{A}\cap C)
    =
    \varnothing.
  \]
  Also, by Lecture 23,
  \[
    \closure{B\cup C}=\closure{B}\cup\closure{C}.
  \]
  Hence
  \[
    A\cap\closure{B\cup C}
    =
    (A\cap\closure{B})\cup(A\cap\closure{C})
    =
    \varnothing.
  \]
  Thus \(A\) and \(B\cup C\) are separated.

  \item Suppose toward a contradiction that \(F\) is disconnected. Then there exist \(A,B\subseteq X\) such that
  \[
    F=A\cup B,
    \qquad
    A,B\neq\varnothing,
  \]
  and \(A\) and \(B\) are separated.

  Since \(F\cap G\) is connected and \(F\cap G\subseteq A\cup B\), Lecture 29 implies that either \(F\cap G\subseteq A\) or \(F\cap G\subseteq B\). After swapping \(A\) and \(B\) if necessary, assume \(F\cap G\subseteq A\).

  We claim that \(B\) and \(G\) are separated. First \(B\cap G=\varnothing\), because any point in \(B\cap G\) would lie in \(F\cap G\subseteq A\), contradicting \(A\cap B=\varnothing\). Since \(G\) is closed, \(\closure{G}=G\), and this gives \(B\cap\closure{G}=\varnothing\). Also, if \(x\in G=\closure{G}\), then either \(x\in G\setminus F\) or \(x\in F\cap G\subseteq A\), so \(x\notin B\). Hence \(\closure{B}\cap G=\varnothing\). Thus \(B\) and \(G\) are separated.

  Since \(A\) and \(B\) are separated, part (a) implies that \(B\) and \(A\cup G\) are separated. Both are nonempty, and
  \[
    F\cup G=(A\cup B)\cup G=(A\cup G)\cup B.
  \]
  Therefore \(F\cup G\) is disconnected, contradicting the hypothesis. Hence \(F\) is connected.
\end{enumerate}
\end{solutionbox}

\newpage
\reflectionpage{Homework 11, Problem 6}

\newpage
\section{Homework 12}

% ---------------------------------------------------------------------------
% HW12 Problem 1
% ---------------------------------------------------------------------------
\newpage
\subsection{Problem 1}

\begin{problemstatementbox}{Homework 12, Problem 1}
Let \((X,d_X)\) and \((Y,d_Y)\) be metric spaces, \(f:X\to Y\), and \(x_0\in X\). Prove that \(f\) is continuous at \(x_0\) if and only if for any open set \(V\subseteq Y\) containing \(f(x_0)\), there is an open set \(U\subseteq X\) containing \(x_0\) so that \(x\in U\) implies \(f(x)\in V\).
\end{problemstatementbox}

\newpage
\firstdiagnosis{Homework 12, Problem 1}

\newpage
\recipecard{Homework 12, Problem 1}

\newpage
\begin{selfexplanationbox}{Homework 12, Problem 1}
\workquestion{1. What is the \(\varepsilon\)-\(\delta\) definition of continuity at \(x_0\)?}{1.0in}
\workquestion{2. In the forward direction, how does openness of \(V\) produce a ball around \(f(x_0)\)?}{1.0in}
\workquestion{3. Why is \(U=B_\delta(x_0)\) the natural choice after applying continuity?}{0.95in}
\workquestion{4. In the reverse direction, why should \(V=B_\varepsilon(f(x_0))\) be used?}{1.0in}
\workquestion{5. How does openness of \(U\) produce the \(\delta\) needed for continuity?}{1.0in}
\end{selfexplanationbox}

\newpage
\begin{skeletonbox}{Homework 12, Problem 1}
\textbf{Forward direction.}
\workquestion{Assume \(f\) is continuous at \(x_0\). Fix open \(V\subseteq Y\) with \(f(x_0)\in V\). Choose \(\varepsilon>0\) so that:}{0.95in}
\workquestion{Use continuity at \(x_0\) to choose \(\delta>0\) so that:}{0.95in}
\workquestion{Set \(U=\)}{0.65in}
\workquestion{Check \(U\) is open, contains \(x_0\), and \(x\in U\Rightarrow f(x)\in V\):}{0.95in}

\textbf{Reverse direction.}
\workquestion{Assume the open-neighborhood condition. Fix \(\varepsilon>0\) and set \(V=\)}{0.8in}
\workquestion{Use the hypothesis to get open \(U\ni x_0\) with \(f(U)\subseteq V\):}{0.85in}
\workquestion{Use openness of \(U\) to choose \(\delta>0\) with \(B_\delta(x_0)\subseteq U\):}{0.85in}
\workquestion{Conclude the \(\varepsilon\)-\(\delta\) condition:}{0.8in}
\end{skeletonbox}

\newpage
\fullproofpage{Homework 12, Problem 1}{Use this page for the forward direction.}

\newpage
\fullproofpage{Homework 12, Problem 1}{Use this page for the reverse direction.}

\newpage
\begin{solutionbox}{Student-Derived Solution (No official solution available yet)}
\(\Rightarrow\) Assume \(f\) is continuous at \(x_0\). Let \(V\subseteq Y\) be open with \(f(x_0)\in V\). Since \(V\) is open, there is \(\varepsilon>0\) such that
\[
  B_\varepsilon(f(x_0))\subseteq V.
\]
By continuity at \(x_0\), there is \(\delta>0\) such that
\[
  d_X(x,x_0)<\delta
  \quad\Longrightarrow\quad
  d_Y(f(x),f(x_0))<\varepsilon.
\]
Let \(U=B_\delta(x_0)\). Then \(U\) is open, \(x_0\in U\), and if \(x\in U\), then \(f(x)\in B_\varepsilon(f(x_0))\subseteq V\).

\(\Leftarrow\) Assume the open-set condition. Let \(\varepsilon>0\), and set \(V=B_\varepsilon(f(x_0))\). Then \(V\) is open and contains \(f(x_0)\), so there is an open set \(U\subseteq X\) with \(x_0\in U\) and \(f(U)\subseteq V\). Since \(U\) is open and \(x_0\in U\), there is \(\delta>0\) such that \(B_\delta(x_0)\subseteq U\). Thus if \(d_X(x,x_0)<\delta\), then \(x\in U\), so \(f(x)\in V\), meaning
\[
  d_Y(f(x),f(x_0))<\varepsilon.
\]
Therefore \(f\) is continuous at \(x_0\).
\end{solutionbox}

\newpage
\reflectionpage{Homework 12, Problem 1}

% ---------------------------------------------------------------------------
% HW12 Problem 2
% ---------------------------------------------------------------------------
\newpage
\subsection{Problem 2}

\begin{problemstatementbox}{Homework 12, Problem 2}
Let \((X,d_X)\) and \((Y,d_Y)\) be metric spaces and \(f:X\to Y\) a function. Prove that \(f\) is continuous if and only if \(f(\closure{A})\subseteq\closure{f(A)}\) for any \(A\subseteq X\).
\end{problemstatementbox}

\newpage
\firstdiagnosis{Homework 12, Problem 2}

\newpage
\recipecard{Homework 12, Problem 2}

\newpage
\begin{selfexplanationbox}{Homework 12, Problem 2}
\workquestion{1. What does \(x\in\closure{A}\) mean sequentially in a metric space?}{1.0in}
\workquestion{2. In the forward direction, how does continuity move a convergent sequence from \(X\) to \(Y\)?}{1.0in}
\workquestion{3. Why does \(f(a_n)\to f(x)\) prove \(f(x)\in\closure{f(A)}\)?}{0.95in}
\workquestion{4. In the reverse direction, why choose \(A=X\setminus f^{-1}(V)\) for an open \(V\subseteq Y\)?}{1.1in}
\workquestion{5. Why does \(x\notin\closure{A}\) produce an open neighborhood \(U\) of \(x\) with \(f(U)\subseteq V\)?}{1.05in}
\end{selfexplanationbox}

\newpage
\begin{skeletonbox}{Homework 12, Problem 2}
\textbf{Forward direction.}
\workquestion{Assume \(f\) is continuous. Fix \(A\subseteq X\) and \(x\in\closure{A}\). Choose a sequence \(a_n\in A\) such that:}{0.95in}
\workquestion{Use continuity to show \(f(a_n)\to\)}{0.75in}
\workquestion{Conclude \(f(x)\in\closure{f(A)}\), hence \(f(\closure{A})\subseteq\closure{f(A)}\):}{0.9in}

\textbf{Reverse direction.}
\workquestion{Assume \(f(\closure{A})\subseteq\closure{f(A)}\) for every \(A\subseteq X\). Fix open \(V\subseteq Y\) with \(f(x)\in V\):}{0.85in}
\workquestion{Define \(A=X\setminus f^{-1}(V)\). Suppose for contradiction that \(x\in\closure{A}\):}{0.9in}
\workquestion{Use the closure hypothesis to force \(f(x)\in\closure{f(A)}\):}{0.8in}
\workquestion{Use \(f(A)\subseteq Y\setminus V\) and \(Y\setminus V\) closed to contradict \(f(x)\in V\):}{0.9in}
\workquestion{From \(x\notin\closure{A}\), build an open \(U\ni x\) with \(U\subseteq f^{-1}(V)\), then apply Problem 1:}{0.95in}
\end{skeletonbox}

\newpage
\fullproofpage{Homework 12, Problem 2}{Use this page for the forward direction.}

\newpage
\fullproofpage{Homework 12, Problem 2}{Use this page for the reverse direction.}

\newpage
\begin{solutionbox}{Student-Derived Solution (No official solution available yet)}
\(\Rightarrow\) Assume \(f\) is continuous. Let \(A\subseteq X\), and let \(x\in\closure{A}\). By the sequential characterization of closure, there is a sequence \(\{a_n\}\subseteq A\) such that \(a_n\to x\). Since \(f\) is continuous,
\[
  f(a_n)\to f(x).
\]
Each \(f(a_n)\in f(A)\), so \(f(x)\in\closure{f(A)}\). Therefore
\[
  f(\closure{A})\subseteq\closure{f(A)}.
\]

\(\Leftarrow\) Assume that for every \(A\subseteq X\),
\[
  f(\closure{A})\subseteq\closure{f(A)}.
\]
We prove continuity using Problem 1. Let \(V\subseteq Y\) be open and suppose \(f(x)\in V\). Set
\[
  A=X\setminus f^{-1}(V).
\]
We claim \(x\notin\closure{A}\). If \(x\in\closure{A}\), then by the hypothesis
\[
  f(x)\in f(\closure{A})\subseteq\closure{f(A)}.
\]
But if \(a\in A\), then \(f(a)\notin V\), so \(f(A)\subseteq Y\setminus V\). Since \(Y\setminus V\) is closed,
\[
  \closure{f(A)}\subseteq Y\setminus V.
\]
This would imply \(f(x)\notin V\), contradicting \(f(x)\in V\). Thus \(x\notin\closure{A}\).

Since \(x\notin\closure{A}\), there exists an open set \(U\subseteq X\) containing \(x\) such that \(U\cap A=\varnothing\). Hence
\[
  U\subseteq X\setminus A=f^{-1}(V),
\]
so \(f(U)\subseteq V\). By Problem 1, \(f\) is continuous at \(x\). Since \(x\in X\) was arbitrary, \(f\) is continuous.
\end{solutionbox}

\newpage
\reflectionpage{Homework 12, Problem 2}

% ---------------------------------------------------------------------------
% HW12 Problem 3
% ---------------------------------------------------------------------------
\newpage
\subsection{Problem 3}

\begin{problemstatementbox}{Homework 12, Problem 3}
Let \(f:\R\to\R\) be the function
\[
  f(x)=\frac{1}{1+x^2}.
\]

\begin{enumerate}[label=(\alph*)]
  \item Prove that \(f\) is continuous.
  \item Find an open set \(U\subseteq\R\) so that \(f(U)\) is not open.
  \item Find a closed set \(F\subseteq\R\) so that \(f(F)\) is not closed.
  \item Find a set \(A\subseteq\R\) so that \(f(\closure{A})\neq\closure{f(A)}\).
\end{enumerate}
\end{problemstatementbox}

\newpage
\firstdiagnosis{Homework 12, Problem 3}

\newpage
\recipecard{Homework 12, Problem 3}

\newpage
\begin{selfexplanationbox}{Homework 12, Problem 3}
\workquestion{1. Which standard limit laws prove continuity of \(f(x)=1/(1+x^2)\)?}{0.95in}
\workquestion{2. What is the range of \(f\) on all of \(\R\)?}{0.9in}
\workquestion{3. Why is \((0,1]\) not open in \(\R\)?}{0.9in}
\workquestion{4. Why is \((0,1]\) not closed in \(\R\)?}{0.9in}
\workquestion{5. For \(A=(0,\infty)\), compare \(f(\closure{A})\) with \(\closure{f(A)}\). Where does the extra endpoint come from?}{1.15in}
\end{selfexplanationbox}

\newpage
\begin{skeletonbox}{Homework 12, Problem 3}
\textbf{Part (a): continuity.}
\workquestion{Use sequences or limit laws: if \(x_n\to x\), then \(x_n^2\to\)}{0.8in}
\workquestion{Conclude \(1+x_n^2\to1+x^2\), and since denominators stay nonzero:}{0.85in}

\textbf{Parts (b) and (c): image counterexamples.}
\workquestion{For an open set, try \(U=\R\). Compute \(f(U)=\)}{0.75in}
\workquestion{Explain why that image is not open:}{0.75in}
\workquestion{For a closed set, try \(F=[0,\infty)\). Compute \(f(F)=\)}{0.75in}
\workquestion{Explain why that image is not closed:}{0.75in}

\textbf{Part (d): closure/image mismatch.}
\workquestion{Try \(A=(0,\infty)\). Then \(\closure{A}=\)}{0.65in}
\workquestion{Compute \(f(\closure{A})\), \(f(A)\), and \(\closure{f(A)}\):}{1.0in}
\end{skeletonbox}

\newpage
\fullproofpage{Homework 12, Problem 3}{Use this page for parts (a) and (b).}

\newpage
\fullproofpage{Homework 12, Problem 3}{Use this page for parts (c) and (d).}

\newpage
\begin{solutionbox}{Student-Derived Solution (No official solution available yet)}
\begin{enumerate}[label=(\alph*)]
  \item Let \(x_n\to x\) in \(\R\). Then by the standard limit laws,
  \[
    x_n^2\to x^2,
    \qquad
    1+x_n^2\to 1+x^2.
  \]
  Since \(1+x^2\neq0\), the reciprocal limit law gives
  \[
    \frac{1}{1+x_n^2}\to \frac{1}{1+x^2}.
  \]
  Hence \(f(x_n)\to f(x)\), so \(f\) is continuous.

  \item Let \(U=\R\), which is open. Then
  \[
    f(U)=(0,1].
  \]
  This is not open in \(\R\), since \(1\in(0,1]\) but no open interval around \(1\) is contained in \((0,1]\).

  \item Let \(F=[0,\infty)\), which is closed in \(\R\). Then
  \[
    f(F)=(0,1].
  \]
  This set is not closed in \(\R\), since \(0\) is a limit point of \((0,1]\) but \(0\notin(0,1]\).

  \item Let \(A=(0,\infty)\). Then \(\closure{A}=[0,\infty)\), so
  \[
    f(\closure{A})=f([0,\infty))=(0,1].
  \]
  On the other hand,
  \[
    f(A)=f((0,\infty))=(0,1),
    \qquad
    \closure{f(A)}=[0,1].
  \]
  Therefore
  \[
    f(\closure{A})=(0,1]\neq[0,1]=\closure{f(A)}.
  \]
\end{enumerate}
\end{solutionbox}

\newpage
\reflectionpage{Homework 12, Problem 3}

% ---------------------------------------------------------------------------
% HW12 Problem 4
% ---------------------------------------------------------------------------
\newpage
\subsection{Problem 4}

\begin{problemstatementbox}{Homework 12, Problem 4}
Consider the metric spaces \((\R,|\cdot|)\) and \((\R^2,d_2)\).

\begin{enumerate}[label=(\alph*)]
  \item Prove that the functions \(f,g:\R^2\to\R\) given by \(f(x,y)=x\) and \(g(x,y)=y\) are continuous. (Hint: Use sequences.)
  \item Prove that if \(h:[0,1]\to\R^2\) is a continuous function such that \(h(0)=(x_0,y_0)\), \(h(1)=(x_1,y_1)\), and \(x_0<0<x_1\), then there exists a point \(t\in(0,1)\) so that \(h(t)\) lies on the \(y\)-axis.
\end{enumerate}
\end{problemstatementbox}

\newpage
\firstdiagnosis{Homework 12, Problem 4}

\newpage
\recipecard{Homework 12, Problem 4}

\newpage
\begin{selfexplanationbox}{Homework 12, Problem 4}
\workquestion{1. How does convergence in \((\R^2,d_2)\) imply coordinatewise convergence?}{1.0in}
\workquestion{2. Why does \(|x_n-x|\le d_2((x_n,y_n),(x,y))\) prove continuity of the first projection?}{1.05in}
\workquestion{3. What scalar function should be formed from the path \(h\) in part (b)?}{0.95in}
\workquestion{4. Why is \(\phi(t)=f(h(t))\) continuous?}{0.9in}
\workquestion{5. How does the Intermediate Value Theorem produce a point on the \(y\)-axis?}{1.1in}
\end{selfexplanationbox}

\newpage
\begin{skeletonbox}{Homework 12, Problem 4}
\textbf{Part (a): projections are continuous.}
\workquestion{Let \((x_n,y_n)\to(x,y)\) in \((\R^2,d_2)\). Write the coordinate inequality for \(x_n\):}{0.95in}
\workquestion{Use squeeze to prove \(x_n\to x\), hence \(f(x_n,y_n)\to f(x,y)\):}{0.9in}
\workquestion{Repeat the coordinate inequality for \(y_n\):}{0.85in}

\textbf{Part (b): path crosses the \(y\)-axis.}
\workquestion{Define \(\phi:[0,1]\to\R\) by \(\phi(t)=\)}{0.7in}
\workquestion{Explain why \(\phi\) is continuous:}{0.85in}
\workquestion{Compute \(\phi(0)\) and \(\phi(1)\), and record their signs:}{0.9in}
\workquestion{Apply the Intermediate Value Theorem to get \(t\in(0,1)\) with \(\phi(t)=0\):}{0.9in}
\workquestion{Translate \(\phi(t)=0\) into "\(h(t)\) lies on the \(y\)-axis":}{0.8in}
\end{skeletonbox}

\newpage
\fullproofpage{Homework 12, Problem 4}{Use this page for part (a).}

\newpage
\fullproofpage{Homework 12, Problem 4}{Use this page for part (b).}

\newpage
\begin{solutionbox}{Student-Derived Solution (No official solution available yet)}
\begin{enumerate}[label=(\alph*)]
  \item We prove continuity of \(f(x,y)=x\). Let \((x_n,y_n)\to(x,y)\) in \((\R^2,d_2)\). Then
  \[
    d_2((x_n,y_n),(x,y))
    =
    \sqrt{(x_n-x)^2+(y_n-y)^2}\to0.
  \]
  Since
  \[
    |x_n-x|\le d_2((x_n,y_n),(x,y)),
  \]
  the squeeze theorem gives \(x_n\to x\). Thus
  \[
    f(x_n,y_n)=x_n\to x=f(x,y),
  \]
  so \(f\) is continuous. The proof for \(g(x,y)=y\) is the same: use
  \[
    |y_n-y|\le d_2((x_n,y_n),(x,y))
  \]
  to get \(y_n\to y\), hence \(g(x_n,y_n)\to g(x,y)\).

  \item Let \(f:\R^2\to\R\) be the first-coordinate projection from part (a), and define
  \[
    \phi:[0,1]\to\R,
    \qquad
    \phi(t)=f(h(t)).
  \]
  Since \(h\) is continuous and \(f\) is continuous, \(\phi=f\circ h\) is continuous. Also
  \[
    \phi(0)=f(h(0))=f(x_0,y_0)=x_0<0,
  \]
  and
  \[
    \phi(1)=f(h(1))=f(x_1,y_1)=x_1>0.
  \]
  By the Intermediate Value Theorem, there exists \(t\in(0,1)\) such that \(\phi(t)=0\). This means the first coordinate of \(h(t)\) is \(0\), so \(h(t)\in\{(0,y):y\in\R\}\). Therefore \(h(t)\) lies on the \(y\)-axis.
\end{enumerate}
\end{solutionbox}

\newpage
\reflectionpage{Homework 12, Problem 4}

% ---------------------------------------------------------------------------
% HW12 Problem 5
% ---------------------------------------------------------------------------
\newpage
\subsection{Problem 5}

\begin{problemstatementbox}{Homework 12, Problem 5}
Going forward, you may assume (no proof necessary) any fact about \(\sin x\) and \(\cos x\) from calculus, such as \(\sin,\cos:\R\to[-1,1]\) are continuous functions.

Prove that the function \(f:\R^2\to\R\) given by
\[
  f(x,y)=x^2+y^2+\sin(x-y)
\]
has a global minimum. (Hint: Notice that we cannot directly apply the extreme value theorem here. However, we can apply it to any compact subset of \(\R^2\).)
\end{problemstatementbox}

\newpage
\firstdiagnosis{Homework 12, Problem 5}

\newpage
\recipecard{Homework 12, Problem 5}

\newpage
\begin{selfexplanationbox}{Homework 12, Problem 5}
\workquestion{1. Why can the extreme value theorem not be applied directly on all of \(\R^2\)?}{1.0in}
\workquestion{2. Why is \(f\) continuous?}{0.95in}
\workquestion{3. What lower bound follows from \(\sin(x-y)\ge -1\)?}{0.95in}
\workquestion{4. Why does the closed unit disk capture a global minimizer?}{1.1in}
\workquestion{5. Why is it useful to find one point inside the disk where \(f<0\)?}{1.0in}
\end{selfexplanationbox}

\newpage
\begin{skeletonbox}{Homework 12, Problem 5}
\textbf{Continuity and coercive estimate.}
\workquestion{Show \(f\) is continuous as a sum/composition of continuous functions:}{0.85in}
\workquestion{Use \(\sin(x-y)\ge -1\) to prove \(f(x,y)\ge\)}{0.75in}
\workquestion{Conclude that if \(x^2+y^2\ge1\), then \(f(x,y)\ge\)}{0.65in}

\textbf{Find a compact region where the minimum occurs.}
\workquestion{Choose \(p=(-1/2,1/2)\). Compute \(f(p)=\)}{0.8in}
\workquestion{Explain why \(f(p)<0\):}{0.75in}
\workquestion{Let \(K=\{(x,y):x^2+y^2\le1\}\). Explain why \(K\) is compact:}{0.85in}
\workquestion{Apply the extreme value theorem to \(f|_K\):}{0.85in}
\workquestion{Compare points inside \(K\) and outside \(K\) to prove the minimizer on \(K\) is global:}{1.0in}
\end{skeletonbox}

\newpage
\fullproofpage{Homework 12, Problem 5}{Use this page for the setup and compact restriction.}

\newpage
\fullproofpage{Homework 12, Problem 5}{Use this page to prove the minimum is global.}

\newpage
\begin{solutionbox}{Student-Derived Solution (No official solution available yet)}
First, \(f\) is continuous on \(\R^2\). The map \((x,y)\mapsto x^2+y^2\) is continuous, the map \((x,y)\mapsto x-y\) is continuous, and \(\sin\) is continuous by the allowed calculus facts. Hence \((x,y)\mapsto \sin(x-y)\) is continuous, and so \(f\) is continuous.

Since \(\sin(x-y)\ge -1\), we have
\[
  f(x,y)=x^2+y^2+\sin(x-y)\ge x^2+y^2-1.
\]
In particular, if \(x^2+y^2\ge1\), then \(f(x,y)\ge0\).

Now consider
\[
  p=\left(-\frac12,\frac12\right).
\]
Then
\[
  f(p)
  =
  \frac14+\frac14+\sin(-1)
  =
  \frac12-\sin(1)
  <0.
\]
Also \(p\) lies in the closed unit disk
\[
  K=\{(x,y)\in\R^2:x^2+y^2\le1\}.
\]
The set \(K\) is closed and bounded in \(\R^2\), hence compact. By the extreme value theorem, since \(f\) is continuous, there exists \(q\in K\) such that
\[
  f(q)\le f(z)
  \qquad\text{for all }z\in K.
\]
We claim \(q\) is a global minimizer on \(\R^2\). If \(z\in K\), this is true by the choice of \(q\). If \(z\notin K\), then \(z=(x,y)\) satisfies \(x^2+y^2>1\), so \(f(z)\ge0\). But
\[
  f(q)\le f(p)<0\le f(z).
\]
Thus \(f(q)\le f(z)\) for every \(z\in\R^2\). Therefore \(f\) has a global minimum.
\end{solutionbox}

\newpage
\reflectionpage{Homework 12, Problem 5}

% ---------------------------------------------------------------------------
% HW12 Problem 6
% ---------------------------------------------------------------------------
\newpage
\subsection{Problem 6}

\begin{problemstatementbox}{Homework 12, Problem 6}
Let \(X\) and \(Y\) be metric spaces and \(f:X\to Y\) a function. Prove that if \(X\) is compact and \(f\) is continuous and bijective, then the inverse function \(f^{-1}:Y\to X\) is also continuous. (Hint: It suffices to show that the preimage of any closed subset of \(X\) under \(f^{-1}\) is closed in \(Y\).)
\end{problemstatementbox}

\newpage
\firstdiagnosis{Homework 12, Problem 6}

\newpage
\recipecard{Homework 12, Problem 6}

\newpage
\begin{selfexplanationbox}{Homework 12, Problem 6}
\workquestion{1. Which closed-set criterion for continuity is suggested by the hint?}{1.0in}
\workquestion{2. If \(C\subseteq X\) is closed, why is \(C\) compact?}{0.95in}
\workquestion{3. Why is \(f(C)\) compact in \(Y\)?}{0.95in}
\workquestion{4. Why is \(f(C)\) closed in \(Y\)?}{0.95in}
\workquestion{5. Why is \((f^{-1})^{-1}(C)=f(C)\) when \(f\) is bijective?}{1.05in}
\end{selfexplanationbox}

\newpage
\begin{skeletonbox}{Homework 12, Problem 6}
\textbf{Use the closed-set criterion for \(f^{-1}\).}
\workquestion{Let \(C\subseteq X\) be closed. Since \(X\) is compact, conclude \(C\) is:}{0.85in}
\workquestion{Since \(f\) is continuous, conclude \(f(C)\) is:}{0.85in}
\workquestion{Since \(Y\) is a metric space, compact subsets of \(Y\) are:}{0.85in}
\workquestion{Use bijectivity to compute \((f^{-1})^{-1}(C)=\)}{0.75in}
\workquestion{Conclude the preimage under \(f^{-1}\) of every closed set \(C\subseteq X\) is closed in \(Y\):}{0.95in}
\workquestion{Apply the closed-set criterion to finish:}{0.75in}
\end{skeletonbox}

\newpage
\fullproofpage{Homework 12, Problem 6}{Use this page for the complete proof.}

\newpage
\begin{solutionbox}{Student-Derived Solution (No official solution available yet)}
We prove that \(f^{-1}:Y\to X\) is continuous using the closed-set criterion.

Let \(C\subseteq X\) be closed. Since \(X\) is compact and \(C\) is closed in \(X\), \(C\) is compact. Since \(f:X\to Y\) is continuous, the image of a compact set under \(f\) is compact. Therefore \(f(C)\) is compact in \(Y\).

Because \(Y\) is a metric space, compact subsets of \(Y\) are closed. Hence \(f(C)\) is closed in \(Y\).

Since \(f\) is bijective,
\[
  (f^{-1})^{-1}(C)=f(C).
\]
Thus the preimage under \(f^{-1}\) of every closed subset \(C\subseteq X\) is closed in \(Y\). By the closed-set criterion for continuity, \(f^{-1}\) is continuous.
\end{solutionbox}

\newpage
\reflectionpage{Homework 12, Problem 6}

\end{document}
